\section{Limites projetivos de pré-esquemas}
\label{section:IV.8}

\subsection{Introdução}
\label{subsection:IV.8.1}

\oldpage[IV-3]{5}

\begin{env}[8.1.1]
\label{IV.8.1.1}
Neste parágrafo, estudaremos sistematicamente a seguinte situação.
Seja $\mathrm{I}$ um conjunto pré-ordenado filtrante crescente, $(\mathrm{A}_\lambda,\varphi_{\mu\lambda})$ um sistema indutivo de anéis cujo conjunto de índices é $\mathrm{I}$, e seja $\mathrm{A}=\varinjlim\mathrm{A}_\lambda$ seu limite indutivo.
Para cada $\alpha\in\mathrm{I}$ e cada $\mathrm{A}_\alpha$-pré-esquema $\mathrm{X}_\alpha$, consideremos os $\mathrm{A}_\lambda$-pré-esquemas $\mathrm{X}_\lambda=\mathrm{X}_\alpha\otimes_{\mathrm{A}_\alpha}\mathrm{A}_\lambda$ para $\lambda\geqslant\alpha$, e o $\mathrm{A}$-pré-esquema $\mathrm{X}=\mathrm{X}_\alpha\otimes_{\mathrm{A}_\alpha}\mathrm{A}$; é claro que os pré-esquemas $\mathrm{X}_\lambda$ (para $\lambda\geqslant\alpha$) formam um \emph{sistema projetivo}, e veremos em \sref{IV.8.2.5} que $\mathrm{X}$ é um \emph{limite projetivo} desse sistema na categoria dos pré-esquemas.
Propomo-nos encontrar condições sobre $\mathrm{X}_\alpha$ ou sobre os $\mathrm{A}_\lambda$ que permitam demonstrar propriedades do seguinte tipo:
para que $\mathrm{X}$ tenha uma propriedade $\mathbf{P}$ (por exemplo, ser próprio sobre $\mathrm{S}=\operatorname{Spec}(\mathrm{A})$, ou irredutível, ou conexo etc.), é necessário e suficiente que exista um índice $\mu\geqslant\alpha$ tal que, para todo $\lambda\geqslant\mu$, $\mathrm{X}_\lambda$ tenha (em relação a $\mathrm{S}_\lambda=\operatorname{Spec}(\mathrm{A}_\lambda)$, conforme o caso) a mesma propriedade $\mathbf{P}$.
Obteremos enunciados análogos para propriedades de $\sh{O}_\mathrm{X}$-Módulos, dos $\mathrm{A}$-morfismos de $\mathrm{A}$-pré-esquemas etc.
Mostraremos também em \sref{IV.8.9.1} que o dado de um $\mathrm{A}$-pré-esquema de apresentação finita \sref{IV.1.6.1} é essencialmente equivalente ao dado de um $\mathrm{A}_\lambda$-pré-esquema de apresentação finita $\mathrm{X}_\lambda$ para algum $\lambda$ suficientemente grande, sendo então $\mathrm{X}$ \emph{isomorfo a} $\mathrm{X}_\lambda\otimes_{\mathrm{A}_\lambda}\mathrm{A}$.
Existem também enunciados análogos para $\sh{O}_\mathrm{X}$-Módulos de apresentação finita, seus homomorfismos, os $\mathrm{A}$-morfismos de $\mathrm{A}$-pré-esquemas de apresentação finita etc.
\end{env}

\begin{env}[8.1.2]
\label{IV.8.1.2}
A utilidade desses resultados aparecerá, por exemplo, nas questões seguintes:
\begin{enumerate}[label=\emph{\alph*)}]
  \item
  Seja $\mathrm{Y}$ um pré-esquema, $y$ um ponto de $\mathrm{Y}$, e $(\mathrm{U}_\lambda)$ o sistema projetivo filtrante decrescente das vizinhanças abertas afins de $y$ em $\mathrm{Y}$; se $\mathrm{A}_\lambda$ é o anel de $\mathrm{U}_\lambda$, os $\mathrm{A}_\lambda$ formam um sistema indutivo filtrante crescente, cujo limite indutivo $\mathrm{A}$ é o anel local $\sh{O}_y$; além disso, se denotarmos por $\mathfrak{p}_\lambda$ o ideal primo $\mathfrak{j}_y$ no anel $\mathrm{A}_\lambda$, o sistema indutivo $(\mathrm{A}_\lambda)$ é cofinal em todo sistema indutivo $(\mathrm{A}_\alpha)_f$, onde $f$ percorre $\mathrm{A}_\alpha-\mathfrak{p}_\alpha$
\oldpage[IV-3]{6}
(para um $\alpha$ fixo), pois os $\mathrm{D}(f)$ formam um sistema fundamental de vizinhanças de $y$ em $\mathrm{U}_\alpha$ e, portanto, em $\mathrm{Y}$.
Por conseguinte, os resultados deste parágrafo implicam que a geometria algébrica dos $\sh{O}_y$-pré-esquemas de apresentação finita (e a teoria dos Módulos de apresentação finita sobre esses pré-esquemas) é essencialmente equivalente à geometria algébrica dos pré-esquemas de apresentação finita em vizinhanças abertas «suficientemente pequenas» de $y$.
Assim, o enunciado \sref{IV.8.10.5},~(xii) implica que, se um morfismo de apresentação finita $f:\mathrm{X}\to\mathrm{Y}$ é tal que $\mathrm{X}\times_\mathrm{Y}\operatorname{Spec}(\sh{O}_y)$ é próprio sobre $\operatorname{Spec}(\sh{O}_y)$, é necessário e suficiente que exista uma vizinhança aberta $\mathrm{U}$ de $y$ em $\mathrm{Y}$ tal que $f^{-1}(\mathrm{U})$ seja próprio sobre $\mathrm{U}$.

Um caso particular importante, clássico em certo sentido, é aquele em que $\mathrm{Y}$ é integral e $y$ é seu ponto genérico, de modo que $\sh{O}_y$ nada mais é do que \emph{o corpo} $\mathrm{R}(\mathrm{Y})=\mathrm{K}$ \emph{de funções racionais sobre} $\mathrm{Y}$.
Os resultados deste parágrafo equivalem então a interpretar a geometria algébrica sobre $\mathrm{K}$ em termos de geometria algébrica sobre subconjuntos abertos não vazios «suficientemente pequenos» de $\mathrm{Y}$, isto é, intuitivamente, em termos de «famílias» de objetos geométricos indexadas pelos pontos de um tal aberto.
Além disso, esse ponto de vista vem sendo habitualmente utilizado há muito tempo, não só na geometria algébrica sobre corpos algebricamente fechados, mas também no estudo aritmético de variedades definidas sobre um \emph{corpo de números} $\mathrm{K}$ (uma extensão finita de $\mathbf{Q}$), considerando este último como o corpo de frações de seu anel de inteiros $\mathrm{A}$ (a «teoria da \emph{redução módulo} $\mathfrak{p}$» ($\mathfrak{p}$ um ideal primo de $\mathrm{A}$); cf.\ \sref[I]{I.3.7}).
Os resultados das seções \textsection\textsection\,8 e~9 fornecem, assim, entre outras coisas, alguns fundamentos da linguagem da «redução módulo $\mathfrak{p}$» em aritmética.

Observemos que, no exemplo aqui considerado, os morfismos $\mathrm{S}_\mu\to\mathrm{S}_\lambda$ (para $\lambda\leqslant\mu$) são as \emph{imersões abertas} canônicas $\mathrm{U}_\mu\hookrightarrow\mathrm{U}_\lambda$ e, \emph{a fortiori}, são morfismos \emph{planos} (embora em geral não sejam fielmente planos), o que explica o interesse dos enunciados que recorrem a tal restrição.

  \item
  Suponhamos que os $\mathrm{A}_\lambda$ sejam \emph{corpos}, de modo que $\mathrm{A}=\varinjlim\mathrm{A}_\lambda$ também é um corpo.
Esse caso ocorre geralmente quando se parte de dados geométricos sobre um corpo arbitrário $\mathrm{K}$, considerado como extensão de um corpo $k$ (por exemplo, o subcorpo primo de $\mathrm{K}$).
É então útil considerar $\mathrm{K}$ como limite indutivo de suas subextensões que são \emph{de tipo finito sobre $k$}, o que permite reduzir muitas questões ao caso em que $\mathrm{K}$ é uma extensão de tipo finito de $k$.
Utilizando também o método esboçado em \emph{a)}, pode-se reduzir, em geral, ao caso de um anel de base $\mathrm{A}$ que seja uma \emph{álgebra íntegra de tipo finito sobre $k$}.

Observemos que, neste exemplo, os morfismos $\mathrm{S}_\mu\to\mathrm{S}_\lambda$ são \emph{fielmente planos}.

  \item
  Suponhamos que nos interessem propriedades, locais sobre $\mathrm{Y}$, de pré-esquemas de apresentação finita sobre um pré-esquema arbitrário $\mathrm{Y}$, que doravante podemos supor ser um esquema afim de anel $\mathrm{A}$.
É então útil considerar $\mathrm{A}$ como limite indutivo de seus subanéis que são \emph{$\mathbf{Z}$-álgebras de tipo finito}, o que permite reduzir muitas questões ao caso em que $\mathrm{Y}$ é o espectro de uma dessas álgebras.
Essa é a explicação do «ponto de vista kroneckeriano», segundo o qual a geometria algébrica se reduz à
\oldpage[IV-3]{7}
geometria algébrica dos pré-esquemas de tipo finito sobre $\mathbf{Z}$ (que às vezes se qualifica de «geometria algébrica absoluta»).
Este exemplo mostra, em particular, que, na maioria das questões «relativas» a um pré-esquema de base $\mathrm{Y}$, pode-se reduzir ao caso em que $\mathrm{Y}$ é \emph{noetheriano}.

Observemos que, neste exemplo, diferentemente dos anteriores, os morfismos $\mathrm{S}_\mu\to\mathrm{S}_\lambda$ em geral não têm nenhuma propriedade particular de regularidade.
\end{enumerate}

No que segue, ao aplicarmos os resultados subsequentes a qualquer uma das três situações particulares que acabamos de descrever, não entraremos em detalhes sobre o procedimento que permite realizar essas aplicações e nos limitaremos a remeter ao que precede.
\end{env}

\begin{env}[8.1.3]
\label{IV.8.1.3}
No exemplo \emph{a)} de \sref{IV.8.1.2}, vimos que, se $\mathrm{Y}$ é um pré-esquema integral com ponto genérico $y$, $f:\mathrm{X}\to\mathrm{Y}$ é um morfismo de apresentação finita, e a fibra genérica $f^{-1}(y)$ é própria sobre $\boldsymbol{k}(y)$, então existe uma vizinhança aberta $\mathrm{U}$ de $y$ em $\mathrm{Y}$ tal que $f^{-1}(\mathrm{U})$ é próprio sobre $\mathrm{U}$; deduz-se \emph{a fortiori} que, para todo $s\in\mathrm{U}$, $f^{-1}(s)$ é próprio sobre $\boldsymbol{k}(s)$.
É necessária, porém, uma recíproca que assegure que, se $f^{-1}(s)$ é próprio sobre $\boldsymbol{k}(s)$ para «suficientemente muitos» pontos $s\in\mathrm{U}$, então $f^{-1}(y)$ é próprio sobre $\boldsymbol{k}(y)$.
Por exemplo, suponhamos que $\mathrm{X}$ e $\mathrm{Y}$ sejam pré-esquemas algébricos sobre um corpo algebricamente fechado $k$ (para fixar as ideias, pode-se tomar como $k$ o corpo $\mathbf{C}$ dos números complexos); às vezes é necessário saber que, se, para todo $s\in\mathrm{Y}$, \emph{racional sobre $k$}, a fibra $f^{-1}(s)$ é própria sobre $\boldsymbol{k}(s)$, então a fibra $f^{-1}(y)$ é própria sobre $\boldsymbol{k}(y)$ e, por conseguinte, $f^{-1}(\mathrm{U})$ é próprio sobre $\mathrm{U}$ para uma vizinhança $\mathrm{U}$ de $y$\footnote{Observemos que tal enunciado é, em última análise, de natureza puramente geométrica, pois recorre apenas a pontos racionais sobre $k$, e não a pontos genéricos; por exemplo, quando $k=\mathbf{C}$, esse enunciado tem um significado topológico evidente para o analista, interpretando «próprio» no sentido topológico do termo para os espaços analíticos subjacentes formados pelos pontos de $\mathrm{X}$ e $\mathrm{Y}$ racionais sobre $\mathbf{C}$.}.
Ora, esse enunciado decorre facilmente do seguinte: o conjunto $\mathrm{E}$ dos pontos $s\in\mathrm{Y}$ tais que $f^{-1}(s)$ é próprio sobre $\boldsymbol{k}(s)$ é \emph{construtível} (e, portanto, coincide com $\mathrm{Y}$ se contém os pontos fechados de $\mathrm{Y}$, pelo teorema dos zeros de Hilbert \sref[I]{I.10.4.8}); isto equivale mais uma vez a dizer que, se $f^{-1}(y)$ \emph{não é próprio} sobre $\boldsymbol{k}(y)$, existe uma vizinhança aberta $\mathrm{U}$ de $y$ tal que $f^{-1}(s)$ não é próprio sobre $\boldsymbol{k}(s)$ para todo $s\in\mathrm{U}$ (cf.\ \sref{IV.9.6.1},~(iv)).
Este exemplo ilustra o interesse em desenvolver sistematicamente \emph{critérios de construtibilidade}: é o que será feito na \textsection\,9.
\end{env}

\subsection{Limites projetivos de pré-esquemas}
\label{subsection:IV.8.2}

\begin{env}[8.2.1]
\label{IV.8.2.1}
Seja $\mathrm{S}_0$ um espaço anelado, $\mathrm{L}$ um conjunto pré-ordenado filtrante crescente, $(\sh{A}_\lambda,\varphi_{\mu\lambda})$ um sistema indutivo de $\sh{O}_{\mathrm{S}_0}$-Álgebras (não necessariamente comutativas), cujo conjunto de índices é $\mathrm{L}$.
Sabemos que, considerado como sistema indutivo de $\sh{O}_{\mathrm{S}_0}$-Módulos, $(\sh{A}_\lambda,\varphi_{\mu\lambda})$ admite um limite indutivo $\sh{A}$; denotemos por $\varphi_\lambda:\sh{A}_\lambda\to\sh{A}$ o homomorfismo canônico de $\sh{O}_{\mathrm{S}_0}$-Módulos.
Seja $m_\lambda:\sh{A}_\lambda\otimes\sh{A}_\lambda\to\sh{A}_\lambda$ o homomorfismo de $\sh{O}_{\mathrm{S}_0}$-Módulos que define a multiplicação na $\sh{O}_{\mathrm{S}_0}$-Álgebra $\sh{A}_\lambda$; a hipótese sobre os $\varphi_{\mu\lambda}$ implica que os $m_\lambda$ formam um sistema indutivo de homomorfismos,
\oldpage[IV-3]{8}
e, como o funtor $\varinjlim$ comuta com os produtos tensoriais, $m=\varinjlim m_\lambda$ é um homomorfismo $\sh{A}\otimes\sh{A}\to\sh{A}$ de $\sh{O}_{\mathrm{S}_0}$-Módulos; passando ao limite nos diagramas comutativos que exprimem a associatividade de $m_\lambda$ e a existência da seção unidade em $\sh{A}_\lambda$, vê-se que $m$ dota $\sh{A}$ de uma estrutura de $\sh{O}_{\mathrm{S}_0}$-Álgebra e que $\varphi_\lambda$ é um homomorfismo de $\sh{O}_{\mathrm{S}_0}$-Álgebras para todo $\lambda\in\mathrm{L}$.
Além disso, $\sh{A}$ é o limite indutivo do sistema $(\sh{A}_\lambda,\varphi_{\mu\lambda})$ \emph{na categoria das $\sh{O}_{\mathrm{S}_0}$-Álgebras}; em outras palavras, para toda $\sh{O}_{\mathrm{S}_0}$-Álgebra $\sh{B}$, a aplicação canônica
\begin{equation}
  \operatorname{Hom}_{\sh{O}_{\mathrm{S}_0}\text{-}\mathrm{Alg.}}(\sh{A},\sh{B})\to\varprojlim\operatorname{Hom}_{\sh{O}_{\mathrm{S}_0}\text{-}\mathrm{Alg.}}(\sh{A}_\lambda,\sh{B})
  \tag{8.2.1.1}
\label{IV.8.2.1.1}
\end{equation}
que associa a cada homomorfismo $f:\sh{A}\to\sh{B}$ de $\sh{O}_{\mathrm{S}_0}$-Álgebras a família $(f\circ\varphi_\lambda)$ é \emph{bijetiva}.
Com efeito, já sabemos que ela é injetiva e identifica $\operatorname{Hom}_{\sh{O}_{\mathrm{S}_0}\text{-}\mathrm{Alg.}}(\sh{A},\sh{B})$ com um subconjunto de $\varprojlim\operatorname{Hom}_{\sh{O}_{\mathrm{S}_0}\text{-}\mathrm{Mod.}}(\sh{A}_\lambda,\sh{B})$; tudo se reduz a verificar que, se $(f_\lambda)$ é um sistema indutivo de homomorfismos de $\sh{O}_{\mathrm{S}_0}$-Álgebras, $f_\lambda:\sh{A}_\lambda\to\sh{B}$, seu limite indutivo $f:\sh{A}\to\sh{B}$, que por definição é um homomorfismo de $\sh{O}_{\mathrm{S}_0}$-\emph{Módulos}, é também um homomorfismo de $\sh{O}_{\mathrm{S}_0}$-\emph{Álgebras}; mas isso decorre da passagem ao limite indutivo no diagrama comutativo de $\sh{O}_{\mathrm{S}_0}$-Módulos
\[
\xymatrix{
  \sh{A}_\lambda\otimes\sh{A}_\lambda \ar[r]^{f_\lambda\otimes f_\lambda} \ar[d]_{m_\lambda} & \sh{B}\otimes\sh{B} \ar[d] \\
  \sh{A}_\lambda \ar[r]_{f_\lambda} & \sh{B}
}
\]
e do fato de que o funtor $\varinjlim$ comuta com os produtos tensoriais.
Observemos, por fim, que, se os $\sh{A}_\lambda$ são $\sh{O}_{\mathrm{S}_0}$-Álgebras comutativas, o mesmo vale para $\sh{A}$.
\end{env}

\begin{env}[8.2.2]
\label{IV.8.2.2}
Suponhamos agora que $\mathrm{S}_0$ seja um \emph{pré-esquema} e que as $\sh{A}_\lambda$ sejam $\sh{O}_{\mathrm{S}_0}$-Álgebras (comutativas) \emph{quase-coerentes}; sabemos então que $\sh{A}=\varinjlim\sh{A}_\lambda$ é uma $\sh{O}_{\mathrm{S}_0}$-Álgebra quase-coerente \sref[I]{I.4.1.1}.
Denotemos por $\mathrm{S}_\lambda$ (resp. $\mathrm{S}$) o espectro da $\sh{O}_{\mathrm{S}_0}$-Álgebra $\sh{A}_\lambda$ (resp. $\sh{A}$) \sref[II]{II.1.3.1}, e sejam $u_{\lambda\mu}:\mathrm{S}_\mu\to\mathrm{S}_\lambda$ (para $\lambda\leqslant\mu$) e $u_\lambda:\mathrm{S}\to\mathrm{S}_\lambda$ os $\mathrm{S}_0$-morfismos correspondentes aos homomorfismos $\varphi_{\mu\lambda}$ e $\varphi_\lambda$, respectivamente \sref[II]{II.1.2.7}; é claro que $(\mathrm{S}_\lambda,u_{\lambda\mu})$ é um \emph{sistema projetivo} na categoria dos $\mathrm{S}_0$-pré-esquemas.
Observemos que os $u_{\lambda\mu}$ e $u_\lambda$ são morfismos \emph{afins} \sref[II]{II.1.6.2}, portanto \emph{quase-compactos} e \emph{separados}.
\end{env}

\begin{proposition}[8.2.3]
\label{IV.8.2.3}
Com as notações de \sref{IV.8.2.2}, os morfismos $u_\lambda:\mathrm{S}\to\mathrm{S}_\lambda$ fazem de $\mathrm{S}$ um limite projetivo do sistema projetivo $(\mathrm{S}_\lambda,u_{\lambda\mu})$ na categoria dos pré-esquemas.
Além disso, se $h:\mathrm{S}_0\to\mathrm{T}$ é um morfismo que faz de todo $\mathrm{S}_0$-pré-esquema um $\mathrm{T}$-pré-esquema, $\mathrm{S}$ é também um limite projetivo do sistema $(\mathrm{S}_\lambda,u_{\lambda\mu})$ na categoria dos $\mathrm{T}$-pré-esquemas.
\end{proposition}

\begin{proof}
Demonstremos primeiro a segunda afirmação do enunciado no caso $\mathrm{T}=\mathrm{S}_0$.
\oldpage[IV-3]{9}
Tudo se reduz a demonstrar que, se $\mathrm{X}$ é um $\mathrm{S}_0$-pré-esquema qualquer, a aplicação canônica
\begin{equation}
  \operatorname{Hom}_{\mathrm{S}_0}(\mathrm{X},\mathrm{S})\to\varprojlim\operatorname{Hom}_{\mathrm{S}_0}(\mathrm{X},\mathrm{S}_\lambda)
  \tag{8.2.3.1}
\label{IV.8.2.3.1}
\end{equation}
que associa a cada $\mathrm{S}_0$-morfismo $v:\mathrm{X}\to\mathrm{S}$ a família $(u_\lambda\circ v)$ é \emph{bijetiva}.
Ora, se $g:\mathrm{X}\to\mathrm{S}_0$ é o morfismo estrutural e pusermos $\sh{B}=g_*(\sh{O}_\mathrm{X})$, que é uma $\sh{O}_{\mathrm{S}_0}$-Álgebra, a aplicação \sref{IV.8.2.3.1} identifica-se canonicamente com \sref{IV.8.2.1.1} \sref[II]{II.1.2.7}, e a conclusão decorre do que vimos em \sref{IV.8.2.1}.

As demais afirmações de \sref{IV.8.2.3} são consequências do seguinte lema geral:
\end{proof}

\begin{lemma}[8.2.4]
\label{IV.8.2.4}
Seja $\mathbf{C}$ uma categoria, $\mathrm{T}$ um objeto de $\mathbf{C}$ e $\mathbf{C}_\mathrm{T}$ a subcategoria dos objetos de $\mathbf{C}$ sobre $\mathrm{T}$.
Seja $(\mathrm{S}_\lambda,u_{\lambda\mu})$ um sistema projetivo em $\mathbf{C}_\mathrm{T}$; então todo limite projetivo desse sistema em $\mathbf{C}_\mathrm{T}$ é também um limite projetivo em $\mathbf{C}$, e reciprocamente.
\end{lemma}

\begin{proof}
Seja $f_\lambda:\mathrm{S}_\lambda\to\mathrm{T}$ o morfismo estrutural.
Suponhamos que $\mathrm{S}$ seja um limite projetivo de $(\mathrm{S}_\lambda,u_{\lambda\mu})$ em $\mathbf{C}$ e denotemos por $u_\lambda:\mathrm{S}\to\mathrm{S}_\lambda$ os morfismos canônicos correspondentes.
Consideremos um sistema projetivo de $\mathrm{T}$-morfismos (em $\mathbf{C}$) $w_\lambda:\mathrm{Y}\to\mathrm{S}_\lambda$, onde $\mathrm{Y}\in\mathbf{C}_\mathrm{T}$.
Por hipótese, existe um único morfismo (em $\mathbf{C}$) $w:\mathrm{Y}\to\mathrm{S}$ tal que $w_\lambda=u_\lambda\circ w$ para todo $\lambda$.
A hipótese de que os $u_{\lambda\mu}$ são $\mathrm{T}$-morfismos implica que os morfismos $f_\lambda\circ u_\lambda:\mathrm{S}\to\mathrm{T}$ são todos iguais, e esse composto $f$ dota $\mathrm{S}$ de uma estrutura de $\mathrm{T}$-objeto.
Se $g:\mathrm{Y}\to\mathrm{T}$ é o morfismo estrutural de $\mathrm{Y}$, temos $f\circ w=f_\lambda\circ u_\lambda\circ w=f_\lambda\circ w_\lambda=g$ para todo $\lambda$, o que demonstra que $w$ é um $\mathrm{T}$-morfismo.
Reciprocamente, suponhamos (com as mesmas notações) que $\mathrm{S}$ seja um limite projetivo de $(\mathrm{S}_\lambda,u_{\lambda\mu})$ em $\mathbf{C}_\mathrm{T}$ e consideremos agora um sistema projetivo de morfismos (em $\mathbf{C}$) $w_\lambda:\mathrm{Y}\to\mathrm{S}_\lambda$.
Os compostos $f_\lambda\circ w_\lambda:\mathrm{Y}\to\mathrm{T}$ são então todos iguais: com efeito, para dois índices quaisquer $\lambda$, $\mu$, existe um índice $\nu$ tal que $\lambda\leqslant\nu$ e $\mu\leqslant\nu$, de onde $f_\nu=f_\lambda\circ u_{\lambda\nu}=f_\mu\circ u_{\mu\nu}$; como $w_\lambda=u_{\lambda\nu}\circ w_\nu$ e $w_\mu=u_{\mu\nu}\circ w_\nu$, temos $f_\lambda\circ w_\lambda=f_\lambda\circ u_{\lambda\nu}\circ w_\nu=f_\nu\circ w_\nu$ e $f_\mu\circ w_\mu=f_\mu\circ u_{\mu\nu}\circ w_\nu=f_\nu\circ w_\nu$.
Se $g:\mathrm{Y}\to\mathrm{T}$ é o único morfismo assim definido, $g$ dota $\mathrm{Y}$ de uma estrutura de $\mathrm{T}$-objeto, e os $w_\lambda$ são $\mathrm{T}$-morfismos; portanto, têm um limite projetivo $w:\mathrm{Y}\to\mathrm{S}$ que é um $\mathrm{T}$-morfismo e, \emph{a fortiori}, um morfismo de $\mathbf{C}$; além disso, a primeira parte do argumento mostra que qualquer limite projetivo $w'$ (em $\mathbf{C}$) do sistema projetivo $(w_\lambda)$ é também um $\mathrm{T}$-morfismo e, por conseguinte, igual a $w$, o que conclui a demonstração do lema.
\end{proof}

\begin{proposition}[8.2.5]
\label{IV.8.2.5}
Nas condições de \sref{IV.8.2.2}, seja $\alpha$ um elemento de $\mathrm{L}$ e seja $\mathrm{X}_\alpha$ um $\mathrm{S}_\alpha$-pré-esquema.
Para todo $\lambda\geqslant\alpha$, ponhamos $\mathrm{X}_\lambda=\mathrm{X}_\alpha\times_{\mathrm{S}_\alpha}\mathrm{S}_\lambda$ e, para $\alpha\leqslant\lambda\leqslant\mu$, ponhamos $v_{\lambda\mu}=1_{\mathrm{X}_\alpha}\times u_{\lambda\mu}$, de modo que $(\mathrm{X}_\lambda,v_{\lambda\mu})$ seja um sistema projetivo de $\mathrm{X}_\alpha$-pré-esquemas, cujo conjunto de índices é formado pelos $\lambda\geqslant\alpha$ de $\mathrm{L}$.
Ponhamos também $\mathrm{X}=\mathrm{X}_\alpha\times_{\mathrm{S}_\alpha}\mathrm{S}$ e $v_\lambda=1_{\mathrm{X}_\alpha}\times u_\lambda$.
Então, os $\mathrm{X}_\alpha$-morfismos $v_\lambda:\mathrm{X}\to\mathrm{X}_\lambda$ fazem de $\mathrm{X}$ um limite projetivo do sistema projetivo $(\mathrm{X}_\lambda,v_{\lambda\mu})$ na categoria dos $\mathrm{X}_\alpha$-pré-esquemas ou na categoria de todos os pré-esquemas.
\end{proposition}

\begin{proof}
Isto também decorrerá do seguinte lema geral:
\end{proof}

\begin{lemma}[8.2.6]
\label{IV.8.2.6}
Seja $\mathbf{C}$ uma categoria na qual existem produtos fibrados, $q:\mathrm{T}'\to\mathrm{T}$ um morfismo de $\mathbf{C}$, e $\mathbf{C}_\mathrm{T}$ (resp. $\mathbf{C}_{\mathrm{T}'}$) a categoria dos objetos de $\mathbf{C}$ sobre
\oldpage[IV-3]{10}
$\mathrm{T}$ (resp. $\mathrm{T}'$).
Seja $(\mathrm{S}_\lambda,u_{\lambda\mu})$ um sistema projetivo (não necessariamente filtrante) em $\mathbf{C}_\mathrm{T}$ e ponhamos $\mathrm{S}'_\lambda=\mathrm{S}_\lambda\times_\mathrm{T}\mathrm{T}'$, $u'_{\lambda\mu}=u_{\lambda\mu}\times 1_{\mathrm{T}'}$, de modo que $(\mathrm{S}'_\lambda,u'_{\lambda\mu})$ seja um sistema projetivo em $\mathbf{C}_{\mathrm{T}'}$.
Então, se $(\mathrm{S},u_\lambda)$ é um limite projetivo de $(\mathrm{S}_\lambda,u_{\lambda\mu})$ em $\mathbf{C}_\mathrm{T}$, $(\mathrm{S}\times_\mathrm{T}\mathrm{T}',u_\lambda\times 1_{\mathrm{T}'})$ é um limite projetivo de $(\mathrm{S}'_\lambda,u'_{\lambda\mu})$ em $\mathbf{C}_{\mathrm{T}'}$.
\end{lemma}

\begin{proof}
Por hipótese, para todo $\lambda$, tem-se um diagrama comutativo
\[
\xymatrix{
  \mathrm{S}' \ar[r]^{u'_\lambda} \ar[d]_p & \mathrm{S}'_\lambda \ar[r]^{h'_\lambda} \ar[d]_{p_\lambda} & \mathrm{T}' \ar[d]^q \\
  \mathrm{S} \ar[r]_{u_\lambda} & \mathrm{S}_\lambda \ar[r]_{h_\lambda} & \mathrm{T}
}
\]
onde pusemos $\mathrm{S}'=\mathrm{S}\times_\mathrm{T}\mathrm{T}'$, $u'_\lambda=u_\lambda\times 1_{\mathrm{T}'}$, $h'_\lambda=h_\lambda\times 1_{\mathrm{T}'}$.
Seja $\mathrm{Y}$ um $\mathrm{T}'$-objeto, $g':\mathrm{Y}\to\mathrm{T}'$, e consideremos um sistema projetivo de $\mathrm{T}'$-morfismos $w'_\lambda:\mathrm{Y}\to\mathrm{S}'_\lambda$.
Então, $\mathrm{Y}$ é um $\mathrm{T}$-objeto para o morfismo $g=q\circ g'$, e os $w_\lambda=p_\lambda\circ w'_\lambda$ são $\mathrm{T}$-morfismos, pois $h_\lambda\circ w_\lambda=h_\lambda\circ p_\lambda\circ w'_\lambda=q\circ h'_\lambda\circ w'_\lambda=q\circ g'$ por hipótese.
Além disso, verifica-se imediatamente que $(w_\lambda)$ é um sistema projetivo.
Por hipótese, existe então um único $\mathrm{T}$-morfismo $w:\mathrm{Y}\to\mathrm{S}$ tal que $u_\lambda\circ w=w_\lambda$ para todo $\lambda$.
Pela definição do produto fibrado, existe um único $\mathrm{T}'$-morfismo $w':\mathrm{Y}\to\mathrm{S}'$ tal que $p\circ w'=w$.
Então, $u_\lambda\circ p\circ w'=u_\lambda\circ w=w_\lambda=p_\lambda\circ w'_\lambda$, o que também se pode escrever $p_\lambda\circ u'_\lambda\circ w'=p_\lambda\circ w'_\lambda$; por outro lado, como $w'$ é um $\mathrm{T}'$-morfismo, $h'_\lambda\circ u'_\lambda\circ w'=g'=h'_\lambda\circ w'_\lambda$. A definição de $\mathrm{S}'_\lambda$ como o produto fibrado $\mathrm{S}_\lambda\times_\mathrm{T}\mathrm{T}'$ dá, portanto, $u'_\lambda\circ w'=w'_\lambda$, e é imediato que $w'$ é o único $\mathrm{T}'$-morfismo que satisfaz essas relações, donde resulta o lema.
\end{proof}

\begin{remark}[8.2.7]
\label{IV.8.2.7}
Dado um espaço anelado $\mathrm{S}$, os limites indutivos relativos a um conjunto pré-ordenado $\mathrm{L}$ (não necessariamente filtrante) existem na categoria das $\sh{O}_\mathrm{S}$-Álgebras comutativas, pois o limite indutivo filtrante existe por \sref{IV.8.2.1} e, por outro lado, para quaisquer dois homomorfismos de $\sh{O}_\mathrm{S}$-Álgebras $\sh{A}\to\sh{B}$, $\sh{A}\to\sh{C}$, o produto tensorial $\sh{B}\otimes_{\sh{A}}\sh{C}$ é a «soma amalgamada» correspondente nessa categoria.

Quando $\mathrm{S}$ é um pré-esquema, sabemos que o produto tensorial $\sh{B}\otimes_{\sh{A}}\sh{C}$ é uma $\sh{O}_\mathrm{S}$-Álgebra quase-coerente sempre que $\sh{A}$, $\sh{B}$ e $\sh{C}$ também o sejam \sref[I]{I.1.3.13}; conclui-se daí que, na categoria das $\sh{O}_\mathrm{S}$-Álgebras quase-coerentes, existem sempre limites indutivos para qualquer conjunto pré-ordenado de índices.
Isto permite generalizar a definição de limite projetivo de pré-esquemas e as proposições \sref{IV.8.2.3} e \sref{IV.8.2.5} ao caso em que o conjunto pré-ordenado $\mathrm{L}$ não seja necessariamente filtrante.
\end{remark}

\begin{env}[8.2.8]
\label{IV.8.2.8}
Com as notações de \sref{IV.8.2.2}, ponhamos $u_{\lambda\mu}=(\psi_{\lambda\mu},\theta_{\lambda\mu})$ e $u_\lambda=(\psi_\lambda,\theta_\lambda)$; $(\mathrm{S}_\lambda,\psi_{\lambda\mu})$ é, assim, um sistema projetivo de espaços topológicos, e $(\psi_\lambda)$ é um sistema projetivo de aplicações contínuas $\mathrm{S}\to\mathrm{S}_\lambda$ entre os espaços subjacentes aos pré-esquemas $\mathrm{S}$ e $\mathrm{S}_\lambda$, respectivamente.
\end{env}

\begin{proposition}[8.2.9]
\label{IV.8.2.9}
Com as notações de \sref{IV.8.2.8}, o limite projetivo do sistema projetivo $(\psi_\lambda)$ de aplicações contínuas é um homeomorfismo do espaço subjacente de $\mathrm{S}$ sobre o limite projetivo do sistema projetivo $(\mathrm{S}_\lambda,\psi_{\lambda\mu})$ de espaços topológicos.
\end{proposition}

\begin{proof}
\oldpage[IV-3]{11}
Seja $\mathrm{T}$ o espaço limite topológico do sistema $(\mathrm{S}_\lambda,\psi_{\lambda\mu})$ e ponhamos $\psi=\varprojlim\psi_\lambda:\mathrm{S}\to\mathrm{T}$.
Podemos restringir-nos ao caso em que $\mathrm{S}_0=\mathrm{S}_\alpha$ para algum $\alpha\in\mathrm{L}$ e em que $\lambda\geqslant\alpha$.

\medskip
\noindent
(i) Mostremos primeiro que a topologia de $\mathrm{S}$ é a imagem inversa, por $\psi$, da topologia de $\mathrm{T}$; em outras palavras, se $\pi_\lambda:\mathrm{T}\to\mathrm{S}_\lambda$ é a aplicação canônica, trata-se de demonstrar que todo aberto de $\mathrm{S}$ é união de abertos da forma $\psi^{-1}(\pi_\lambda^{-1}(\mathrm{U}_\lambda))$, em que $\mathrm{U}_\lambda$ é aberto em $\mathrm{S}_\lambda$.
Ora, por definição, todo aberto de $\mathrm{S}$ é união de abertos obtidos do seguinte modo: considera-se um aberto afim $\mathrm{U}_0$ de $\mathrm{S}_0$, de anel $\mathrm{A}_0$, de modo que $\psi_\alpha^{-1}(\mathrm{U}_0)$ seja o aberto afim de $\mathrm{S}$ de anel $\mathrm{A}=\Gamma(\mathrm{U}_0,\sh{A})$; em seguida, toma-se um elemento $f\in\mathrm{A}$ e considera-se, em $\operatorname{Spec}(\mathrm{A})$, identificado com $\psi_\alpha^{-1}(\mathrm{U}_0)$, o aberto $\mathrm{D}(f)$.
São esses os abertos $\mathrm{D}(f)$ que formam uma base da topologia de $\mathrm{S}$ \sref[II]{II.1.3.1}.
Ora, se pusermos $\mathrm{A}_\lambda=\Gamma(\mathrm{U}_0,\sh{A}_\lambda)$, temos $\mathrm{A}=\varinjlim\mathrm{A}_\lambda$ \sref[I]{I.1.3.9}; existe, portanto, um índice $\lambda$ tal que $f$ é a imagem canônica de um elemento $f_\lambda\in\mathrm{A}_\lambda$; então $\mathrm{D}(f)=\psi_\lambda^{-1}(\mathrm{D}(f_\lambda))$ \sref[I]{I.1.2.2}, e, como $\psi_\lambda=\pi_\lambda\circ\psi$, fica demonstrada nossa afirmação.

\medskip
\noindent
(ii) Mostremos agora que $\psi$ é \emph{bijetiva}, o que concluirá a demonstração.
Como $\mathrm{S}$ é um espaço de Kolmogoroff, já se deduz de (i) que $\psi$ é injetiva, e resta demonstrar que $\psi$ é sobrejetiva.
Evidentemente, podemos substituir $\mathrm{T}$ por um aberto $\pi_\alpha^{-1}(\mathrm{U}_0)$, em que $\mathrm{U}_0$ é um aberto afim de $\mathrm{S}_\alpha=\mathrm{S}_0$; portanto, podemos restringir-nos ao caso em que os $\mathrm{S}_\lambda$ e $\mathrm{S}$ sejam afins, isto é, em que $\sh{A}_\lambda$ seja o feixe de álgebras associado a uma $\mathrm{A}_0$-álgebra $\mathrm{A}_\lambda$, e $\sh{A}$ o feixe de álgebras associado a $\mathrm{A}=\varinjlim\mathrm{A}_\lambda$; continuaremos denotando por $\varphi_{\mu\lambda}:\mathrm{A}_\lambda\to\mathrm{A}_\mu$ e $\varphi_\lambda:\mathrm{A}_\lambda\to\mathrm{A}$ os homomorfismos canônicos.
Por definição, um elemento de $\mathrm{T}$ é uma família $(\mathfrak{p}_\lambda)_{\lambda\in\mathrm{L}}$, em que $\mathfrak{p}_\lambda$ é um ideal primo de $\mathrm{A}_\lambda$ e em que $\mathfrak{p}_\lambda=\varphi_{\mu\lambda}^{-1}(\mathfrak{p}_\mu)$ para $\lambda\leqslant\mu$.
Sabemos (\sref{IV.5.13.3} e \sref{IV.5.13.1}) que existe um ideal primo $\mathfrak{p}$ de $\mathrm{A}$ tal que $\mathfrak{p}_\lambda=\varphi_\lambda^{-1}(\mathfrak{p})$ para todo $\lambda\in\mathrm{L}$, o que conclui a demonstração.
\end{proof}

Em particular, demonstramos assim o seguinte

\begin{corollary}[8.2.10]
\label{IV.8.2.10}
Seja $(\mathrm{A}_\lambda)_{\lambda\in\mathrm{L}}$ um sistema indutivo filtrante de anéis, seja $\mathrm{A}=\varinjlim\mathrm{A}_\lambda$, e sejam $\varphi_\lambda:\mathrm{A}_\lambda\to\mathrm{A}$ os homomorfismos canônicos.
A aplicação canônica $\mathfrak{p}\mapsto(\varphi_\lambda^{-1}(\mathfrak{p}))$ é um homeomorfismo de $\operatorname{Spec}(\mathrm{A})$ sobre o espaço topológico $\varprojlim\operatorname{Spec}(\mathrm{A}_\lambda)$.
\end{corollary}

\begin{corollary}[8.2.11]
\label{IV.8.2.11}
Com as notações de \sref{IV.8.2.8}, para todo aberto quase-compacto $\mathrm{U}$ de $\mathrm{S}$, existem um índice $\lambda$ e um aberto quase-compacto $\mathrm{U}_\lambda$ de $\mathrm{S}_\lambda$ tais que $\mathrm{U}=\psi_\lambda^{-1}(\mathrm{U}_\lambda)$.
\end{corollary}

\begin{proof}
Isto decorre do fato de que, pela definição da topologia do limite projetivo, os $\psi_\lambda^{-1}(\mathrm{U}_\lambda)$ ($\mathrm{U}_\lambda$ aberto quase-compacto de $\mathrm{S}_\lambda$) formam uma base da topologia de $\mathrm{S}$, e do fato de que o conjunto de índices $\mathrm{L}$ é filtrante.
\end{proof}

\begin{corollary}[8.2.12]
\label{IV.8.2.12}
Com as notações de \sref{IV.8.2.8}, o limite indutivo do sistema indutivo de homomorfismos $\theta_\lambda^\sharp:\psi_\lambda^*(\sh{O}_{\mathrm{S}_\lambda})\to\sh{O}_\mathrm{S}$ de feixes de anéis sobre $\mathrm{S}$ é um isomorfismo
\begin{equation}
  \varinjlim_\lambda\psi_\lambda^*(\sh{O}_{\mathrm{S}_\lambda})\xrightarrow{\sim}\sh{O}_\mathrm{S}.
  \label{IV.8.2.12.1}
\tag{8.2.12.1}
\end{equation}
\end{corollary}

\begin{proof}
Podemos supor, evidentemente, que os $\mathrm{S}_\lambda$ sejam afins; com as notações da demonstração de \sref{IV.8.2.9}, tudo se reduz a verificar que o limite indutivo do sistema indutivo de aplicações canônicas $(\mathrm{A}_\lambda)_{\mathfrak{p}_\lambda}\to\mathrm{A}_\mathfrak{p}$ é um isomorfismo, o que não é senão \sref{IV.5.13.3},~(ii).
\end{proof}

\begin{proposition}[8.2.13]
\label{IV.8.2.13}
Suponhamos que os morfismos $u_{\lambda\mu}:\mathrm{S}_\mu\to\mathrm{S}_\lambda$ sejam imersões abertas, de modo que $\mathrm{S}_\mu$ se identifica com o subesquema induzido em um aberto de $\mathrm{S}_\lambda$ para $\lambda\leqslant\mu$.
Então, para todo $\alpha\in\mathrm{L}$, o espaço subjacente do pré-esquema $\mathrm{S}$ identifica-se com o subespaço de $\mathrm{S}_\alpha$ dado pela interseção dos $\mathrm{S}_\lambda$ para $\lambda\geqslant\alpha$, e o feixe estrutural $\sh{O}_\mathrm{S}$ identifica-se com o feixe induzido \textnormal{(}$\mathrm{G}$, II, $1.5$\textnormal{)} por $\sh{O}_{\mathrm{S}_\alpha}$ sobre essa interseção; mais geralmente, para todo $\sh{O}_{\mathrm{S}_\alpha}$-Módulo $\sh{F}_\alpha$, $u_\alpha^*(\sh{F}_\alpha)$ identifica-se com o $\sh{O}_\mathrm{S}$-Módulo induzido por $\sh{F}_\alpha$ sobre $\mathrm{S}$.
\end{proposition}

\begin{proof}
\oldpage[IV-3]{12}
A primeira afirmação decorre de \sref{IV.8.2.9}, levando em conta a definição de limite projetivo de espaços topológicos; além disso, todos os $\psi_\lambda^*(\sh{O}_{\mathrm{S}_\lambda})$ são iguais ao feixe induzido por $\sh{O}_{\mathrm{S}_\alpha}$ sobre $\mathrm{S}$ por definição \sref[0]{0.1.3.7.1}, e, com as notações da demonstração de \sref{IV.8.2.9}, os homomorfismos $(\mathrm{A}_\lambda)_{\mathfrak{p}_\lambda}\to(\mathrm{A}_\mu)_{\mathfrak{p}_\mu}$ são bijetivos para um sistema $(\mathfrak{p}_\lambda)$ de ideais primos correspondente a um único ponto de $\mathrm{S}$; a afirmação relativa a $\sh{O}_\mathrm{S}$ decorre então de \sref{IV.8.2.12}.
A última afirmação decorre então da definição de $u_\alpha^*(\sh{F}_\alpha)$ \sref[0]{0.1.4.3.1}.
\end{proof}

\begin{remark}[8.2.14]
\label{IV.8.2.14}
Os resultados de \sref{IV.8.2.9} e \sref{IV.8.2.12} mostram que $\mathrm{S}$ é um limite projetivo do sistema projetivo $(\mathrm{S}_\lambda,u_{\lambda\mu})$ na categoria de \emph{todos} os espaços anelados (ou de todos os espaços localmente anelados).
Com efeito, seja $\mathrm{Y}$ um espaço anelado e consideremos um sistema projetivo de morfismos de espaços anelados $w_\lambda:\mathrm{Y}\to\mathrm{S}_\lambda$.
Escrevendo $w_\lambda=(\rho_\lambda,\omega_\lambda)$, os $\rho_\lambda$ formam um sistema projetivo de aplicações contínuas e, por \sref{IV.8.2.9}, seu limite projetivo identifica-se com uma aplicação contínua $\rho:\mathrm{Y}\to\mathrm{S}$ tal que $\rho_\lambda=\psi_\lambda\circ\rho$. Por outro lado, os homomorfismos
\[
 \omega_\lambda^\sharp:\rho_\lambda^*(\sh{O}_{\mathrm{S}_\lambda})\longrightarrow\sh{O}_\mathrm{Y}
\]
formam um sistema indutivo de homomorfismos de feixes de anéis. Como
\[
 \rho_\lambda^*(\sh{O}_{\mathrm{S}_\lambda})
 =\rho^*\bigl(\psi_\lambda^*(\sh{O}_{\mathrm{S}_\lambda})\bigr)
\]
e o funtor $\rho^*$ é exato, o limite indutivo dos $\rho_\lambda^*(\sh{O}_{\mathrm{S}_\lambda})$ é $\rho^*(\sh{O}_\mathrm{S})$ por \sref{IV.8.2.12}. Existe, portanto, um único homomorfismo
\[
 \omega^\sharp:\rho^*(\sh{O}_\mathrm{S})\longrightarrow\sh{O}_\mathrm{Y}
\]
tal que $\omega_\lambda^\sharp=\omega^\sharp\circ\rho^*(\theta_\lambda^\sharp)$ para todo $\lambda$, o que demonstra a afirmação.
\end{remark}

\subsection{Subconjuntos construtíveis de um limite projetivo de pré-esquemas}
\label{subsection:IV.8.3}

\begin{env}[8.3.1]
\label{IV.8.3.1}
Em tudo o que segue neste parágrafo, supomos satisfeitas as condições de \sref{IV.8.2.2} e conservamos as notações.
\end{env}

\begin{theorem}[8.3.2]
\label{IV.8.3.2}
Para todo $\lambda$, sejam $\mathrm{E}_\lambda$, $\mathrm{F}_\lambda$ dois subconjuntos de $\mathrm{S}_\lambda$.
Ponhamos
\begin{equation}
  \mathrm{E}=\bigcap_\lambda u_\lambda^{-1}(\mathrm{E}_\lambda),\qquad \mathrm{F}=\bigcup_\lambda u_\lambda^{-1}(\mathrm{F}_\lambda).
  \tag{8.3.2.1}
\label{IV.8.3.2.1}
\end{equation}
Suponhamos satisfeitas as seguintes condições:
\begin{enumerate}[label=(\roman*)]
  \item Para todo $\lambda$, $\mathrm{E}_\lambda$ é pró-construtível e $\mathrm{F}_\lambda$ é ind-construtível \sref{IV.1.9.4}.
  \item Para $\lambda\leqslant\mu$, tem-se
  \[
    \mathrm{E}_\mu\subset u_{\lambda\mu}^{-1}(\mathrm{E}_\lambda),\qquad \mathrm{F}_\mu\supset u_{\lambda\mu}^{-1}(\mathrm{F}_\lambda).
  \]
  \item Existe $\alpha\in\mathrm{L}$ tal que $\mathrm{S}_\alpha$ é quase-compacto (o que implica que $\mathrm{S}_\lambda$ é quase-compacto para $\lambda\geqslant\alpha$).
\end{enumerate}
Então, as seguintes propriedades são equivalentes:
\begin{enumerate}[label=\emph{\alph*)}]
  \item $\mathrm{E}\subset\mathrm{F}$.
  \item Existe $\lambda\geqslant\alpha$ tal que $u_\lambda^{-1}(\mathrm{E}_\lambda)\subset u_\lambda^{-1}(\mathrm{F}_\lambda)$ (e então também se tem $u_\mu^{-1}(\mathrm{E}_\mu)\subset u_\mu^{-1}(\mathrm{F}_\mu)$ para
\oldpage[IV-3]{13}
$\mu\geqslant\lambda$).
  \item Existe $\lambda\geqslant\alpha$ tal que $\mathrm{E}_\lambda\subset\mathrm{F}_\lambda$ (e então também se tem $\mathrm{E}_\mu\subset\mathrm{F}_\mu$ para $\mu\geqslant\lambda$).
\end{enumerate}
As observações entre parênteses de \emph{b)} e \emph{c)} decorrem de (ii).
Ponhamos
\[
  \mathrm{G}_\lambda=\mathrm{E}_\lambda\cap(\mathrm{S}_\lambda-\mathrm{F}_\lambda),\qquad \mathrm{G}=\mathrm{E}\cap(\mathrm{S}-\mathrm{F}).
\]
Então, $\mathrm{G}_\lambda$ é um subconjunto pró-construtível de $\mathrm{S}_\lambda$ \sref{IV.1.9.5},~(ii), e, em virtude de \sref{IV.8.3.2.1} e da condição (ii), tem-se
\[
  \mathrm{G}_\mu\subset u_{\lambda\mu}^{-1}(\mathrm{G}_\lambda)\qquad\text{para }\lambda\leqslant\mu
\]
\[
  \mathrm{G}=\bigcap_\lambda u_\lambda^{-1}(\mathrm{G}_\lambda).
\]
Ficamos, assim, reduzidos a demonstrar o caso particular de \sref{IV.8.3.2} correspondente a $\mathrm{F}_\lambda=\varnothing$ para todo $\lambda$:
\end{theorem}

\begin{corollary}[8.3.3]
\label{IV.8.3.3}
Para todo $\lambda$, seja $\mathrm{E}_\lambda$ um subconjunto pró-construtível de $\mathrm{S}_\lambda$ tal que, para $\lambda\leqslant\mu$, $\mathrm{E}_\mu\subset u_{\lambda\mu}^{-1}(\mathrm{E}_\lambda)$.
Suponhamos que exista $\alpha\in\mathrm{L}$ tal que $\mathrm{S}_\alpha$ seja quase-compacto.
Então, as seguintes condições são equivalentes:
\begin{enumerate}[label=\emph{\alph*)}]
  \item $\mathrm{E}=\bigcap_\lambda u_\lambda^{-1}(\mathrm{E}_\lambda)=\varnothing$.
  \item Existe $\lambda$ tal que $u_\lambda^{-1}(\mathrm{E}_\lambda)=\varnothing$ (e então $u_\mu^{-1}(\mathrm{E}_\mu)=\varnothing$ para $\mu\geqslant\lambda$).
  \item Existe $\lambda$ tal que $\mathrm{E}_\lambda=\varnothing$ (e então $\mathrm{E}_\mu=\varnothing$ para $\mu\geqslant\lambda$).
\end{enumerate}
\end{corollary}

\begin{proof}
É claro que \emph{c)} implica \emph{a)}.
Demonstremos que \emph{a)} implica \emph{b)}: como $\mathrm{S}_\lambda$ é quase-compacto, o mesmo vale para $\mathrm{S}$ \sref{IV.8.2.2}; como $\mathrm{E}_\lambda$ é pró-construtível, o mesmo vale para $u_\lambda^{-1}(\mathrm{E}_\lambda)$ \sref{IV.1.9.5},~(vi); portanto, a família filtrante decrescente de subconjuntos pró-construtíveis $u_\lambda^{-1}(\mathrm{E}_\lambda)$ tem interseção vazia e, por \sref{IV.1.9.9}, um deles é vazio.

Finalmente, demonstremos que \emph{b)} implica \emph{c)}.
Como $\mathrm{S}_\alpha$ é quase-compacto e $\mathrm{L}$ é filtrante, podemos substituir $\mathrm{S}_\alpha$ por um aberto afim; suponhamos então \sref{IV.8.2.2} que $\mathrm{S}$ e os $\mathrm{S}_\lambda$ para $\lambda\geqslant\alpha$ sejam afins; tem-se então \sref{IV.1.9.2.1}, para $\lambda\geqslant\alpha$,
\[
  u_\lambda(\mathrm{S})=\bigcap_{\mu\geqslant\lambda}u_{\lambda\mu}(\mathrm{S}_\mu)
\]
de onde
\[
  \mathrm{E}_\lambda\cap u_\lambda(\mathrm{S})=\bigcap_{\mu\geqslant\lambda}(\mathrm{E}_\lambda\cap u_{\lambda\mu}(\mathrm{S}_\mu)).
\]
Ora, como $u_\lambda$ e os $u_{\lambda\mu}$ são quase-compactos, $u_\lambda(\mathrm{S})$ e $u_{\lambda\mu}(\mathrm{S}_\mu)$ são pró-construtíveis em $\mathrm{S}_\lambda$ \sref{IV.1.9.5},~(viii); portanto, os conjuntos $\mathrm{E}_\lambda\cap u_{\lambda\mu}(\mathrm{S}_\mu)$ para $\mu\geqslant\lambda$ formam uma família filtrante decrescente de subconjuntos pró-construtíveis de $\mathrm{S}_\lambda$.
Como $\mathrm{S}_\lambda$ é quase-compacto, a hipótese \emph{b)} implica que a interseção dessa família é vazia e, por \sref{IV.1.9.9}, um dos conjuntos $\mathrm{E}_\lambda\cap u_{\lambda\mu}(\mathrm{S}_\mu)$ é vazio; portanto, na inclusão $\mathrm{E}_\mu\subset u_{\lambda\mu}^{-1}(\mathrm{E}_\lambda)$, o membro direito é vazio e, por conseguinte, também o membro esquerdo.
\end{proof}

\begin{corollary}[8.3.4]
\label{IV.8.3.4}
Para todo $\lambda$, seja $\mathrm{F}_\lambda$ um subconjunto ind-construtível de $\mathrm{S}_\lambda$ tal que, para $\lambda\leqslant\mu$, se tenha $u_{\lambda\mu}^{-1}(\mathrm{F}_\lambda)\subset\mathrm{F}_\mu$.
Suponhamos que exista $\alpha\in\mathrm{L}$ tal que $\mathrm{S}_\alpha$ seja quase-compacto.
Então, as seguintes condições são equivalentes:
\begin{enumerate}[label=\emph{\alph*)}]
  \item O conjunto $\mathrm{F}=\bigcup_\lambda u_\lambda^{-1}(\mathrm{F}_\lambda)$ é igual a $\mathrm{S}$.
  \item Existe $\lambda$ tal que $u_\lambda^{-1}(\mathrm{F}_\lambda)=\mathrm{S}$ (e então $u_\mu^{-1}(\mathrm{F}_\mu)=\mathrm{S}$ para $\mu\geqslant\lambda$).
\oldpage[IV-3]{14}
  \item Existe $\lambda$ tal que $\mathrm{F}_\lambda=\mathrm{S}_\lambda$ (e então $\mathrm{F}_\mu=\mathrm{S}_\mu$ para $\mu\geqslant\lambda$).
\end{enumerate}
Isto decorre imediatamente de \sref{IV.8.3.3}, passando aos complementares.
\end{corollary}

\begin{corollary}[8.3.5]
\label{IV.8.3.5}
Para todo $\lambda$, sejam $\mathrm{E}_\lambda$, $\mathrm{F}_\lambda$ dois subconjuntos construtíveis de $\mathrm{S}_\lambda$ tais que, para $\mu\geqslant\lambda$, se tenha $\mathrm{E}_\mu\subset u_{\lambda\mu}^{-1}(\mathrm{E}_\lambda)$ e $\mathrm{F}_\mu\supset u_{\lambda\mu}^{-1}(\mathrm{F}_\lambda)$.
Suponhamos que exista $\alpha$ tal que $\mathrm{S}_\alpha$ seja quase-compacto.
Então, para que $\mathrm{E}\subset\mathrm{F}$ \textnormal{(}resp. $\mathrm{E}=\mathrm{F}$\textnormal{)}, é necessário e suficiente que exista $\lambda$ tal que $\mathrm{E}_\lambda\subset\mathrm{F}_\lambda$ \textnormal{(}resp. $\mathrm{E}_\lambda=\mathrm{F}_\lambda$\textnormal{)}, caso em que também se tem $\mathrm{E}_\mu\subset\mathrm{F}_\mu$ \textnormal{(}resp. $\mathrm{E}_\mu=\mathrm{F}_\mu$\textnormal{)} para $\mu\geqslant\lambda$.

Este é um caso particular de \sref{IV.8.3.2}.
\end{corollary}

Em particular:
\begin{corollary}[8.3.6]
\label{IV.8.3.6}
Suponhamos que exista $\alpha$ tal que $\mathrm{S}_\alpha$ seja quase-compacto.
Para que $\mathrm{S}=\varnothing$, é necessário e suficiente que exista $\lambda$ tal que $\mathrm{S}_\lambda=\varnothing$.
\end{corollary}

\begin{corollary}[8.3.7]
\label{IV.8.3.7}
Para todo $\lambda$, tem-se
\begin{equation}
  u_\lambda(\mathrm{S})=\bigcap_{\mu\geqslant\lambda}u_{\lambda\mu}(\mathrm{S}_\mu).
  \tag{8.3.7.1}
\label{IV.8.3.7.1}
\end{equation}
\end{corollary}

\begin{proof}
É claro que o primeiro membro de \sref{IV.8.3.7.1} está contido no segundo.
Seja $s$ um ponto de $\mathrm{S}_\lambda$ e ponhamos $\mathrm{X}_\lambda=\operatorname{Spec}(\boldsymbol{k}(s))$; consideremos o sistema projetivo $(\mathrm{X}_\mu,v_{\mu\nu})$, em que $\mathrm{X}_\mu=\mathrm{X}_\lambda\times_{\mathrm{S}_\lambda}\mathrm{S}_\mu$ e $v_{\mu\nu}=1\times u_{\mu\nu}$ para $\lambda\leqslant\mu\leqslant\nu$; seu limite projetivo é $\mathrm{X}=\mathrm{X}_\lambda\times_{\mathrm{S}_\lambda}\mathrm{S}$, e $v_\lambda=1\times u_\lambda$ é a aplicação canônica $\mathrm{X}\to\mathrm{X}_\lambda$ \sref{IV.8.2.5}.
Se $s\in\bigcap_{\mu\geqslant\lambda}u_{\lambda\mu}(\mathrm{S}_\mu)$, isso implica que $\mathrm{X}_\mu\neq\varnothing$ para todo $\mu\geqslant\lambda$ \sref[I]{I.3.4.8}; decorre então de \sref{IV.8.3.6} que $\mathrm{X}\neq\varnothing$ e, portanto, $s\in u_\lambda(\mathrm{S})$ por \sref[I]{I.3.4.8}.
\end{proof}

\begin{proposition}[8.3.8]
\label{IV.8.3.8}
\begin{enumerate}[label=(\roman*)]
  \item Para que o morfismo $u_\lambda:\mathrm{S}\to\mathrm{S}_\lambda$ seja dominante (resp. sobrejetivo), é necessário e suficiente que, para $\mu\geqslant\lambda$, o morfismo $u_{\lambda\mu}:\mathrm{S}_\mu\to\mathrm{S}_\lambda$ seja dominante (resp. sobrejetivo).
  \item Se, para algum índice $\lambda$, os morfismos $u_{\lambda\mu}:\mathrm{S}_\mu\to\mathrm{S}_\lambda$ forem planos (resp. fielmente planos) para todo $\mu\geqslant\lambda$, então o morfismo $u_\lambda:\mathrm{S}\to\mathrm{S}_\lambda$ será plano (resp. fielmente plano).
  \item Suponhamos que os morfismos $u_\mu:\mathrm{S}\to\mathrm{S}_\mu$ sejam sobrejetivos para $\mu\geqslant\lambda$.
Para que $u_\lambda$ seja um morfismo aberto (resp. universalmente aberto), é necessário e suficiente que, para todo $\mu\geqslant\lambda$, $u_{\lambda\mu}$ seja um morfismo aberto (resp. universalmente aberto).
\end{enumerate}
\end{proposition}

\begin{proof}
\begin{enumerate}[label=(\roman*)]
  \item Como $u_\lambda(\mathrm{S})\subset u_{\lambda\mu}(\mathrm{S}_\mu)$ para todo $\mu\geqslant\lambda$, a necessidade das condições é trivial, e decorre imediatamente de \sref{IV.8.3.7.1} que, se os $u_{\lambda\mu}$ forem sobrejetivos, $u_\lambda$ também o será.
Suponhamos agora que os $u_{\lambda\mu}$ sejam dominantes para $\mu\geqslant\lambda$ e consideremos em $\mathrm{S}_\lambda$ um subconjunto aberto, quase-compacto e não vazio $\mathrm{U}_\lambda$; então, os $\mathrm{U}_\mu=u_{\lambda\mu}^{-1}(\mathrm{U}_\lambda)$ para $\mu\geqslant\lambda$ formam um sistema projetivo cujo limite projetivo é $\mathrm{U}=u_\lambda^{-1}(\mathrm{U}_\lambda)$ \sref{IV.8.2.5}.
Por hipótese, os $\mathrm{U}_\mu$ são todos não vazios, e o mesmo vale para $\mathrm{U}$ por \sref{IV.8.3.6}, o que demonstra que $u_\lambda$ é dominante.

  \item Em virtude de (i), basta considerar o caso em que os $u_{\lambda\mu}$ sejam planos; pode-se então restringir ao caso em que $\mathrm{S}_\lambda$ seja afim e, portanto, também o sejam os $\mathrm{S}_\mu$ para $\mu\geqslant\lambda$ e $\mathrm{S}$, e a afirmação decorre de \sref{IV.2.1.2} e \sref[0]{0.6.2.3}.

  \item Em virtude de \sref{IV.8.2.5} e de \sref[I]{I.3.5.2}, basta tratar o caso dos morfismos abertos.
Como $u_\lambda=u_{\lambda\mu}\circ u_\mu$ e $u_\mu$ é sobrejetivo, sabemos que, se $u_\lambda$ for aberto, $u_{\lambda\mu}$ também o será para $\mu\geqslant\lambda$ (Bourbaki, \emph{Top.\ g\'{e}n.}, chap.~I\textsuperscript{er}, 3\textsuperscript{e}~ed., \textsection\,5, n\textsuperscript{o}\,1, prop.~1).
\oldpage[IV-3]{15}
Reciprocamente, para demonstrar que $u_\lambda$ é aberto quando todos os $u_{\lambda\mu}$ são abertos para $\mu\geqslant\lambda$, basta verificar que, para todo aberto quase-compacto $\mathrm{V}$ de $\mathrm{S}$, $u_\lambda(\mathrm{V})$ é aberto em $\mathrm{S}_\lambda$; existe então $\mu\geqslant\lambda$ tal que $\mathrm{V}=u_\mu^{-1}(\mathrm{V}_\mu)$, em que $\mathrm{V}_\mu$ é aberto em $\mathrm{S}_\mu$ \sref{IV.8.2.11}; como $u_\mu$ é sobrejetivo, tem-se $\mathrm{V}_\mu=u_\mu(\mathrm{V})$, e $u_\lambda(\mathrm{V})=u_{\lambda\mu}(u_\mu(\mathrm{V}))$ é, portanto, aberto por hipótese.
\end{enumerate}

Observemos que pode ocorrer que todos os $u_{\lambda\mu}$ sejam abertos sem que $u_\lambda$ seja aberto, quando os $u_\mu$ não são sobrejetivos.
Com efeito, obtém-se um exemplo considerando um anel íntegro $\mathrm{A}$ que não seja um corpo e seu corpo de frações $\mathrm{K}$, que é o limite indutivo dos $\mathrm{A}_f$, em que $f$ percorre $\mathrm{A}-\{0\}$; tomando $\mathrm{S}_1=\operatorname{Spec}(\mathrm{A})$, $\mathrm{S}_f=\operatorname{Spec}(\mathrm{A}_f)$ e $\mathrm{S}=\varprojlim\mathrm{S}_f=\operatorname{Spec}(\mathrm{K})$, o morfismo $\mathrm{S}\to\mathrm{S}_1$ não é aberto, embora os morfismos $\mathrm{S}_f\to\mathrm{S}_1$ sejam abertos.
\end{proof}

\begin{env}[8.3.9]
\label{IV.8.3.9}
Para todo pré-esquema $\mathrm{X}$, denotaremos, como de costume, por $\mathfrak{P}(\mathrm{X})$ o conjunto de todos os subconjuntos do conjunto subjacente de $\mathrm{X}$, e por $\mathfrak{C}(\mathrm{X})$ (resp. $\mathfrak{D}_c(\mathrm{X})$, $\mathfrak{F}_c(\mathrm{X})$, $\mathfrak{LF}_c(\mathrm{X})$) o conjunto dos subconjuntos construtíveis (resp. construtíveis e abertos, construtíveis e fechados, e construtíveis e localmente fechados) de $\mathrm{X}$.
É claro que $(\mathfrak{P}(\mathrm{S}_\lambda),u_{\lambda\mu}^{-1})$ é um \emph{sistema indutivo de conjuntos} e que as aplicações $u_\lambda^{-1}:\mathfrak{P}(\mathrm{S}_\lambda)\to\mathfrak{P}(\mathrm{S})$ formam um sistema indutivo de aplicações; portanto, passando ao limite indutivo, obtém-se uma aplicação canônica
\begin{equation}
  u_\mathfrak{P}:\varinjlim\mathfrak{P}(\mathrm{S}_\lambda)\to\mathfrak{P}(\mathrm{S}).
  \tag{8.3.9.1}
\label{IV.8.3.9.1}
\end{equation}

Além disso, decorre de \sref{IV.1.8.2} que $u_{\lambda\mu}^{-1}$ envia $\mathfrak{C}(\mathrm{S}_\lambda)$ (resp. $\mathfrak{D}_c(\mathrm{S}_\lambda)$, $\mathfrak{F}_c(\mathrm{S}_\lambda)$, $\mathfrak{LF}_c(\mathrm{S}_\lambda)$) em $\mathfrak{C}(\mathrm{S}_\mu)$ (resp. $\mathfrak{D}_c(\mathrm{S}_\mu)$, $\mathfrak{F}_c(\mathrm{S}_\mu)$, $\mathfrak{LF}_c(\mathrm{S}_\mu)$), e que $u_\lambda^{-1}$ envia $\mathfrak{C}(\mathrm{S}_\lambda)$ (resp. $\mathfrak{D}_c(\mathrm{S}_\lambda)$, $\mathfrak{F}_c(\mathrm{S}_\lambda)$, $\mathfrak{LF}_c(\mathrm{S}_\lambda)$) em $\mathfrak{C}(\mathrm{S})$ (resp. $\mathfrak{D}_c(\mathrm{S})$, $\mathfrak{F}_c(\mathrm{S})$, $\mathfrak{LF}_c(\mathrm{S})$).
Temos, pois, por restrição de \sref{IV.8.3.9.1}, as aplicações canônicas
\begin{align}
  &\varinjlim\mathfrak{C}(\mathrm{S}_\lambda)\to\mathfrak{C}(\mathrm{S}) \tag{8.3.9.2}\label{IV.8.3.9.2}\\
  &\varinjlim\mathfrak{D}_c(\mathrm{S}_\lambda)\to\mathfrak{D}_c(\mathrm{S}) \tag{8.3.9.3}\label{IV.8.3.9.3}\\
  &\varinjlim\mathfrak{F}_c(\mathrm{S}_\lambda)\to\mathfrak{F}_c(\mathrm{S}) \tag{8.3.9.4}\label{IV.8.3.9.4}\\
  &\varinjlim\mathfrak{LF}_c(\mathrm{S}_\lambda)\to\mathfrak{LF}_c(\mathrm{S}). \tag{8.3.9.5}
\label{IV.8.3.9.5}
\end{align}
\end{env}

\begin{env}[8.3.10]
\label{IV.8.3.10}
Seja $g_\alpha:\mathrm{X}_\alpha\to\mathrm{S}_\alpha$ um morfismo; com as notações de \sref{IV.8.2.5}, tem-se, por um lado, uma aplicação canônica $v_\mathfrak{P}:\varinjlim\mathfrak{P}(\mathrm{X}_\lambda)\to\mathfrak{P}(\mathrm{X})$; por outro lado, têm-se morfismos de projeção $g_\lambda:\mathrm{X}_\lambda\to\mathrm{S}_\lambda$ para todo $\lambda\geqslant\alpha$ e um morfismo de projeção $g:\mathrm{X}\to\mathrm{S}$.
É claro que os $g_\lambda^{-1}:\mathfrak{P}(\mathrm{S}_\lambda)\to\mathfrak{P}(\mathrm{X}_\lambda)$ formam um sistema indutivo de aplicações e que os diagramas
\[
\xymatrix{
  \mathfrak{P}(\mathrm{S}_\lambda) \ar[r]^{u_\lambda^{-1}} \ar[d]_{g_\lambda^{-1}} & \mathfrak{P}(\mathrm{S}) \ar[d]^{g^{-1}} \\
  \mathfrak{P}(\mathrm{X}_\lambda) \ar[r]_{v_\lambda^{-1}} & \mathfrak{P}(\mathrm{X})
}
\]
\oldpage[IV-3]{16}
são comutativos; passando ao limite indutivo, deduz-se daí um diagrama comutativo
\begin{equation}
\xymatrix{
  \varinjlim\mathfrak{P}(\mathrm{S}_\lambda) \ar[r]^{u_\mathfrak{P}} \ar[d]_{\varinjlim g_\lambda^{-1}} & \mathfrak{P}(\mathrm{S}) \ar[d]^{g^{-1}} \\
  \varinjlim\mathfrak{P}(\mathrm{X}_\lambda) \ar[r]_{v_\mathfrak{P}} & \mathfrak{P}(\mathrm{X})
}
\tag{8.3.10.1}
\label{IV.8.3.10.1}
\end{equation}
e decorre de \sref{IV.1.8.2} que se têm diagramas análogos substituindo $\mathfrak{P}$ por $\mathfrak{C}$, $\mathfrak{D}_c$, $\mathfrak{F}_c$ ou $\mathfrak{LF}_c$.
\end{env}

Decorre de \sref{IV.8.3.5} que, sob a hipótese de que, para algum $\alpha\in\mathrm{L}$, $\mathrm{S}_\alpha$ seja \emph{quase-compacto}, a aplicação canônica \sref{IV.8.3.9.2} é \emph{injetiva}.
Além disso:

\begin{theorem}[8.3.11]
\label{IV.8.3.11}
Suponhamos que exista $\alpha\in\mathrm{L}$ tal que $\mathrm{S}_\alpha$ seja quase-compacto e quase-separado.
Então as aplicações canônicas \sref{IV.8.3.9.2}, \sref{IV.8.3.9.3}, \sref{IV.8.3.9.4} e \sref{IV.8.3.9.5} são \emph{bijetivas}.
\end{theorem}
\begin{proof}
Em virtude da observação precedente, resta demonstrar que essas aplicações são \emph{sobrejetivas}; como todo subconjunto construtível de $\mathrm{S}$ é união finita de conjuntos da forma $\mathrm{U}\cap\complement\mathrm{V}$, em que $\mathrm{U}$ e $\mathrm{V}$ são abertos e construtíveis, basta demonstrar que \sref{IV.8.3.9.3} é sobrejetiva para que o mesmo valha para \sref{IV.8.3.9.2} (e também para \sref{IV.8.3.9.4}, passando aos complementares).
Ora, como os morfismos $\mathrm{S}_\lambda\to\mathrm{S}_\alpha$ e $\mathrm{S}\to\mathrm{S}_\alpha$ são afins e, portanto, separados, os $\mathrm{S}_\lambda$ para $\lambda\geqslant\alpha$ e $\mathrm{S}$ são quase-compactos e quase-separados \sref{IV.1.2.2}, e sabe-se que os subconjuntos abertos construtíveis de tal pré-esquema não são outros senão os abertos quase-compactos \sref{IV.1.8.1}.
A conclusão decorre então de \sref{IV.8.2.11}, salvo para a aplicação \sref{IV.8.3.9.5}.
Para demonstrar que esta última aplicação é sobrejetiva, consideremos um subconjunto localmente fechado e construtível $\mathrm{Z}$ de $\mathrm{S}$; portanto, $\mathrm{Z}$ é quase-compacto ($\mathbf{0}_{\mathrm{III}}$,~9.1.4).
Como, por hipótese, todo ponto $x\in\mathrm{Z}$ admite uma vizinhança aberta quase-compacta $\mathrm{V}_x$ em $\mathrm{S}$ tal que $\mathrm{Z}\cap\mathrm{V}_x$ é fechado em $\mathrm{V}_x$, podemos cobrir $\mathrm{Z}$ por um número finito de $\mathrm{V}_x$; em outras palavras, existe um aberto quase-compacto $\mathrm{U}$ que contém $\mathrm{Z}$ tal que $\mathrm{Z}$ é fechado em $\mathrm{U}$; como $\mathrm{Z}$ é construtível em $\mathrm{S}$, também é construtível em $\mathrm{U}$ ($\mathbf{0}_{\mathrm{III}}$,~9.1.8).
Sabemos por \sref{IV.8.2.11} que existem $\lambda$ e um aberto quase-compacto $\mathrm{U}_\lambda$ de $\mathrm{S}_\lambda$ tais que $\mathrm{U}=u_\lambda^{-1}(\mathrm{U}_\lambda)$.
Aplicando a $\mathrm{U}$ (que é o limite projetivo dos $\mathrm{U}_\mu=u_{\lambda\mu}^{-1}(\mathrm{U}_\lambda)$ para $\mu\geqslant\lambda$) o fato de que a aplicação \sref{IV.8.3.9.4} é sobrejetiva, vê-se que existem $\mu\geqslant\lambda$ e um subconjunto fechado construtível $\mathrm{Z}_\mu$ em $\mathrm{U}_\mu$ tais que $\mathrm{Z}=u_\mu^{-1}(\mathrm{Z}_\mu)$.
Mas, como a imersão canônica $\mathrm{U}_\mu\to\mathrm{S}_\mu$ é quase-compacta por hipótese \sref{IV.1.2.7}, ela é de apresentação finita \sref{IV.1.6.2},~(i), e $\mathrm{Z}_\mu$ também é um subconjunto construtível \emph{de} $\mathrm{S}_\mu$ em virtude de \sref{IV.1.8.4} e \sref{IV.1.8.1}; como $\mathrm{Z}_\mu$ é evidentemente localmente fechado em $\mathrm{S}_\mu$, isso conclui a demonstração.
\end{proof}

\begin{corollary}[8.3.12]
\label{IV.8.3.12}
Suponhamos que exista $\alpha$ tal que $\mathrm{S}_\alpha$ seja quase-compacto e, para todo $\lambda$, seja $\mathrm{Z}_\lambda$ um subconjunto construtível de $\mathrm{S}_\lambda$ tal que $\mathrm{Z}_\mu=u_{\lambda\mu}^{-1}(\mathrm{Z}_\lambda)$ para $\mu\geqslant\lambda$.
Se $\mathrm{Z}=u_\lambda^{-1}(\mathrm{Z}_\lambda)$,
\oldpage[IV-3]{17}
então, para que $\mathrm{Z}$ seja aberto (resp. fechado, resp. localmente fechado) em $\mathrm{S}$, é necessário e suficiente que exista $\lambda\geqslant\alpha$ tal que $\mathrm{Z}_\lambda$ também o seja em $\mathrm{S}_\lambda$.
\end{corollary}

\begin{proof}
Cubra $\mathrm{S}_\alpha$ com um número finito de abertos afins $\mathrm{U}_\alpha^{(i)}$; então os $\mathrm{U}^{(i)}=u_\alpha^{-1}(\mathrm{U}_\alpha^{(i)})$ formam um recobrimento aberto por afins de $\mathrm{S}$, e, para que $\mathrm{Z}$ seja aberto (resp. fechado, resp. localmente fechado) em $\mathrm{S}$, é necessário e suficiente que cada $\mathrm{Z}\cap\mathrm{U}^{(i)}$ também o seja em $\mathrm{U}^{(i)}$.
Como $\mathrm{L}$ é filtrante e cada $\mathrm{Z}\cap\mathrm{U}^{(i)}$ é construtível em $\mathrm{U}^{(i)}$ ($\mathbf{0}_{\mathrm{III}}$,~9.1.8), podemos restringir-nos a demonstrar o corolário quando $\mathrm{S}_\alpha$ é afim e, portanto, quase-compacto e quase-separado; mas então isso decorre de \sref{IV.8.3.11}.
\end{proof}

\begin{proposition}[8.3.13]
\label{IV.8.3.13}
Suponhamos que os morfismos $u_{\lambda\mu}:\mathrm{S}_\mu\to\mathrm{S}_\lambda$ sejam planos para $\lambda\leqslant\mu$ e que exista $\alpha$ tal que $\mathrm{S}_\alpha$ seja quase-compacto.
Para todo $\lambda$, sejam $\mathrm{Z}'_\lambda$, $\mathrm{Z}''_\lambda$ dois subconjuntos pró-construtíveis de $\mathrm{S}_\lambda$ tais que $\mathrm{Z}'_\mu=u_{\lambda\mu}^{-1}(\mathrm{Z}'_\lambda)$ e $\mathrm{Z}''_\mu=u_{\lambda\mu}^{-1}(\mathrm{Z}''_\lambda)$ para $\mu\geqslant\lambda$; suponhamos, além disso, que $\overline{\mathrm{Z}'_\alpha}$ seja construtível em $\mathrm{S}_\alpha$.
Ponhamos $\mathrm{Z}'=u_\lambda^{-1}(\mathrm{Z}'_\lambda)$, $\mathrm{Z}''=u_\lambda^{-1}(\mathrm{Z}''_\lambda)$; para que $\mathrm{Z}''\subset\overline{\mathrm{Z}'}$, é necessário e suficiente que exista $\lambda\geqslant\alpha$ tal que $\mathrm{Z}''_\lambda\subset\overline{\mathrm{Z}'_\lambda}$.
\end{proposition}

\begin{proof}
Com efeito, sabe-se que $u_\lambda:\mathrm{S}\to\mathrm{S}_\lambda$ também é um morfismo plano para todo $\lambda$ \sref{IV.8.3.8}; como $\mathrm{Z}'_\lambda$ é pró-construtível, o fecho de $\mathrm{Z}'_\mu$ para $\mu\geqslant\lambda$ (resp. de $\mathrm{Z}'$) em $\mathrm{S}_\mu$ (resp. $\mathrm{S}$) é igual a $u_{\lambda\mu}^{-1}(\overline{\mathrm{Z}'_\lambda})$ (resp. $u_\lambda^{-1}(\overline{\mathrm{Z}'_\lambda})$) \sref{IV.2.3.10}.
Como $u_{\lambda\mu}^{-1}(\overline{\mathrm{Z}'_\lambda})$ e $u_\lambda^{-1}(\overline{\mathrm{Z}'_\lambda})$ são construtíveis \sref{IV.1.8.2}, a conclusão decorre de \sref{IV.8.3.2}.
\end{proof}

\subsection{Critérios de irredutibilidade e conexidade para os limites projetivos de pré-esquemas}
\label{subsection:IV.8.4}

\begin{proposition}[8.4.1]
\label{IV.8.4.1}
Suponhamos que exista um índice $\alpha$ tal que $\mathrm{S}_\alpha$ seja quase-compacto.
\begin{enumerate}[label=(\roman*)]
  \item Se $\mathrm{S}$ não for irredutível e se, além disso, o espaço subjacente de $\mathrm{S}$ for noetheriano e $\mathrm{S}_\alpha$ for quase-separado, existe $\lambda\geqslant\alpha$ tal que, para $\mu\geqslant\lambda$, $\mathrm{S}_\mu$ não é irredutível.
  \item Se $\mathrm{S}$ não for conexo, existe $\lambda\geqslant\alpha$ tal que, para $\mu\geqslant\lambda$, $\mathrm{S}_\mu$ não é conexo.
\end{enumerate}
\end{proposition}

\begin{proof}
Suponhamos que $\mathrm{S}$ seja a união de dois subconjuntos fechados próprios $\mathrm{S}'$, $\mathrm{S}''$ de $\mathrm{S}$ (resp.\ de dois subconjuntos fechados não vazios e disjuntos).
No caso (i), $\mathrm{S}'$ e $\mathrm{S}''$ são construtíveis, pois o espaço $\mathrm{S}$ é noetheriano.
Em virtude de \sref{IV.8.3.11}, existem $\lambda\geqslant\alpha$ e dois subconjuntos fechados construtíveis $\mathrm{S}'_\lambda$, $\mathrm{S}''_\lambda$ de $\mathrm{S}_\lambda$ tais que $\mathrm{S}'=u_\lambda^{-1}(\mathrm{S}'_\lambda)$, $\mathrm{S}''=u_\lambda^{-1}(\mathrm{S}''_\lambda)$; como $\mathrm{S}=\mathrm{S}'\cup\mathrm{S}''$, decorre de \sref{IV.8.3.11} que se pode supor $\mathrm{S}_\lambda=\mathrm{S}'_\lambda\cup\mathrm{S}''_\lambda$; como $\mathrm{S}'_\lambda$ e $\mathrm{S}''_\lambda$ são distintos de $\mathrm{S}_\lambda$, isso demonstra que $\mathrm{S}_\lambda$ não é irredutível.

No caso (ii), $\mathrm{S}'$ e $\mathrm{S}''$ são abertos e quase-compactos; portanto, em virtude de \sref{IV.8.2.11}, existem $\lambda\geqslant\alpha$ e dois subconjuntos abertos quase-compactos $\mathrm{S}'_\lambda$, $\mathrm{S}''_\lambda$ de $\mathrm{S}_\lambda$ tais que $\mathrm{S}'=u_\lambda^{-1}(\mathrm{S}'_\lambda)$, $\mathrm{S}''=u_\lambda^{-1}(\mathrm{S}''_\lambda)$.
Além disso, como $\mathrm{S}'$ e $\mathrm{S}''$ são abertos e fechados em $\mathrm{S}$, ambos são pró-construtíveis e ind-construtíveis \sref{IV.1.9.6} e, portanto, construtíveis \sref{IV.1.9.11}; decorre então de \sref{IV.8.3.5} que se pode supor $\lambda$ escolhido de modo que $\mathrm{S}_\lambda=\mathrm{S}'_\lambda\cup\mathrm{S}''_\lambda$ e $\mathrm{S}'_\lambda\cap\mathrm{S}''_\lambda=\varnothing$, o que mostra que $\mathrm{S}_\lambda$ não é conexo.
\end{proof}

\begin{proposition}[8.4.2]
\label{IV.8.4.2}
Suponhamos que o espaço subjacente de $\mathrm{S}$ seja noetheriano e que se satisfaça uma das duas condições seguintes:
\begin{enumerate}[label=\emph{\alph*)}]
  \item Para $\lambda\leqslant\mu$, $u_{\lambda\mu}:\mathrm{S}_\mu\to\mathrm{S}_\lambda$ é dominante, e existe $\alpha$ tal que $\mathrm{S}_\alpha$ é quase-compacto.
\oldpage[IV-3]{18}
  \item Existe $\alpha$ tal que o espaço subjacente de $\mathrm{S}_\alpha$ é noetheriano e, para $\mu\geqslant\lambda$, $u_{\lambda\mu}$ é um homeomorfismo de $\mathrm{S}_\mu$ sobre um subespaço de $\mathrm{S}_\lambda$.
\end{enumerate}
Nessas condições, existe $\lambda$ tal que, para todo $\mu\geqslant\lambda$:
\begin{enumerate}[label=(\roman*)]
  \item Para toda componente irredutível $\mathrm{Y}_i$ de $\mathrm{S}$ ($1\leqslant i\leqslant m$), $\overline{u_\mu(\mathrm{Y}_i)}$ é uma componente irredutível de $\mathrm{S}_\mu$, e a aplicação $\mathrm{Y}_i\mapsto\overline{u_\mu(\mathrm{Y}_i)}$ é uma bijeção do conjunto das componentes irredutíveis de $\mathrm{S}$ sobre o conjunto das componentes irredutíveis de $\mathrm{S}_\mu$.
  \item Para toda componente conexa $\mathrm{C}_j$ de $\mathrm{S}$ ($1\leqslant j\leqslant n$), $\overline{u_\mu(\mathrm{C}_j)}$ é uma componente conexa de $\mathrm{S}_\mu$, e a aplicação $\mathrm{C}_j\mapsto\overline{u_\mu(\mathrm{C}_j)}$ é uma bijeção do conjunto das componentes conexas de $\mathrm{S}$ sobre o conjunto das componentes conexas de $\mathrm{S}_\mu$.
\end{enumerate}
\end{proposition}

\begin{proof}
Demonstremos primeiro o

\begin{lemma}[8.4.2.1]
\label{IV.8.4.2.1}
Sob a condição \emph{a)} ou \emph{b)} de \sref{IV.8.4.2}, existe $\lambda$ tal que, para $\mu\geqslant\lambda$, $u_\mu:\mathrm{S}\to\mathrm{S}_\mu$ é dominante.
\end{lemma}

No caso \emph{a)}, isso já foi demonstrado sem supor noetheriano o espaço $\mathrm{S}$ \sref{IV.8.3.8},~(i).
No caso \emph{b)}, ponhamos $\mathrm{Z}_\alpha=\overline{u_\alpha(\mathrm{S})}$; como subconjunto fechado do espaço noetheriano $\mathrm{S}_\alpha$, $\mathrm{Z}_\alpha$ é construtível e, como $u_\alpha^{-1}(\mathrm{Z}_\alpha)=\mathrm{S}$, decorre de \sref{IV.8.3.5} que existe $\lambda\geqslant\alpha$ tal que, para $\mu\geqslant\lambda$, $\mathrm{Z}_\mu=u_{\alpha\mu}^{-1}(\mathrm{Z}_\alpha)=\mathrm{S}_\mu$.
Mas, como $u_{\alpha\mu}$ é um homeomorfismo de $\mathrm{S}_\mu$ sobre um subespaço de $\mathrm{S}_\alpha$ e a aplicação composta $\mathrm{S}\to\mathrm{Z}_\mu\to\mathrm{Z}_\alpha$ é dominante, o mesmo vale para $\mathrm{S}\to\mathrm{Z}_\mu=\mathrm{S}_\mu$.

Demonstrado este lema, pode-se supor que, para todo $\lambda$, $u_\lambda$ seja um morfismo dominante.
\begin{enumerate}[label=(\roman*)]
  \item Cada $\mathrm{S}_\lambda$ é a união dos $\overline{u_\lambda(\mathrm{Y}_i)}$, que são irredutíveis.
Por outro lado, se $\mathrm{U}_i$ é o aberto de $\mathrm{S}$ complementar da união dos $\mathrm{Y}_j$ de índice $j\neq i$ ($1\leqslant i\leqslant m$), os $\mathrm{U}_i$ são dois a dois disjuntos e $\mathrm{Y}_i=\overline{\mathrm{U}_i}$ \sref[0]{0.2.1.6}.
Como o espaço subjacente de $\mathrm{S}$ é noetheriano, os $\mathrm{U}_i$ são quase-compactos; portanto, existem um índice $\lambda$ e abertos $\mathrm{U}_{i\lambda}$ de $\mathrm{S}_\lambda$ tais que $\mathrm{U}_i=u_\lambda^{-1}(\mathrm{U}_{i\lambda})$ para $1\leqslant i\leqslant m$ \sref{IV.8.2.11}.
Daí decorre que, se pusermos $\mathrm{U}_{i\mu}=u_{\lambda\mu}^{-1}(\mathrm{U}_{i\lambda})$ para $\lambda\leqslant\mu$, os $\mathrm{U}_{i\mu}$ são dois a dois disjuntos, pois os $\mathrm{U}_i=u_\mu^{-1}(\mathrm{U}_{i\mu})$ o são e $u_\mu$ é dominante.
Consequentemente, nenhum dos fechos $\overline{\mathrm{U}_{i\mu}}$ está contido em outro, e $u_\mu(\mathrm{U}_i)$ é denso em $\mathrm{U}_{i\mu}$, pois $u_\mu$ é dominante; daí $\overline{\mathrm{U}_{i\mu}}=\overline{u_\mu(\mathrm{Y}_i)}$, o que demonstra que os $\overline{\mathrm{U}_{i\mu}}$ são as componentes irredutíveis de $\mathrm{S}_\mu$ \sref[0]{0.2.1.7} e conclui a demonstração.

  \item Como o espaço subjacente de $\mathrm{S}$ é noetheriano, os $\mathrm{C}_j$ são abertos e fechados em $\mathrm{S}$ e quase-compactos; o mesmo raciocínio que em (i) mostra que existem $\lambda$ e abertos $\mathrm{V}_{j\lambda}$ de $\mathrm{S}_\lambda$ tais que $\mathrm{C}_j=u_\lambda^{-1}(\mathrm{V}_{j\lambda})$ para $1\leqslant j\leqslant n$.
Vê-se igualmente, como em (i), que, se pusermos $\mathrm{V}_{j\mu}=u_{\lambda\mu}^{-1}(\mathrm{V}_{j\lambda})$ para $\mu\geqslant\lambda$, os $\mathrm{V}_{j\mu}$ são dois a dois disjuntos e $u_\mu(\mathrm{C}_j)$ é denso em $\mathrm{V}_{j\mu}$; isso implica que $\mathrm{V}_{j\mu}$ é conexo.
Além disso, decorre de \sref{IV.8.3.4} que, para $\mu$ suficientemente grande, a união dos $\mathrm{V}_{j\mu}$ é $\mathrm{S}_\mu$, pois todo aberto de um pré-esquema é ind-construtível \sref{IV.1.9.6}.
Os $\mathrm{V}_{j\mu}$ são, portanto, as componentes conexas de $\mathrm{S}_\mu$, e isso conclui a demonstração.
\end{enumerate}
\end{proof}

Observemos que, se os morfismos $u_{\lambda\mu}$ são \emph{imersões}, satisfazem, em particular, a condição \emph{b)} de \sref{IV.8.4.2}.

\oldpage[IV-3]{19}
\begin{corollary}[8.4.3]
\label{IV.8.4.3}
Suponhamos que se satisfaça uma das condições \emph{a)}, \emph{b)} de \sref{IV.8.4.2}, sendo noetheriano o espaço subjacente de $\mathrm{S}$; então, para que $\mathrm{S}$ seja irredutível, é necessário e suficiente que exista $\lambda$ tal que $\mathrm{S}_\mu$ o seja para todo $\mu\geqslant\lambda$.
\end{corollary}

\begin{proposition}[8.4.4]
\label{IV.8.4.4}
Suponhamos que exista $\alpha$ tal que $\mathrm{S}_\alpha$ seja quase-compacto e que, para $\lambda\leqslant\mu$, $u_{\lambda\mu}:\mathrm{S}_\mu\to\mathrm{S}_\lambda$ seja dominante.
Então, para que $\mathrm{S}$ seja conexo, é necessário e suficiente que exista $\lambda$ tal que $\mathrm{S}_\mu$ o seja para todo $\mu\geqslant\lambda$.
\end{proposition}

\begin{proof}
A condição é suficiente em virtude de \sref{IV.8.4.1}; por outro lado, vimos em \sref{IV.8.3.8},~(i), que $u_\lambda:\mathrm{S}\to\mathrm{S}_\lambda$ é dominante para $\lambda$ suficientemente grande; portanto, se $\mathrm{S}$ é conexo, o mesmo ocorre com $\mathrm{S}_\lambda$, pois $u_\lambda(\mathrm{S})$ é denso em $\mathrm{S}_\lambda$ e é conexo.
\end{proof}

\begin{corollary}[8.4.5]
\label{IV.8.4.5}
Seja $k$ um corpo e $\mathrm{X}$ um $k$-pré-esquema quase-compacto.
Para que $\mathrm{X}$ seja geometricamente conexo \sref{IV.4.5.2}, é necessário e suficiente que, para toda extensão finita e separável $\mathrm{K}$ de $k$, $\mathrm{X}\otimes_k\mathrm{K}$ seja conexo.
\end{corollary}

\begin{proof}
A condição é trivialmente necessária.
Para ver que é suficiente, devemos demonstrar que, se $\Omega$ é um fecho algébrico de $k$, $\mathrm{X}\otimes_k\Omega$ é conexo \sref{IV.4.5.1}.
Ora, $\Omega$ é um limite indutivo filtrante das subextensões finitas $k'$ de $k$ e, para $k\subset k'\subset k''\subset\Omega$, o morfismo $\mathrm{X}\otimes_k k''\to\mathrm{X}\otimes_k k'$ é sobrejetivo.
Assim, em virtude de \sref{IV.8.4.4}, basta demonstrar que $\mathrm{X}\otimes_k k'$ é conexo para toda extensão finita $k'$ de $k$.
Mas, se $\mathrm{K}$ é a maior extensão separável contida em $k'$, o morfismo $\mathrm{X}\otimes_k k'\to\mathrm{X}\otimes_k\mathrm{K}$ é finito, sobrejetivo e radicial e, portanto, \sref{IV.2.4.5} um homeomorfismo; como $\mathrm{X}\otimes_k\mathrm{K}$ é conexo por hipótese, o mesmo vale para $\mathrm{X}\otimes_k k'$.
\end{proof}

\begin{remark}[8.4.6]
\label{IV.8.4.6}
\begin{enumerate}[label=(\roman*)]
  \item A demonstração de \sref{IV.8.4.2} mostra que a conclusão desta proposição permanece válida se supusermos que o espaço subjacente de $\mathrm{S}$ é noetheriano, que existe $\alpha$ tal que $\mathrm{S}_\alpha$ é quase-compacto e, finalmente, que os $u_\lambda:\mathrm{S}\to\mathrm{S}_\lambda$ são dominantes.
  \item Por outro lado, a conclusão de \sref{IV.8.4.2} pode falhar quando os $u_\lambda$ não são dominantes para $\lambda$ suficientemente grande, mesmo quando os $\mathrm{S}_\lambda$ e $\mathrm{S}$ são noetherianos, como mostra o exemplo seguinte.
Tomemos $\mathbf{N}$ como conjunto de índices, sendo todos os $\mathrm{S}_n$ iguais a $\operatorname{Spec}(\mathrm{A}\times\mathrm{K})=\operatorname{Spec}(\mathrm{A})\amalg\operatorname{Spec}(\mathrm{K})$, em que $\mathrm{K}$ é um corpo, $\mathrm{A}$ uma $\mathrm{K}$-álgebra qualquer, e todos os morfismos $u_{n,n+1}$ iguais ao mesmo morfismo correspondente ao homomorfismo $(x,y)\mapsto(j(y),y)$ de $\mathrm{A}\times\mathrm{K}$ em si mesmo, em que $j:\mathrm{K}\to\mathrm{A}$ é o homomorfismo canônico.
Verifica-se facilmente que o limite indutivo deste sistema de anéis é $\mathrm{K}$, sendo o homomorfismo canônico $u_n$ correspondente à segunda projeção $\mathrm{A}\times\mathrm{K}\to\mathrm{K}$.
Vê-se assim que $\mathrm{S}=\operatorname{Spec}(\mathrm{K})$ é irredutível, embora nenhum dos $\mathrm{S}_n$ seja conexo.
\end{enumerate}
\end{remark}

\subsection{Módulos de apresentação finita sobre um limite projetivo de pré-esquemas}
\label{subsection:IV.8.5}

\begin{env}[8.5.1]
\label{IV.8.5.1}
Conservamos sempre as notações de \sref{IV.8.2.2}; além disso, restringir-nos-emos ao caso em que $\mathrm{S}_0$ seja um dos $\mathrm{S}_\lambda$, ao qual sempre se pode reduzir.

\emph{Sempre que, nesta seção, considerarmos uma família $(\sh{F}_\lambda)$, em que, para todo $\lambda\in\mathrm{L}$, $\sh{F}_\lambda$ é um $\sh{O}_{\mathrm{S}_\lambda}$-Módulo, ficará subentendido que essa família satisfaz a condição}
\begin{equation}
  \sh{F}_\mu=u_{\lambda\mu}^*(\sh{F}_\lambda) \qquad\text{para }\lambda\leqslant\mu.
  \tag{8.5.1.1}
\label{IV.8.5.1.1}
\end{equation}
\oldpage[IV-3]{20}
Poremos então
\begin{equation}
  \sh{F}=u_\lambda^*(\sh{F}_\lambda)
  \label{IV.8.5.1.2}
\tag{8.5.1.2}
\end{equation}
que é um $\sh{O}_\mathrm{S}$-Módulo que não depende do índice $\lambda\in\mathrm{L}$, em virtude da hipótese \sref{IV.8.5.1.1}.

Sejam $(\sh{F}_\lambda)$ e $(\sh{G}_\lambda)$ duas famílias desse tipo de $\sh{O}_{\mathrm{S}_\lambda}$-Módulos.
É claro que as aplicações $f_\lambda\mapsto u_{\lambda\mu}^*(f_\lambda)$ de $\operatorname{Hom}_{\sh{O}_{\mathrm{S}_\lambda}}(\sh{F}_\lambda,\sh{G}_\lambda)$ em $\operatorname{Hom}_{\sh{O}_{\mathrm{S}_\mu}}(\sh{F}_\mu,\sh{G}_\mu)$ definem um \emph{sistema indutivo} de grupos abelianos $(\operatorname{Hom}_{\sh{O}_{\mathrm{S}_\lambda}}(\sh{F}_\lambda,\sh{G}_\lambda),u_{\lambda\mu}^*)$, e que as aplicações $f_\lambda\mapsto u_\lambda^*(f_\lambda)$ formam um sistema indutivo de homomorfismos de grupos abelianos
\[
  u_\lambda^*:\operatorname{Hom}_{\sh{O}_{\mathrm{S}_\lambda}}(\sh{F}_\lambda,\sh{G}_\lambda)\to\operatorname{Hom}_\mathrm{S}(\sh{F},\sh{G});
\]
daí, passando ao limite indutivo, obtém-se um homomorfismo canônico de grupos abelianos
\begin{equation}
  u_{\sh{F},\sh{G}}^*:\varinjlim\operatorname{Hom}_{\sh{O}_{\mathrm{S}_\lambda}}(\sh{F}_\lambda,\sh{G}_\lambda)\to\operatorname{Hom}_\mathrm{S}(\sh{F},\sh{G}).
  \tag{8.5.1.3}
\label{IV.8.5.1.3}
\end{equation}

Observemos que, quando $\sh{F}_\lambda=\sh{O}_{\mathrm{S}_\lambda}$, a condição \sref{IV.8.5.1.1} é satisfeita e se tem $\sh{F}=\sh{O}_\mathrm{S}$; portanto, o homomorfismo \sref{IV.8.5.1.3} fornece \sref[0]{0.5.1.1} um homomorfismo canônico de grupos abelianos
\begin{equation}
  u_{\sh{G}}:\varinjlim\Gamma(\mathrm{S}_\lambda,\sh{G}_\lambda)\to\Gamma(\mathrm{S},\sh{G}).
  \tag{8.5.1.4}
\label{IV.8.5.1.4}
\end{equation}
\end{env}

\begin{theorem}[8.5.2]
\label{IV.8.5.2}
\begin{enumerate}[label=(\roman*)]
  \item Suponhamos que $\mathrm{S}_0$ seja quase-compacto (resp.\ quase-compacto e quase-separado) e que, para algum $\lambda\in\mathrm{L}$, $\sh{F}_\lambda$ seja quase-coerente e de tipo finito (resp.\ de apresentação finita), e $\sh{G}_\lambda$ seja quase-coerente.
Então o homomorfismo $u_{\sh{F},\sh{G}}^*$ é injetivo (resp.\ bijetivo).
  \item Suponhamos que $\mathrm{S}_0$ seja quase-compacto e quase-separado.
Para todo $\sh{O}_\mathrm{S}$-Módulo $\sh{F}$ quase-coerente e de apresentação finita, existem $\lambda\in\mathrm{L}$ e um $\sh{O}_{\mathrm{S}_\lambda}$-Módulo $\sh{F}_\lambda$ quase-coerente e de apresentação finita tais que $\sh{F}$ é isomorfo a $u_\lambda^*(\sh{F}_\lambda)$.
\end{enumerate}
\end{theorem}

\begin{proof}
\begin{enumerate}[label=(\roman*)]
  \item Podemos evidentemente restringir-nos ao caso em que $\mathrm{S}_0=\mathrm{S}_\lambda$, pois os morfismos $u_{0\lambda}:\mathrm{S}_\lambda\to\mathrm{S}_0$ são afins e, portanto, quase-compactos e separados.
Consideremos primeiro o caso em que $\mathrm{S}_0=\operatorname{Spec}(\mathrm{A}_0)$ é \emph{afim}.
Então a afirmação (i) equivale ao

\begin{lemma}[8.5.2.1]
\label{IV.8.5.2.1}
Seja $\mathrm{A}_0$ um anel, $(\mathrm{A}_\lambda)_{\lambda\in\mathrm{L}}$ um sistema indutivo de $\mathrm{A}_0$-álgebras, $\mathrm{A}=\varinjlim\mathrm{A}_\lambda$; sejam $\mathrm{M}_0$, $\mathrm{N}_0$ dois $\mathrm{A}_0$-módulos, e ponhamos $\mathrm{M}_\lambda=\mathrm{M}_0\otimes_{\mathrm{A}_0}\mathrm{A}_\lambda$, $\mathrm{N}_\lambda=\mathrm{N}_0\otimes_{\mathrm{A}_0}\mathrm{A}_\lambda$, $\mathrm{M}=\mathrm{M}_0\otimes_{\mathrm{A}_0}\mathrm{A}=\varinjlim\mathrm{M}_\lambda$, $\mathrm{N}=\mathrm{N}_0\otimes_{\mathrm{A}_0}\mathrm{A}=\varinjlim\mathrm{N}_\lambda$.
Se $\mathrm{M}_0$ é de tipo finito (resp.\ de apresentação finita), o homomorfismo canônico
\begin{equation}
  \varinjlim\operatorname{Hom}_{\mathrm{A}_\lambda}(\mathrm{M}_\lambda,\mathrm{N}_\lambda)\to\operatorname{Hom}_\mathrm{A}(\mathrm{M},\mathrm{N})
  \tag{8.5.2.2}
\label{IV.8.5.2.2}
\end{equation}
é injetivo (resp.\ bijetivo).
\end{lemma}

Sabemos, com efeito (Bourbaki, \emph{Alg.}, cap.~II, 3\textsuperscript{a}~ed., \textsection\,5, n\textsuperscript{o}\,1), que existem isomorfismos canônicos funtoriais
\[
  \operatorname{Hom}_{\mathrm{A}_\lambda}(\mathrm{M}_\lambda,\mathrm{N}_\lambda)\xrightarrow{\sim}\operatorname{Hom}_{\mathrm{A}_0}(\mathrm{M}_0,\mathrm{N}_\lambda),\qquad\operatorname{Hom}_\mathrm{A}(\mathrm{M},\mathrm{N})\xrightarrow{\sim}\operatorname{Hom}_{\mathrm{A}_0}(\mathrm{M}_0,\mathrm{N})
\]
de modo que o homomorfismo \sref{IV.8.5.2.2} nada mais é, a menos de isomorfismos canônicos, do que o homomorfismo canônico
\begin{equation}
  \varinjlim\operatorname{Hom}_{\mathrm{A}_0}(\mathrm{M}_0,\mathrm{N}_\lambda)\to\operatorname{Hom}_{\mathrm{A}_0}(\mathrm{M}_0,\varinjlim\mathrm{N}_\lambda),
  \tag{8.5.2.3}
\label{IV.8.5.2.3}
\end{equation}
\oldpage[IV-3]{21}
que associa a cada sistema indutivo de homomorfismos de $\mathrm{A}_0$-módulos $\theta_\lambda:\mathrm{M}_0\to\mathrm{N}_\lambda$ seu limite indutivo.

Ora, se $\mathrm{M}_0$ é de tipo finito (resp.\ de apresentação finita), existe uma sequência exata $\mathrm{A}_0^m\to\mathrm{M}_0\to 0$ (resp.\ $\mathrm{A}_0^n\to\mathrm{A}_0^m\to\mathrm{M}_0\to 0$); como é claro que \sref{IV.8.5.2.3} é bijetivo quando $\mathrm{M}_0$ é da forma $\mathrm{A}_0^q$, basta utilizar a exatidão à esquerda do funtor $\mathrm{M}_0\mapsto\operatorname{Hom}_{\mathrm{A}_0}(\mathrm{M}_0,\mathrm{P})$ e a exatidão do funtor $\varinjlim$ (na categoria dos grupos abelianos) para concluir.

\medskip
Passemos ao caso em que $\mathrm{S}_0$ é quase-compacto, e seja $(\mathrm{U}_i)$ uma cobertura finita de $\mathrm{S}_0$ por abertos afins; para cada $\lambda$, os $\mathrm{U}_{i\lambda}=u_{0\lambda}^{-1}(\mathrm{U}_i)$ formam uma cobertura de $\mathrm{S}_\lambda$ por abertos afins, e os $\mathrm{V}_i=u_0^{-1}(\mathrm{U}_i)$ uma cobertura de $\mathrm{S}$ por abertos afins.
Para demonstrar que $u_{\sh{F},\sh{G}}^*$ é injetivo, basta demonstrar que, se $f_\lambda:\sh{F}_\lambda\to\sh{G}_\lambda$ satisfaz $f=u_\lambda^*(f_\lambda)=0$, então existe $\mu\geqslant\lambda$ tal que $u_{\lambda\mu}^*(f_\lambda)=0$.
Em virtude do lema \sref{IV.8.5.2.1}, para cada $i$ existe um índice $\lambda_i$ tal que $f_\mu|_{\mathrm{U}_{i\mu}}=0$ para $\mu\geqslant\lambda_i$.
Portanto, basta tomar $\mu$ maior que todos os $\lambda_i$.

Suponhamos, além disso, que $\mathrm{S}_0$ seja quase-separado e que $\sh{F}_0$ seja de apresentação finita, e seja $f:\sh{F}\to\sh{G}$ um homomorfismo de $\sh{O}_\mathrm{S}$-Módulos.
Em virtude do lema \sref{IV.8.5.2.1}, para cada $i$ existem um índice $\lambda_i$ e um homomorfismo $f_{\lambda_i}^{(i)}:\sh{F}_{\lambda_i}|\mathrm{U}_{i,\lambda_i}\to\sh{G}_{\lambda_i}|\mathrm{U}_{i,\lambda_i}$ tais que $u_{\lambda_i}^*(f_{\lambda_i}^{(i)})=f|\mathrm{V}_i$.
Como $\mathrm{L}$ é filtrante, podemos supor, além disso, que todos os $\lambda_i$ sejam iguais a um único $\lambda$.
Ora, $\mathrm{S}_\lambda$ é quase-separado \sref{IV.1.2.3}, e $\sh{F}_\lambda$ é um $\sh{O}_{\mathrm{S}_\lambda}$-Módulo quase-coerente e de apresentação finita \sref[0]{0.5.2.5}; como, para cada par de índices $i,j$, $\mathrm{U}_{ij\lambda}=\mathrm{U}_{i\lambda}\cap\mathrm{U}_{j\lambda}$ é quase-compacto e, por definição,
\[
  u_\lambda^*(f_\lambda^{(i)}|\mathrm{U}_{ij\lambda})=u_\lambda^*(f_\lambda^{(j)}|\mathrm{U}_{ij\lambda})=f|(\mathrm{V}_i\cap\mathrm{V}_j),
\]
decorre do que foi demonstrado anteriormente que existe um índice $\lambda_{ij}$ tal que
\[
  u_{\lambda\mu}^*(f_\lambda^{(i)}|\mathrm{U}_{ij\lambda})=u_{\lambda\mu}^*(f_\lambda^{(j)}|\mathrm{U}_{ij\lambda})\quad\text{para }\mu\geqslant\lambda_{ij}.
\]
Tomando $\mu$ maior que todos os $\lambda_{ij}$, vê-se que $u_{\lambda\mu}^*(f_\lambda^{(i)})$ e $u_{\lambda\mu}^*(f_\lambda^{(j)})$ coincidem em $\mathrm{U}_{i\mu}\cap\mathrm{U}_{j\mu}$ para cada par $(i,j)$ e, portanto, definem um homomorfismo $f_\mu:\sh{F}_\mu\to\sh{G}_\mu$ tal que $f=u_\mu^*(f_\mu)$.

\end{enumerate}
\end{proof}

Antes de passar à demonstração de (ii), observemos os seguintes corolários de (i):

\begin{corollary}[8.5.2.4]
\label{IV.8.5.2.4}
Suponhamos que $\mathrm{S}_0$ seja quase-compacto, que $\sh{F}_\lambda$ seja quase-coerente e de tipo finito e que $\sh{G}_\lambda$ seja quase-coerente e de apresentação finita.
Seja $f_\lambda:\sh{F}_\lambda\to\sh{G}_\lambda$ um homomorfismo.
Para que $f=u_\lambda^*(f_\lambda):\sh{F}\to\sh{G}$ seja um isomorfismo, é necessário e suficiente que exista $\mu\geqslant\lambda$ tal que $f_\mu=u_{\lambda\mu}^*(f_\lambda):\sh{F}_\mu\to\sh{G}_\mu$ o seja.
\end{corollary}

\begin{proof}
Podemos sempre supor $\mathrm{S}_\lambda=\mathrm{S}_0$; sendo a questão local em $\mathrm{S}_0$, podemos, além disso (como $\mathrm{S}_0$ é quase-compacto e $\mathrm{L}$ é filtrante), reduzir-nos ao caso em que $\mathrm{S}_0$ é afim, portanto quase-separado.
Sendo a condição trivialmente suficiente, resta demonstrar que ela é necessária: por hipótese, existe um homomorfismo de $\sh{O}_\mathrm{S}$-Módulos $g:\sh{G}\to\sh{F}$ tal que $g\circ f=1_\sh{F}$ e $f\circ g=1_\sh{G}$.
Como $\sh{G}$ é de apresentação finita, existem um índice $\nu\geqslant\lambda$ e um homomorfismo $g_\nu:\sh{G}_\nu\to\sh{F}_\nu$ tais que $g=u_\nu^*(g_\nu)$ em virtude de \sref{IV.8.5.2},~(i); tem-se, portanto,
\[
  u_\nu^*(g_\nu\circ f_\nu)=1_\sh{F}
\]
\oldpage[IV-3]{22}
e $u_\nu^*(f_\nu\circ g_\nu)=1_\sh{G}$; como $\sh{F}_\nu$ e $\sh{G}_\nu$ são de tipo finito, decorre de \sref{IV.8.5.2},~(i), que existe $\mu\geqslant\nu$ tal que $g_\mu\circ f_\mu=1_{\sh{F}_\mu}$ e $f_\mu\circ g_\mu=1_{\sh{G}_\mu}$, de onde resulta o corolário.
\end{proof}

\begin{corollary}[8.5.2.5]
\label{IV.8.5.2.5}
Suponhamos que $\mathrm{S}_0$ seja quase-compacto e quase-separado.
Suponhamos que $\sh{F}_\lambda$ e $\sh{G}_\lambda$ sejam $\sh{O}_{\mathrm{S}_\lambda}$-Módulos quase-coerentes e de apresentação finita.
Para que $\sh{F}$ e $\sh{G}$ sejam isomorfos, é necessário e suficiente que exista $\mu\geqslant\lambda$ tal que $\sh{F}_\mu$ e $\sh{G}_\mu$ sejam isomorfos.
Além disso, para todo isomorfismo $f:\sh{F}\xrightarrow{\sim}\sh{G}$, existem $\nu\geqslant\mu$ e um isomorfismo $f_\nu:\sh{F}_\nu\xrightarrow{\sim}\sh{G}_\nu$ tais que $f=u_\nu^*(f_\nu)$.
\end{corollary}

\begin{proof}
Isto decorre de \sref{IV.8.5.2.4} e de \sref{IV.8.5.2},~(i), pois todo homomorfismo $f:\sh{F}\to\sh{G}$ é da forma $u_\mu^*(f_\mu)$ para algum $\mu\geqslant\lambda$ e algum homomorfismo $f_\mu:\sh{F}_\mu\to\sh{G}_\mu$.
\end{proof}

\begin{proof}[Demonstração de \sref{IV.8.5.2},~(ii)]
Consideremos novamente, em primeiro lugar, o caso em que $\mathrm{S}_0=\operatorname{Spec}(\mathrm{A}_0)$ é afim.
Então a afirmação equivale ao lema \sref{IV.5.13.7.1}.

No caso geral, partindo de uma cobertura finita por abertos afins $(\mathrm{U}_i)$ de $\mathrm{S}_0$, decorre de \sref{IV.5.13.7.1} que, para cada $i$, existem um índice $\lambda(i)$ e um $\sh{O}_{\mathrm{U}_{i,\lambda(i)}}$-Módulo quase-coerente e de apresentação finita $\sh{F}_i$ tais que $u_{\lambda(i)}^*(\sh{F}_i)=\sh{F}|_{\mathrm{V}_i}$ (com as notações de (i)).
Como $\mathrm{L}$ é filtrante, podemos supor, além disso, que todos os $\lambda_i$ sejam iguais a um único $\lambda$.
Como $\mathrm{U}_{ij\lambda}=\mathrm{U}_{i\lambda}\cap\mathrm{U}_{j\lambda}$ é quase-compacto e quase-separado \sref{IV.1.2.7}, decorre de \sref{IV.8.5.2.5} que, para cada par $(i,j)$, existem um índice $\lambda_{ij}\geqslant\lambda$ e um isomorfismo
\[
  \theta_{ij}:u_{\lambda,\lambda_{ij}}^*(\sh{F}_i|\mathrm{U}_{ij\lambda})\xrightarrow{\sim}u_{\lambda,\lambda_{ij}}^*(\sh{F}_j|\mathrm{U}_{ij\lambda})
\]
tal que $u_{\lambda_{ij}}^*(\theta_{ij})$ é o automorfismo identidade de $\sh{F}|(\mathrm{V}_i\cap\mathrm{V}_j)$; podemos supor, além disso, que todos os $\lambda_{ij}$ sejam iguais a $\lambda$.
Mudando as notações, podemos, portanto, supor que, para cada par $(i,j)$, exista um isomorfismo $\theta_{ij}:\sh{F}_i|_{\mathrm{U}_{ij\lambda}}\xrightarrow{\sim}\sh{F}_j|_{\mathrm{U}_{ij\lambda}}$ tal que $u_\lambda^*(\theta_{ij})$ seja o automorfismo identidade de $\sh{F}|_{(\mathrm{V}_i\cap\mathrm{V}_j)}$.
Finalmente, para quaisquer três índices $i$, $j$, $k$, se pusermos $\mathrm{U}_{ijk,\lambda}=\mathrm{U}_{i\lambda}\cap\mathrm{U}_{j\lambda}\cap\mathrm{U}_{k\lambda}$, $\mathrm{U}_{ijk,\lambda}$ é quase-compacto e, se $\theta'_{ij}$, $\theta'_{jk}$ e $\theta'_{ik}$ denotarem as restrições de $\theta_{ij}$, $\theta_{jk}$ e $\theta_{ik}$ a $\mathrm{U}_{ijk,\lambda}$, tem-se $u_\lambda^*(\theta'_{jk}\circ\theta'_{ij})=u_\lambda^*(\theta'_{ik})$.
Portanto, em virtude de (i), existe um índice $\mu\geqslant\lambda$ tal que $u_{\lambda\mu}^*(\theta'_{ik})=u_{\lambda\mu}^*(\theta'_{jk}\circ\theta'_{ij})$; por conseguinte, os isomorfismos $u_{\lambda\mu}^*(\theta_{ij}):u_{\lambda\mu}^*(\sh{F}_i)|_{\mathrm{U}_{ij\mu}}\xrightarrow{\sim}u_{\lambda\mu}^*(\sh{F}_j)|_{\mathrm{U}_{ij\mu}}$ satisfazem a condição de colagem e definem em $\mathrm{S}_\mu$ um $\sh{O}_{\mathrm{S}_\mu}$-Módulo quase-coerente e de apresentação finita $\sh{F}_\mu$ \sref[0]{0.3.3.1} tal que $\sh{F}$ é isomorfo a $u_\mu^*(\sh{F}_\mu)$.
\end{proof}

\begin{env}[8.5.3]
\label{IV.8.5.3}
O resultado de \sref{IV.8.5.2} também pode ser expresso dizendo que, se $\mathrm{S}_0$ é quase-compacto e quase-separado, a categoria dos $\sh{O}_\mathrm{S}$-Módulos quase-coerentes e de apresentação finita é determinada, a menos de equivalência, pelos dados das categorias dos $\sh{O}_{\mathrm{S}_\lambda}$-Módulos quase-coerentes e de apresentação finita, dos funtores $u_{\lambda\mu}^*$ entre essas categorias e dos isomorfismos de transição $u_{\mu\nu}^*\circ u_{\lambda\mu}^*\xrightarrow{\sim}u_{\lambda\nu}^*$.
De maneira figurada, pode-se dizer que o dado de um $\sh{O}_\mathrm{S}$-Módulo quase-coerente e de apresentação finita $\sh{F}$ é ``funtorialmente'' equivalente ao dado de um $\sh{O}_{\mathrm{S}_\lambda}$-Módulo quase-coerente e de apresentação finita $\sh{F}'_\lambda$ para algum $\lambda$ suficientemente grande e que, se outro $\sh{O}_{\mathrm{S}_\mu}$-Módulo quase-coerente e de apresentação finita $\sh{F}''_\mu$ também tiver $\sh{F}$ como imagem inversa, então $\sh{F}'_\lambda$ e $\sh{F}''_\mu$ terão a mesma imagem inversa em algum $\mathrm{S}_\nu$ adequado ($\nu\geqslant\lambda$, $\nu\geqslant\mu$).

Doravante, interpretaremos desse ponto de vista diversas noções relacionadas com os $\sh{O}_\mathrm{S}$-Módulos quase-coerentes.
\end{env}

\oldpage[IV-3]{23}
\begin{corollary}[8.5.4]
\label{IV.8.5.4}
Suponhamos que $\mathrm{S}_0$ seja quase-compacto e quase-separado; então, para todo $\sh{O}_{\mathrm{S}_\lambda}$-Módulo quase-coerente $\sh{G}_\lambda$, o homomorfismo canônico \sref{IV.8.5.1.4} é bijetivo.
\end{corollary}

\begin{proof}
Com efeito, basta aplicar \sref{IV.8.5.2},~(i), ao caso $\sh{F}_\lambda=\sh{O}_{\mathrm{S}_\lambda}$, que é de apresentação finita.
\end{proof}

\begin{proposition}[8.5.5]
\label{IV.8.5.5}
Suponhamos que $\mathrm{S}_0$ seja quase-compacto e que $\sh{F}_\lambda$ seja um $\sh{O}_{\mathrm{S}_\lambda}$-Módulo quase-coerente e de apresentação finita.
Para que $\sh{F}$ seja localmente livre (resp.\ localmente livre de posto $n$), é necessário e suficiente que exista $\mu\geqslant\lambda$ tal que $\sh{F}_\mu$ o seja.
\end{proposition}

\begin{proof}
A condição é trivialmente suficiente; demonstremos que é necessária.
Se $\sh{F}$ é localmente livre (resp.\ localmente livre de posto $n$), existe uma cobertura finita por abertos afins $(\mathrm{V}_i)$ de $\mathrm{S}$ tal que $\sh{F}|_{\mathrm{V}_i}$ seja isomorfo a $\sh{O}_\mathrm{S}^{n_i}|_{\mathrm{V}_i}$ (resp.\ $\sh{O}_\mathrm{S}^n|_{\mathrm{V}_i}$) para cada $i$.
Em virtude de \sref{IV.8.2.11}, existem um índice $\nu\geqslant\lambda$ e, para cada $i$, um aberto quase-compacto $\mathrm{U}_{i\nu}$ de $\mathrm{S}_\nu$ tal que $\mathrm{V}_i=u_\nu^{-1}(\mathrm{U}_{i\nu})$.
Como $\mathrm{S}_\nu$ é quase-compacto, cada $\mathrm{U}_{i\nu}$ é uma união finita de abertos afins; assim, reduzimo-nos ao caso em que $\mathrm{S}_0$ é afim e $\sh{F}=\sh{O}_\mathrm{S}^n$; sabemos então que existe $\mu\geqslant\lambda$ tal que $\sh{F}_\mu$ seja isomorfo a $\sh{O}_{\mathrm{S}_\mu}^n$ \sref{IV.8.5.2.5}.
\end{proof}

\begin{proposition}[8.5.6]
\label{IV.8.5.6}
Suponhamos que $\mathrm{S}_0$ seja quase-compacto e consideremos uma sequência
\[
  \sh{F}_\lambda\to\sh{G}_\lambda\to\sh{H}_\lambda\to 0
\]
de homomorfismos de $\sh{O}_{\mathrm{S}_\lambda}$-Módulos quase-coerentes, em que $\sh{F}_\lambda$ e $\sh{G}_\lambda$ são de tipo finito e $\sh{H}_\lambda$ é de apresentação finita.
Para que a sequência correspondente $\sh{F}\to\sh{G}\to\sh{H}\to 0$ seja exata, é necessário e suficiente que exista $\mu\geqslant\lambda$ tal que a sequência $\sh{F}_\mu\to\sh{G}_\mu\to\sh{H}_\mu\to 0$ o seja (caso em que o mesmo vale para $\nu\geqslant\mu$).
\end{proposition}

\begin{proof}
A suficiência da condição e a última afirmação decorrem do fato de o funtor $u_\lambda^*$ (resp.\ $u_{\mu\nu}^*$) ser exato à direita.
Para demonstrar que a condição é necessária, observemos que, da hipótese e de \sref{IV.8.5.2},~(i), decorre que existe $\nu\geqslant\lambda$ tal que a composição $\sh{F}_\nu\to\sh{G}_\nu\to\sh{H}_\nu$ é nula.
Se pusermos $\sh{H}'_\nu=\operatorname{Coker}(\sh{F}_\nu\to\sh{G}_\nu)$, temos, portanto, um homomorfismo $f_\nu:\sh{H}'_\nu\to\sh{H}_\nu$; por hipótese, $u_\nu^*(f_\nu)$ é um isomorfismo, e decorre de \sref{IV.8.5.2.4} que existe $\mu\geqslant\nu$ tal que $u_{\nu\mu}^*(f_\nu)$ é um isomorfismo, o que conclui a demonstração.
\end{proof}

\begin{corollary}[8.5.7]
\label{IV.8.5.7}
Suponhamos que $\mathrm{S}_0$ seja quase-compacto, que $\sh{F}_\lambda$ seja quase-coerente, que $\sh{G}_\lambda$ seja quase-coerente e de tipo finito, e seja $f_\lambda:\sh{F}_\lambda\to\sh{G}_\lambda$ um homomorfismo.
Para que $f=u_\lambda^*(f_\lambda):\sh{F}\to\sh{G}$ seja sobrejetivo, é necessário e suficiente que exista $\mu\geqslant\lambda$ tal que $f_\mu=u_{\lambda\mu}^*(f_\lambda):\sh{F}_\mu\to\sh{G}_\mu$ o seja.
\end{corollary}

\begin{proof}
Este é o caso particular de \sref{IV.8.5.6} aplicado à sequência $0\to\sh{H}_\lambda\to 0\to 0$, em que $\sh{H}_\lambda=\operatorname{Coker}(f_\lambda)$, que é quase-coerente e de tipo finito (levando em conta que $\sh{H}=\operatorname{Coker}(f)$ e $\sh{H}_\mu=\operatorname{Coker}(f_\mu)$, em virtude da exatidão à direita dos funtores $u_\lambda^*$ e $u_{\lambda\mu}^*$).
\end{proof}

\begin{corollary}[8.5.8]
\label{IV.8.5.8}
Suponhamos que $\mathrm{S}_0$ seja quase-compacto e que os morfismos $u_{\lambda\mu}:\mathrm{S}_\mu\to\mathrm{S}_\lambda$ sejam planos.
Então:
\begin{enumerate}[label=(\roman*)]
  \item Seja $\sh{F}_\lambda\xrightarrow{f_\lambda}\sh{G}_\lambda\xrightarrow{g_\lambda}\sh{H}_\lambda$ uma sequência de homomorfismos de $\sh{O}_{\mathrm{S}_\lambda}$-Módulos quase-coerentes, tais que $\operatorname{Im} f_\lambda$ e $\operatorname{Ker} g_\lambda$ sejam de tipo finito.
Para que a sequência correspondente $\sh{F}\xrightarrow{f}\sh{G}\xrightarrow{g}\sh{H}$ seja
\oldpage[IV-3]{24}
exata, é necessário e suficiente que exista $\mu\geqslant\lambda$ tal que a sequência $\sh{F}_\mu\xrightarrow{f_\mu}\sh{G}_\mu\xrightarrow{g_\mu}\sh{H}_\mu$ o seja (caso em que o mesmo vale para $\nu\geqslant\mu$).
  \item Seja $f_\lambda:\sh{F}_\lambda\to\sh{G}_\lambda$ um homomorfismo de $\sh{O}_{\mathrm{S}_\lambda}$-Módulos quase-coerentes, e suponhamos que $\operatorname{Ker} f_\lambda$ seja de tipo finito.
Para que $f=u_\lambda^*(f_\lambda):\sh{F}\to\sh{G}$ seja injetivo, é necessário e suficiente que exista $\mu\geqslant\lambda$ tal que $f_\mu=u_{\lambda\mu}^*(f_\lambda):\sh{F}_\mu\to\sh{G}_\mu$ o seja.
\end{enumerate}
\end{corollary}

\begin{proof}
(i) Levando em conta \sref{IV.8.3.8},~(ii), observemos que, por planicidade, $\operatorname{Im} f$ e $\operatorname{Ker} g$ (resp. $\operatorname{Im} f_\mu$ e $\operatorname{Ker} g_\mu$ para $\mu\geqslant\lambda$) são as imagens inversas de $\operatorname{Im} f_\lambda$ e $\operatorname{Ker} g_\lambda$ \sref[0\textsubscript{I}]{0.6.7.2}. Suponhamos que a sequência $\sh{F}\xrightarrow{f}\sh{G}\xrightarrow{g}\sh{H}$ seja exata. Como $\operatorname{Im} f_\lambda$ é de tipo finito, decorre de \sref{IV.8.5.2},~(i), que existe $\mu\geqslant\lambda$ tal que a composição $\sh{F}_\mu\xrightarrow{f_\mu}\sh{G}_\mu\xrightarrow{g_\mu}\sh{H}_\mu$ é nula. Mudando as notações, podemos, portanto, supor que $g_\lambda\circ f_\lambda=0$. Como $\operatorname{Ker} g_\lambda$ é de tipo finito, decorre então de \sref{IV.8.5.7} que existe $\mu\geqslant\lambda$ tal que o homomorfismo $\sh{F}_\mu\to\operatorname{Ker} g_\mu$ seja sobrejetivo, o que demonstra (i).

(ii) Este é o caso particular de (i) aplicado à sequência $0\to\sh{F}_\lambda\to\sh{G}_\lambda$.
\end{proof}

\begin{lemma}[8.5.9]
\label{IV.8.5.9}
Suponhamos que $\mathrm{S}_0$ seja quase-compacto e que $\sh{F}_\lambda$ seja quase-coerente e de tipo finito. Sejam $\sh{G}'_\lambda$ e $\sh{G}''_\lambda$ dois quocientes quase-coerentes de $\sh{F}_\lambda$, supondo-se, além disso, que $\sh{G}'_\lambda$ seja de apresentação finita. Para que $\sh{G}''$ seja um quociente de $\sh{G}'$, é necessário e suficiente que exista $\mu\geqslant\lambda$ tal que $\sh{G}''_\mu$ seja um quociente de $\sh{G}'_\mu$.
\end{lemma}

\begin{proof}
Por hipótese, existem dois homomorfismos sobrejetivos $p'_\lambda:\sh{F}_\lambda\to\sh{G}'_\lambda$ e $p''_\lambda:\sh{F}_\lambda\to\sh{G}''_\lambda$. Pela exatidão à direita de $u_\lambda^*$ e $u_{\lambda\mu}^*$, os homomorfismos $p'=u_\lambda^*(p'_\lambda)$, $p''=u_\lambda^*(p''_\lambda)$, $p'_\mu=u_{\lambda\mu}^*(p'_\lambda)$ e $p''_\mu=u_{\lambda\mu}^*(p''_\lambda)$ também são sobrejetivos. Além disso, se existir um homomorfismo $f:\sh{G}'\to\sh{G}''$ (resp. $f_\mu:\sh{G}'_\mu\to\sh{G}''_\mu$) tal que $p''=f\circ p'$ (resp. $p''_\mu=f_\mu\circ p'_\mu$), tal homomorfismo será necessariamente único; isso mostra que a questão é local em $\mathrm{S}_\lambda$ e, por conseguinte, como $\mathrm{S}_0$ é quase-compacto e $\mathbf{L}$ é filtrante, podemos supor que $\mathrm{S}_\lambda$ seja afim e, portanto, quase-separado. A condição do enunciado é claramente suficiente. Reciprocamente, como $\sh{G}'_\lambda$ é de apresentação finita e $\mathrm{S}_\lambda$ é quase-compacto e quase-separado, decorre de \sref{IV.8.5.2},~(i), que, se existir $f:\sh{G}'\to\sh{G}''$ tal que $p''=f\circ p'$, então existirão $\mu\geqslant\lambda$ e $f_\mu:\sh{G}'_\mu\to\sh{G}''_\mu$ tais que $f=u_\mu^*(f_\mu)$ e $p''_\mu=f_\mu\circ p'_\mu$, o que demonstra o lema.
\end{proof}

\begin{env}[8.5.10]
\label{IV.8.5.10}
No restante desta seção, para todo módulo quase-coerente $\sh{F}$ sobre um pré-esquema, denotemos por $\mathfrak{Q}(\sh{F})$ o conjunto dos módulos quocientes de $\sh{F}$ que são de apresentação finita. Se $\sh{F}_\lambda$ é quase-coerente e $\sh{G}_\lambda\in\mathfrak{Q}(\sh{F}_\lambda)$, segue da exatidão à direita de $u_{\lambda\mu}^*$ e $u_\lambda^*$ que $\sh{G}_\mu\in\mathfrak{Q}(\sh{F}_\mu)$ para $\mu\geqslant\lambda$ e $\sh{G}\in\mathfrak{Q}(\sh{F})$. É claro que $(\mathfrak{Q}(\sh{F}_\lambda),u_{\lambda\mu}^*)$ é um sistema indutivo de conjuntos e que as aplicações $u_\lambda^*:\mathfrak{Q}(\sh{F}_\lambda)\to\mathfrak{Q}(\sh{F})$ formam um sistema indutivo de aplicações; portanto, passando ao limite indutivo, obtém-se uma aplicação canônica
\[
u_{\mathfrak{Q}}:\varinjlim\mathfrak{Q}(\sh{F}_\lambda)\longrightarrow\mathfrak{Q}(\sh{F}).
\tag{8.5.10.1}\label{IV.8.5.10.1}
\]
Além disso, se $(\sh{F}'_\lambda)$ é uma segunda família de $\sh{O}_{\mathrm{S}_\lambda}$-módulos quase-coerentes e, para todo $\lambda$, $\sh{F}'_\lambda$ é um quociente de $\sh{F}_\lambda$, então $\sh{F}'$ é um quociente de $\sh{F}$, e temos um diagrama comutativo

\oldpage[IV-3]{25}

\[
\begin{tikzcd}
\varinjlim\mathfrak{Q}(\sh{F}_\lambda) \ar[r,"u_{\mathfrak{Q}}"] \ar[d] & \mathfrak{Q}(\sh{F}) \ar[d] \\
\varinjlim\mathfrak{Q}(\sh{F}'_\lambda) \ar[r,"u_{\mathfrak{Q}}"'] & \mathfrak{Q}(\sh{F}')
\end{tikzcd}
\tag{8.5.10.2}\label{IV.8.5.10.2}
\]
\end{env}

\begin{proposition}[8.5.11]
\label{IV.8.5.11}
Suponhamos que $\mathrm{S}_0$ seja quase-compacto (resp. quase-compacto e quase-separado). Suponhamos que $\sh{F}_\lambda$ seja quase-coerente e de tipo finito (resp. de apresentação finita) para todo $\lambda$; então a aplicação canônica \eqref{IV.8.5.10.1} é injetiva (resp. bijetiva).
\end{proposition}

\begin{proof}
A primeira afirmação segue, mais precisamente, do Lema \sref{IV.8.5.9}. Para a segunda, considere um $\sh{O}_{\mathrm{S}}$-Módulo quociente $\sh{G}$ de $\sh{F}$ que seja de apresentação finita. Segue de \sref{IV.8.5.2},~(ii) que existem $\lambda'\in\mathbf{L}$ e um $\sh{O}_{\mathrm{S}_{\lambda'}}$-Módulo quase-coerente $\sh{G}_{\lambda'}$ de apresentação finita tais que $\sh{G}=u_{\lambda'}^*(\sh{G}_{\lambda'})$; como $\mathbf{L}$ é filtrado, podemos substituir $\lambda$ e $\lambda'$ por um índice que majore ambos e supor que $\lambda'=\lambda$. Considere então o homomorfismo canônico $p:\sh{F}\to\sh{G}$. Segue de \sref{IV.8.5.2},~(i) que existem $\mu\geqslant\lambda$ e um homomorfismo $p_\mu:\sh{F}_\mu\to\sh{G}_\mu$ tais que $p=u_\mu^*(p_\mu)$. Por \sref{IV.8.5.7}, podemos, além disso, supor que $\mu$ é suficientemente grande para que $p_\mu$ seja sobrejetivo, o que completa a demonstração.
\end{proof}

\subsection{Subpré-esquemas de apresentação finita de um limite projetivo de pré-esquemas}
\label{subsection:IV.8.6}

\begin{env}[8.6.1]
\label{IV.8.6.1}
Dado um pré-esquema $Y$, denotamos nesta seção por $\mathfrak{Spr}(Y)$ o conjunto ordenado
\sref[I]{I.4.1.10} dos subpré-esquemas de $Y$ que são \emph{de apresentação finita}
\sref{IV.1.6.1}, e por $\mathfrak{Spro}(Y)$ (resp.\ $\mathfrak{Sprf}(Y)$) o subconjunto de $\mathfrak{Spr}(Y)$ formado pelos subpré-esquemas induzidos em abertos (resp.\ em subpré-esquemas fechados) de $Y$, de apresentação finita sobre $Y$; equivale a dizer que um subpré-esquema $Z$ de $Y$ pertence a $\mathfrak{Spro}(Y)$ (resp.\ $\mathfrak{Sprf}(Y)$) ou que é induzido em um aberto e que o espaço subjacente é \emph{retrocompacto} em $Y$ (resp.\ que é fechado e que o ideal $\sh{J}$ de $\sh{O}_Y$ que o define é \emph{de tipo finito}, o que também significa que $j_*(\sh{O}_Z)\in\mathfrak{Q}(\sh{O}_Y)$, onde $j:Z\to Y$ é a injeção canônica) \sref{IV.1.6.1} e \sref{IV.1.4.5}.
\end{env}

\begin{env}[8.6.2]
\label{IV.8.6.2}
Conservamos sempre as notações de \sref{IV.8.2.2} e supomos que $S_0$ seja um dos $S_\lambda$.
Seja $Y_\lambda$ um subpré-esquema de $S_\lambda$; então $Y_\mu=u_{\lambda\mu}^{-1}(Y_\lambda)$ (resp.\ $Y=u_\lambda^{-1}(Y_\lambda)$) é um subpré-esquema de $S_\mu$ para $\mu\geqslant\lambda$ (resp.\ de $S$). Ele é induzido em um aberto (resp.\ é fechado) se $Y_\lambda$ o for \sref[I]{I.4.3.2}, e é de apresentação finita sobre $S_\mu$ (resp.\ $S$) se $Y_\lambda$ for de apresentação finita sobre $S_\lambda$ \sref{IV.1.6.2},~(iii). Segue-se que $(\mathfrak{Spr}(S_\lambda),u_{\lambda\mu}^{-1})$ (resp.\ $(\mathfrak{Spro}(S_\lambda),u_{\lambda\mu}^{-1})$, $(\mathfrak{Sprf}(S_\lambda),u_{\lambda\mu}^{-1})$) é um \emph{sistema indutivo} de conjuntos e que as aplicações $u_\lambda^{-1}:\mathfrak{Spr}(S_\lambda)\to\mathfrak{Spr}(S)$ (resp.\ $\mathfrak{Spro}(S_\lambda)\to\mathfrak{Spro}(S)$, $\mathfrak{Sprf}(S_\lambda)\to\mathfrak{Sprf}(S)$) formam um sistema indutivo de aplicações; de onde, passando ao limite indutivo, obtêm-se aplicações canônicas
\end{env}

\begin{align}
  \varinjlim \mathfrak{Spr}(S_\lambda) &\to \mathfrak{Spr}(S) \tag{8.6.2.1}\label{IV.8.6.2.1}\\
  \varinjlim \mathfrak{Spro}(S_\lambda) &\to \mathfrak{Spro}(S) \tag{8.6.2.2}\label{IV.8.6.2.2}\\
  \varinjlim \mathfrak{Sprf}(S_\lambda) &\to \mathfrak{Sprf}(S). \tag{8.6.2.3}\label{IV.8.6.2.3}
\end{align}

\oldpage[IV-3]{26}

Recordemos, por \sref[I]{I.4.1.10}, que $\mathfrak{Spr}(X)$, para todo pré-esquema $X$, é um conjunto ordenado pela relação «$Z$ é um subpré-esquema do subpré-esquema $Y$», que escrevemos $Z\leqslant Y$.
As aplicações $u_{\lambda\mu}^{-1}$ e $u_\lambda^{-1}$ são \emph{crescentes} relativamente às relações de ordem correspondentes em $\mathfrak{Spr}(S_\lambda)$, $\mathfrak{Spr}(S_\mu)$ e $\mathfrak{Spr}(S)$.
Além disso, escrevendo $\eta\leqslant\eta'$ para dois elementos desse conjunto sempre que existam $\lambda$ e dois elementos $Y_\lambda$, $Y'_\lambda$ de $\mathfrak{Spr}(S_\lambda)$, cujas imagens canônicas sejam $\eta$ e $\eta'$, e tais que $Y_\lambda\leqslant Y'_\lambda$, verifica-se facilmente que isso não depende da escolha dos representantes e que se obtém, com efeito, uma \emph{relação de ordem}.
O fato de as $u_\lambda^{-1}$ serem crescentes implica imediatamente que a aplicação canônica \eqref{IV.8.6.2.1} é \emph{crescente}; o mesmo vale evidentemente para \eqref{IV.8.6.2.2} e \eqref{IV.8.6.2.3}.

\begin{proposition}[8.6.3]
\label{IV.8.6.3}
Suponhamos que $S_0$ seja quase-compacto (resp.\ quase-compacto e quase-separado).
Então as aplicações \eqref{IV.8.6.2.1}, \eqref{IV.8.6.2.2} e \eqref{IV.8.6.2.3} são injetivas (resp.\ bijetivas).
\end{proposition}

\begin{proof}
Levando em conta as observações de \sref{IV.8.6.1}, as afirmações relativas a \eqref{IV.8.6.2.3} decorrem de \sref{IV.8.5.11} aplicado a $\sh{F}_\lambda=\sh{O}_{S_\lambda}$; do mesmo modo, as afirmações relativas a \eqref{IV.8.6.2.2} são casos particulares de \sref{IV.8.3.5} e \sref{IV.8.3.11}, levando em conta que os $S_\lambda$ e $S$ são quase-compactos.
Resta considerar a aplicação \eqref{IV.8.6.2.1}.
Demonstremos primeiro que ela é sobrejetiva quando $S_0$ é quase-compacto e quase-separado.
Seja $Z$ um subpré-esquema de $S$, de apresentação finita sobre $S$; como $S$ é quase-compacto, $Z$ também o é, e existe, portanto, um aberto quase-compacto $U$ de $S$ tal que $Z$ seja um subpré-esquema de $U$, fechado em $U$, de apresentação finita sobre $U$.
Existem então um índice $\lambda$ e um aberto quase-compacto $U_\lambda$ de $S_\lambda$ tais que $U=u_\lambda^{-1}(U_\lambda)$ \sref{IV.8.2.11}; como $S_\lambda$ é quase-separado, $U_\lambda$ também o é \sref{IV.1.2.7}, e podemos, portanto, restringir-nos ao caso $U=S$; mas, nesse caso, reduzimo-nos ao fato de que \eqref{IV.8.6.2.3} é sobrejetiva.

Por fim, para ver que \eqref{IV.8.6.2.1} é injetiva quando $S_0$ é quase-compacto, bastará demonstrar o resultado mais preciso a seguir:
\end{proof}

\begin{corollary}[8.6.3.1]
\label{IV.8.6.3.1}
Suponhamos que $S_0$ seja quase-compacto e sejam $Z'_\lambda$, $Z''_\lambda$ dois subpré-esquemas de $S_\lambda$, ambos de apresentação finita.
Para que $Z'=u_\lambda^{-1}(Z'_\lambda)$ seja majorado por $Z''=u_\lambda^{-1}(Z''_\lambda)$ \sref[I]{I.4.1.10}, é necessário e suficiente que exista $\mu\geqslant\lambda$ tal que $Z'_\mu=u_{\lambda\mu}^{-1}(Z'_\lambda)$ seja majorado por $Z''_\mu=u_{\lambda\mu}^{-1}(Z''_\lambda)$.
\end{corollary}

\begin{proof}
É trivial que a condição é suficiente.
Para ver que ela é necessária, observemos primeiro que os conjuntos subjacentes $Z'_\lambda$ e $Z''_\lambda$ são localmente construtíveis em $S_\lambda$ pela hipótese \sref{IV.1.8.4}; portanto, a hipótese implica que existe um $\nu\geqslant\lambda$ tal que $Z'_\nu\subset Z''_\nu$ \sref{IV.8.3.5}; substituindo $\lambda$ por $\nu$, podemos, portanto, supor desde já que $Z'_\lambda\subset Z''_\lambda$.
Além disso, pela hipótese \sref{IV.1.6.1}, os subespaços $Z'_\lambda$ e $Z''_\lambda$ de $S_\lambda$ são quase-compactos; para cada ponto $x\in Z'_\lambda$, existe uma vizinhança aberta quase-compacta $V(x)$ em $S_\lambda$ tal que $V(x)\cap Z'_\lambda$ e $V(x)\cap Z''_\lambda$ são fechados em $V(x)$.
Cobrindo $Z'_\lambda$ por um número finito de tais vizinhanças, vemos que existe um aberto quase-compacto $U_\lambda$ de $S_\lambda$ que contém $Z'_\lambda$ e tal que $Z'_\lambda$ e $Z''_\lambda\cap U_\lambda$ são fechados em $U_\lambda$.
Se denotarmos por $Y_\lambda$ o subpré-esquema de $S_\lambda$ induzido por $Z''_\lambda$ em $U_\lambda\cap Z''_\lambda$, é claro que, com as notações usuais, $Y_\mu$ (resp.\ $Y$) é induzido por $Z''_\mu$ em $U_\mu\cap Z''_\mu$ para $\mu\geqslant\lambda$ (resp.\ por $Z''$ em $U\cap Z''$); além disso, $Z'$ é majorado por $Y$ \sref[I]{I.4.4.1}, e, como basta demonstrar que $Z'_\mu$ é majorado por $Y_\mu$ para $\mu$
\oldpage[IV-3]{27}
 suficientemente grande, vemos finalmente que (substituindo $S_\lambda$ por $U_\lambda$) nos reduzimos ao caso em que $Z'_\lambda$ e $Z''_\lambda$ são subpré-esquemas fechados de $S_\lambda$.
Mas isso já foi demonstrado, pois \eqref{IV.8.6.2.3} é crescente e injetiva.
\end{proof}

\begin{corollary}[8.6.4]
\label{IV.8.6.4}
Suponhamos que $S_0$ seja quase-compacto e seja $Z_\lambda$ um subpré-esquema de $S_\lambda$, de apresentação finita sobre $S_\lambda$.
Para que $Z=u_\lambda^{-1}(Z_\lambda)$ seja o subpré-esquema induzido em um aberto de $S$ (resp.\ um subpré-esquema fechado de $S$), é necessário e suficiente que exista $\mu\geqslant\lambda$ tal que $Z_\mu=u_{\lambda\mu}^{-1}(Z_\lambda)$ seja o subpré-esquema induzido em um aberto de $S_\mu$ (resp.\ um subpré-esquema fechado de $S_\mu$).
\end{corollary}

\begin{proof}
Seja $(U_\lambda^{(i)})$ uma cobertura finita por abertos afins de $S_\lambda$, e ponhamos $U^{(i)}_\mu=u_{\lambda\mu}^{-1}(U_\lambda^{(i)})$ para $\mu\geqslant\lambda$ e $U^{(i)}=u_\lambda^{-1}(U_\lambda^{(i)})$.
Se $Z$ é aberto (resp.\ fechado) em $S$, $Z\cap U^{(i)}$ é aberto (resp.\ fechado) em $U^{(i)}$; reciprocamente, se cada $Z_\mu\cap U^{(i)}_\mu$ é aberto (resp.\ fechado) em $U^{(i)}_\mu$, então $Z_\mu$ é aberto (resp.\ fechado) em $S_\mu$.
Como $L$ é filtrado, basta demonstrar o corolário quando $S_\lambda$ é afim e, portanto, quase-separado; o resultado decorre do fato de que as aplicações \eqref{IV.8.6.2.1}, \eqref{IV.8.6.2.2} e \eqref{IV.8.6.2.3} são bijetivas.
\end{proof}

\subsection{Critérios para que um limite projetivo de pré-esquemas seja um pré-esquema reduzido (resp.\ integral)}
\label{subsection:IV.8.7}

Conservamos as hipóteses e notações de \sref{IV.8.2.2} e supomos sempre que $S_0$ seja um dos $S_\lambda$.

\begin{proposition}[8.7.1]
\label{IV.8.7.1}
Suponhamos que $S$ não seja reduzido.
Então existe $\lambda_0$ tal que, para $\lambda\geqslant\lambda_0$, $S_\lambda$ não é reduzido.
\end{proposition}

\begin{proof}
Como a questão é local em $S_0$, podemos supor que $S_0=\operatorname{Spec}(A_0)$ seja afim, de onde $S=\operatorname{Spec}(A)$, onde $A=\varinjlim A_\lambda$ é o limite indutivo de um sistema indutivo de $A_0$-álgebras $(A_\lambda)$.
Sabemos então, por \sref{IV.5.13.2}, que o nilradical de $A$ é o limite indutivo dos nilradicais dos $A_\lambda$; se ele não é nulo, um dos $A_\lambda$ contém, portanto, um elemento nilpotente $a_\lambda\neq 0$ cuja imagem em $A$ é nilpotente e não nula, e, consequentemente, a imagem de $a_\lambda$ em $A_\mu$ é nilpotente e não nula para $\mu\geqslant\lambda$.
\end{proof}

\begin{proposition}[8.7.2]
\label{IV.8.7.2}
Suponhamos que uma das seguintes hipóteses seja satisfeita:
\begin{enumerate}[label=\emph{\alph*)}]
\item $S_0$ é quase-compacto, o nilradical $\sh{N}_0$ de $\sh{O}_{S_0}$ é um ideal de tipo finito (o que ocorre, por exemplo, quando $S_0$ é noetheriano), e os morfismos $u_{\lambda\mu}:S_\mu\to S_\lambda$ são imersões abertas.
\item Os morfismos $u_{\lambda\mu}:S_\mu\to S_\lambda$ são fielmente planos.
\end{enumerate}
Nessas condições, para que $S$ seja reduzido, é necessário e suficiente que exista $\lambda_0$ tal que $S_\lambda$ seja reduzido para $\lambda\geqslant\lambda_0$.

Além disso, no caso b), se $S$ é reduzido, todos os $S_\lambda$ são reduzidos.
\end{proposition}

\begin{proof}
A última afirmação decorre do fato de que o morfismo $S\to S_\lambda$ é então fielmente plano para todo $\lambda$ \sref{IV.8.3.8}, e de \sref{IV.2.1.13}.
Por outro lado, \sref{IV.8.7.1} demonstra que a condição do enunciado é suficiente (sem hipótese sobre $S_0$ nem sobre os $u_{\lambda\mu}$).
Resta, portanto, ver que a condição é necessária sob a hipótese $a)$; então, por \sref{IV.8.2.13}, o espaço subjacente de $S$ identifica-se com a interseção dos espaços subjacentes dos $S_\lambda$ (identificando-se os $S_\lambda$ com subpré-esquemas induzidos em abertos de $S_0$), e o feixe estrutural $\sh{O}_S$ identifica-se com
\oldpage[IV-3]{28}
o feixe induzido em $S$ por cada um dos $\sh{O}_{S_\lambda}$; em particular, para todo $s\in S$, o anel local $\sh{O}_s$ é o mesmo para $S$ e para todos os $S_\lambda$.
Se $\sh{N}_\lambda$ é o nilradical de $\sh{O}_{S_\lambda}$, o nilradical $\sh{N}$ de $\sh{O}_S$ tem, portanto, em cada ponto de $S$, o mesmo talo $\sh{N}_s$ (o nilradical de $\sh{O}_s$) que $u_\lambda^*(\sh{N}_\lambda)$ (induzido em $S$ por $\sh{N}_\lambda$).
A hipótese de que $S$ é reduzido implica, portanto, $u_\lambda^*(\sh{N}_\lambda)=0$; como $\sh{N}_0$ é suposto de tipo finito, o mesmo vale para os $\sh{N}_\lambda$, e a conclusão decorre de \sref{IV.8.5.7}.
\end{proof}

\begin{corollary}[8.7.3]
\label{IV.8.7.3}
Suponhamos que uma das seguintes hipóteses seja satisfeita:
\begin{enumerate}[label=\emph{\alph*)}]
\item $S_0$ é um pré-esquema noetheriano e os morfismos $u_{\lambda\mu}:S_\mu\to S_\lambda$ são imersões abertas.
\item Os morfismos $u_{\lambda\mu}:S_\mu\to S_\lambda$ são fielmente planos.
\end{enumerate}
Então, para que $S$ seja integral, é necessário e suficiente que exista $\lambda_0$ tal que $S_\lambda$ seja integral para $\lambda\geqslant\lambda_0$.
\end{corollary}

\begin{proof}
Dizer que um pré-esquema é integral significa que ele é simultaneamente reduzido e irredutível. O caráter reduzido decorre de \sref{IV.8.7.2}. No caso (a), a irredutibilidade decorre de \sref{IV.8.4.3}. No caso (b), \sref{IV.8.3.8} mostra que todo morfismo $S\to S_\lambda$ é fielmente plano, logo sobrejetivo; assim, a irredutibilidade de $S$ implica a irredutibilidade de todo $S_\lambda$. Reciprocamente, a irredutibilidade dos $S_\lambda$ a partir de certo índice implica a irredutibilidade de $S$ pelo argumento geral de limite projetivo utilizado na demonstração de \sref{IV.8.4.1}.
\end{proof}

\begin{remark}[8.7.4]
\label{IV.8.7.4}
Se não se faz nenhuma hipótese sobre os $u_{\lambda\mu}$, pode ocorrer que $S$ seja integral embora todos os $S_\lambda$ sejam não reduzidos e não conexos, como mostra o exemplo \sref{IV.8.4.6},~(ii), tomando o anel $A$ não reduzido.
\end{remark}

\subsection{Pré-esquemas de apresentação finita sobre um limite projetivo de pré-esquemas}
\label{subsection:IV.8.8}

\begin{env}[8.8.1]
\label{IV.8.8.1}
Conservando sempre as notações e hipóteses de \sref{IV.8.2.2}, suporemos dados nesta seção dois $S_\alpha$-pré-esquemas $X_\alpha$, $Y_\alpha$, o que define \sref{IV.8.2.5} dois sistemas projetivos de pré-esquemas $(X_\lambda,v_{\lambda\mu})$ e $(Y_\lambda,w_{\lambda\mu})$, pondo $X_\lambda=X_\alpha\times_{S_\alpha}S_\lambda$, $Y_\lambda=Y_\alpha\times_{S_\alpha}S_\lambda$, $v_{\lambda\mu}=1_{X_\alpha}\times u_{\lambda\mu}$, $w_{\lambda\mu}=1_{Y_\alpha}\times u_{\lambda\mu}$ (para $\alpha\leqslant\lambda\leqslant\mu$), cujos limites projetivos são respectivamente $X=X_\alpha\times_{S_\alpha}S$, $Y=Y_\alpha\times_{S_\alpha}S$, sendo os morfismos canônicos $v_\lambda:X\to X_\lambda$ e $w_\lambda:Y\to Y_\lambda$ respectivamente iguais a $1_{X_\alpha}\times u_\lambda$ e $1_{Y_\alpha}\times u_\lambda$.
Para $\alpha\leqslant\lambda\leqslant\mu$, existe uma aplicação canônica $e_{\mu\lambda}:\operatorname{Hom}_{S_\lambda}(X_\lambda,Y_\lambda)\to\operatorname{Hom}_{S_\mu}(X_\mu,Y_\mu)$, que a todo $S_\lambda$-morfismo $f_\lambda:X_\lambda\to Y_\lambda$ associa $f_\mu=f_\lambda\times 1_{S_\mu}:X_\lambda\times_{S_\lambda}S_\mu\to Y_\lambda\times_{S_\lambda}S_\mu$, e é claro que $(\operatorname{Hom}_{S_\lambda}(X_\lambda,Y_\lambda),e_{\mu\lambda})$ é um \emph{sistema indutivo} de conjuntos.
Analogamente, existe uma aplicação canônica $e_\lambda:\operatorname{Hom}_{S_\lambda}(X_\lambda,Y_\lambda)\to\operatorname{Hom}_S(X,Y)$ que a $f_\lambda$ associa $f=f_\lambda\times 1_S:X_\lambda\times_{S_\lambda}S\to Y_\lambda\times_{S_\lambda}S$, e $(e_\lambda)$ é um sistema indutivo de aplicações; daí se obtém, passando ao limite indutivo, uma aplicação canônica, funtorial em $S_\alpha$, $X_\alpha$ e $Y_\alpha$:
\[
  e:\varinjlim\operatorname{Hom}_{S_\lambda}(X_\lambda,Y_\lambda)\to\operatorname{Hom}_S(X,Y).
  \tag{8.8.1.1}\label{IV.8.8.1.1}
\]
\end{env}

\begin{theorem}[8.8.2]
\label{IV.8.8.2}
\begin{enumerate}[label=(\roman*)]
\item Suponhamos que $X_\alpha$ seja quase-compacto (resp.\ quase-compacto e quase-separado), e que $Y_\alpha$ seja localmente de tipo finito (resp.\ localmente de apresentação finita) sobre $S_\alpha$.
Então a aplicação \eqref{IV.8.8.1.1} é injetiva (resp.\ bijetiva).
\item Suponhamos que $S_0$ seja quase-compacto e quase-separado.
Para todo pré-esquema $X$ de apresentação finita sobre $S$, existem um $\lambda\in L$, um pré-esquema $X_\lambda$ de apresentação finita sobre $S_\lambda$ e um $S$-isomorfismo $X\xrightarrow{\sim}X_\lambda\times_{S_\lambda}S$.
\end{enumerate}
\end{theorem}

\begin{proof}
\oldpage[IV-3]{29}
(i) Consideremos primeiro o caso em que $S_0=\operatorname{Spec}(A_0)$, $X_\alpha=\operatorname{Spec}(B_\alpha)$ e $Y_\alpha=\operatorname{Spec}(C_\alpha)$ são afins; então os $S_\lambda=\operatorname{Spec}(A_\lambda)$ e $S=\operatorname{Spec}(A)$ também são afins, com $A=\varinjlim A_\lambda$, e as afirmações de (i) equivalem ao

\begin{lemma}[8.8.2.1]
\label{IV.8.8.2.1}
Seja $A_0$ um anel, $(A_\lambda)_{\lambda\in L}$ um sistema indutivo de $A_0$-álgebras, $A=\varinjlim A_\lambda$; sejam $B_\alpha$ uma $A_\alpha$-álgebra e $C_\alpha$ uma $A_\alpha$-álgebra de tipo finito (resp.\ de apresentação finita).
Então o homomorfismo canônico
\[
  \varinjlim_\lambda\operatorname{Hom}_{A_\lambda\text{-alg.}}(C_\alpha\otimes_{A_\alpha}A_\lambda,B_\alpha\otimes_{A_\alpha}A_\lambda)\to\operatorname{Hom}_{A\text{-alg.}}(C_\alpha\otimes_{A_\alpha}A,B_\alpha\otimes_{A_\alpha}A)
  \tag{8.8.2.2}\label{IV.8.8.2.2}
\]
é injetivo (resp.\ bijetivo).
\end{lemma}

\begin{proof}
Sabemos que existem isomorfismos canônicos funtoriais
\[
  \operatorname{Hom}_{A_\lambda\text{-alg.}}(C_\alpha\otimes_{A_\alpha}A_\lambda,B_\alpha\otimes_{A_\alpha}A_\lambda)\xrightarrow{\sim}\operatorname{Hom}_{A_\alpha\text{-alg.}}(C_\alpha,B_\alpha\otimes_{A_\alpha}A_\lambda)
\]
\[
  \operatorname{Hom}_{A\text{-alg.}}(C_\alpha\otimes_{A_\alpha}A,B_\alpha\otimes_{A_\alpha}A)\xrightarrow{\sim}\operatorname{Hom}_{A_\alpha\text{-alg.}}(C_\alpha,B_\alpha\otimes_{A_\alpha}A)
\]
em virtude da propriedade universal do produto tensorial de duas álgebras.
Basta, portanto, demonstrar o
\end{proof}

\begin{lemma}[8.8.2.3]
\label{IV.8.8.2.3}
Sejam $E$ um anel, $G$ uma $E$-álgebra de tipo finito (resp.\ de apresentação finita) e $(F_\lambda)$ um sistema indutivo de $E$-álgebras.
Então o homomorfismo canônico
\[
  \varinjlim\operatorname{Hom}_{E\text{-alg.}}(G,F_\lambda)\to\operatorname{Hom}_{E\text{-alg.}}(G,\varinjlim F_\lambda)
\]
que associa, a todo sistema indutivo de homomorfismos $\theta_\lambda:G\to F_\lambda$ de $E$-álgebras, seu limite indutivo, é injetivo (resp.\ bijetivo).
\end{lemma}

\begin{proof}
Suponhamos primeiro que a $E$-álgebra $G$ seja de tipo finito, e seja $(t_i)_{1\leqslant i\leqslant n}$ um sistema de geradores dessa $E$-álgebra; demonstremos que, se $(\theta'_\lambda)$ e $(\theta''_\lambda)$ são dois sistemas indutivos de homomorfismos $G\to F_\lambda$, com $\varinjlim\theta'_\lambda=\varinjlim\theta''_\lambda$, existe $\mu\geqslant\lambda$ tal que $\theta'_\mu=\theta''_\mu$.
Com efeito, se $\varphi_{\mu\lambda}:F_\lambda\to F_\mu$ e $\varphi_\lambda:F_\lambda\to F=\varinjlim F_\lambda$ são os homomorfismos do sistema indutivo $(F_\lambda)$, por hipótese, para cada índice $i$, existe um índice $\lambda_i$ tal que $\varphi_{\lambda_i}(\theta'_{\lambda_i}(t_i))=\varphi_{\lambda_i}(\theta''_{\lambda_i}(t_i))$, e pode-se supor que todos os $\lambda_i$ sejam iguais a um único $\lambda$; deduz-se, do mesmo modo, a existência de $\mu\geqslant\lambda$ tal que $\varphi_{\mu\lambda}(\theta'_\lambda(t_i))=\varphi_{\mu\lambda}(\theta''_\lambda(t_i))$ para $1\leqslant i\leqslant n$, isto é, $\theta'_\mu(t_i)=\theta''_\mu(t_i)$ para $1\leqslant i\leqslant n$, de onde $\theta'_\mu=\theta''_\mu$.

Suponhamos, em segundo lugar, que $G$ seja de apresentação finita, de modo que se tenha
\[
  G=E[T_1,\ldots,T_n]/\mathfrak{J},
\]
onde $\mathfrak{J}$ é um ideal de tipo finito, sendo $t_i$ a classe de $T_i$ módulo $\mathfrak{J}$.
Seja $(P_j)_{1\leqslant j\leqslant m}$ um sistema de geradores de $\mathfrak{J}$.
Seja dado um homomorfismo de $E$-álgebras $\theta:G\to F$; ponhamos $b_i=\theta(t_i)$; por definição, tem-se então $P_j(b_1,b_2,\ldots,b_n)=\theta(P_j(t_1,\ldots,t_n))=0$ para $1\leqslant j\leqslant m$.
Ora, existem $\lambda$ e elementos $x_1,\ldots,x_n$ de $F_\lambda$ tais que $b_i=\varphi_\lambda(x_i)$ para $1\leqslant i\leqslant n$; tem-se, portanto, $\varphi_\lambda(P_j(x_1,\ldots,x_n))=P_j(b_1,\ldots,b_n)=0$, e, consequentemente, existe $\mu\geqslant\lambda$ tal que $\varphi_{\mu\lambda}(P_j(x_1,\ldots,x_n))=P_j(\varphi_{\mu\lambda}(x_1),\ldots,\varphi_{\mu\lambda}(x_n))=0$ para $1\leqslant j\leqslant m$; conclui-se que existe um homomorfismo de $E$-álgebras $\theta_\mu:G\to F_\mu$ tal que $\theta_\mu(t_i)=\varphi_{\mu\lambda}(x_i)$
\oldpage[IV-3]{30}
para $1\leqslant i\leqslant n$; deduz-se, portanto, para todo $\nu\geqslant\mu$, um homomorfismo de $E$-álgebras $\theta_\nu=\varphi_{\nu\mu}\circ\theta_\mu$ de $G$ em $F_\nu$, e é claro que $\theta$ é o limite indutivo desse sistema indutivo de homomorfismos, o que conclui a demonstração do lema.
\end{proof}

Passemos agora ao caso em que $X_\alpha$ é quase-compacto e $Y_\alpha$ localmente de tipo finito sobre $S_\alpha$.
Ponhamos $Z_\alpha=X_\alpha\times_{S_\alpha}Y_\alpha$ e introduzamos o sistema projetivo correspondente dos $Z_\lambda=Z_\alpha\times_{S_\alpha}S_\lambda=X_\lambda\times_{S_\lambda}Y_\lambda$ e seu limite $Z=Z_\alpha\times_{S_\alpha}S=X\times_S Y$; as bijeções canônicas \sref[I]{I.3.3.14} fornecem diagramas comutativos
\[
\xymatrix{
  \operatorname{Hom}_{S_\lambda}(X_\lambda,Y_\lambda) \ar[r]^-{e_{\mu\lambda}} \ar[d]^{\wr} & \operatorname{Hom}_{S_\mu}(X_\mu,Y_\mu) \ar[r]^-{e_\mu} \ar[d]^{\wr} & \operatorname{Hom}_S(X,Y) \ar[d]^{\wr}\\
  \operatorname{Hom}_{X_\lambda}(X_\lambda,Z_\lambda) \ar[r]_-{e_{\mu\lambda}} & \operatorname{Hom}_{X_\mu}(X_\mu,Z_\mu) \ar[r]_-{e_\mu} & \operatorname{Hom}_X(X,Z)
}
\]
e, por conseguinte, reduzimo-nos a demonstrar que \eqref{IV.8.8.1.1} é injetiva no caso particular em que $S_\alpha=X_\alpha$ (levando em conta \sref[I]{I.1.3.4}).
Além disso, como $X_\alpha$ é quase-compacto, e portanto reunião finita de abertos afins, podemos supor que $X_\alpha$ seja afim (pois $L$ é filtrante). Sejam, pois, dados dois $X_\alpha$-morfismos $f'_\alpha:X_\alpha\to Y_\alpha$, $f''_\alpha:X_\alpha\to Y_\alpha$ tais que $f'$ e $f''$ sejam $X$-morfismos iguais de $X$ em $Y$; devemos demonstrar que $f'_\mu=f''_\mu$ para $\mu\geqslant\alpha$ suficientemente grande. Como $X_\alpha$ é quase-compacto, também o é $f'_\alpha(X_\alpha)\cup f''_\alpha(X_\alpha)$ e, como $Y_\alpha$ é localmente de tipo finito sobre $X_\alpha$, essa união está contida em uma união finita de abertos afins $V_{i\alpha}$ de $Y_\alpha$, de tipo finito sobre $X_\alpha$. Ponhamos $U'_{i\alpha}=f_\alpha^{'-1}(V_{i\alpha})$, $U''_{i\alpha}=f_\alpha^{''-1}(V_{i\alpha})$, $U_{i\alpha}=U'_{i\alpha}\cap U''_{i\alpha}$ e $U_\alpha=\bigcup_iU_{i\alpha}$. A hipótese $f'=f''$ implica $v_\alpha^{-1}(U'_{i\alpha})=v_\alpha^{-1}(U''_{i\alpha})$, pois esses dois conjuntos são respectivamente iguais a $f^{'-1}(w_\alpha^{-1}(V_{i\alpha}))$ e $f^{''-1}(w_\alpha^{-1}(V_{i\alpha}))$. Como a união dos $V_{i\alpha}$ contém $f'_\alpha(X_\alpha)\cup f''_\alpha(X_\alpha)$, tem-se $v_\alpha^{-1}(U_\alpha)=X$. Como $X_\alpha$ é quase-compacto e todo aberto de $X_\lambda$ é ind-construtível, decorre de \sref{IV.8.3.4} que existe um índice $\mu\geqslant\alpha$ tal que os $U_{i\mu}=v_{\alpha\mu}^{-1}(U_{i\alpha})$ formem uma cobertura de $X_\mu$. Substituindo $\alpha$ por $\mu$, podemos, portanto, supor que os $U_{i\alpha}$ cobrem $X_\alpha$; isso implica que, para todo $x\in X_\alpha$, existe uma vizinhança aberta afim $W(x)$ contida em algum dos $U_{i\alpha}$, isto é, tal que as restrições de $f'_\alpha$ e $f''_\alpha$ a $W(x)$ levam $W(x)$ a um mesmo aberto afim $V_{i\alpha}$. Como $X_\alpha$ é quase-compacto, ele é coberto por um número finito de abertos afins $W(x_k)$; em virtude do lema~\sref{IV.8.8.2.1} e do fato de que $L$ é filtrante, existe $\lambda\geqslant\alpha$ tal que as restrições de $f'_\lambda$ e $f''_\lambda$ a cada um dos abertos $v_{\alpha\lambda}^{-1}(W(x_k))$ sejam iguais, de onde $f'_\lambda=f''_\lambda$.

Suponhamos agora que $X_\alpha$ seja quase-compacto e quase-separado e que $Y_\alpha$ seja localmente de apresentação finita sobre $S_\alpha$, e demonstremos que \eqref{IV.8.8.1.1} é sobrejetiva.
Seja, pois, dado um $X$-morfismo $f:X\to Y$.
Como $X$ é quase-compacto, também o é
\oldpage[IV-3]{31}
$f(X)$ e, consequentemente, existe um aberto quase-compacto $Y'$ de $Y$ que contém $f(X)$; existem então $\lambda\geqslant\alpha$ e um aberto quase-compacto $Y'_\lambda$ de $Y_\lambda$ tais que $Y'=w_\lambda^{-1}(Y'_\lambda)$ \sref{IV.8.2.11}.
Substituindo, se necessário, $\alpha$ por $\lambda$ e $Y_\lambda$ por $Y'_\lambda$, podemos, portanto, restringir-nos ao caso em que $Y_\alpha$ é \emph{quase-compacto}, e, por conseguinte, também $Y$ e os $Y_\lambda$.
Como $Y$ é quase-compacto, é uma união finita de abertos afins $V_i$, e, consequentemente, $X$ é a união dos abertos $f^{-1}(V_i)$. Como todo ponto de $X$ possui uma vizinhança aberta quase-compacta contida em algum dos $f^{-1}(V_i)$ e $X$ é quase-compacto, podemos, levando em conta \sref{IV.8.2.11} e substituindo $\alpha$, se necessário, por um índice $\lambda\geqslant\alpha$, supor que $Y$ seja uma união finita de abertos $V_i=w_\alpha^{-1}(V_{i\alpha})$, onde os $V_{i\alpha}$ são abertos afins de $Y_\alpha$; consequentemente, $X$ é a união dos abertos $f^{-1}(V_i)$. Como todo ponto de $X$ possui uma vizinhança aberta quase-compacta contida em algum dos $f^{-1}(V_i)$ e $X$ é quase-compacto, podemos cobrir $X$ por um número finito de tais vizinhanças e, repetindo alguns dos $V_i$, se necessário, supor que exista uma cobertura $(U_i)$ de $X$ por abertos quase-compactos, com o mesmo conjunto de índices que $(V_i)$ e tal que $f(U_i)\subset V_i$ para todo $i$. Além disso, utilizando \sref{IV.8.2.11} (e substituindo $\alpha$, se necessário, por um índice $\lambda\geqslant\alpha$), podemos supor que $U_i=v_\alpha^{-1}(U_{i\alpha})$, onde os $U_{i\alpha}$ são abertos quase-compactos de $X_\alpha$; utilizando também \sref{IV.8.3.4} como acima, podemos supor que $(U_{i\alpha})$ seja uma \emph{cobertura} de $X_\alpha$.

Bastará então demonstrar que existem $\lambda\geqslant\alpha$ e, para cada $i$, um morfismo $f_{i\lambda}:U_{i\lambda}\to V_{i\lambda}$ (onde $U_{i\lambda}=v_{\alpha\lambda}^{-1}(U_{i\alpha})$ e $V_{i\lambda}=w_{\alpha\lambda}^{-1}(V_{i\alpha})$) tal que o morfismo correspondente $f_i=e_\lambda(f_{i\lambda}):U_i\to V_i$ seja igual à \emph{restrição} de $f$ a $U_i$.
Com efeito, se assim for, como $X_\lambda$ é quase-separado \sref[I]{I.1.2.3}, os $U_{i\lambda}\cap U_{j\lambda}$ são \emph{quase-compactos}, e o resultado de unicidade demonstrado anteriormente (aplicável, pois $Y_\lambda$ é localmente de tipo finito sobre $S_\lambda$) mostra que existe $\mu\geqslant\lambda$ tal que $f_{i\mu}$ e $f_{j\mu}$ coincidam em $U_{i\mu}\cap U_{j\mu}$ para todo par $(i,j)$ e definam, consequentemente, um $S_\mu$-morfismo $f_\mu:X_\mu\to Y_\mu$ tal que $e_\mu(f_\mu)=f$.

Reduzimo-nos, assim, ao caso em que $Y_\alpha$ é \emph{afim} e, como, além disso, podemos supor que as imagens dos $V_{i\alpha}$ em $S_\alpha$ estejam contidas em abertos afins, também podemos supor que $S_\alpha$ seja afim; sejam, pois, $S_\alpha=\operatorname{Spec}(A_\alpha)$, $Y_\alpha=\operatorname{Spec}(C_\alpha)$, com $C_\alpha$ uma $A_\alpha$-álgebra de apresentação finita, e sejam $S=\operatorname{Spec}(A)$, $Y=\operatorname{Spec}(C)$, onde $A=\varinjlim A_\lambda$ e $C=\varinjlim(C_\alpha\otimes_{A_\alpha}A_\lambda)=C_\alpha\otimes_{A_\alpha}A$.
Tem-se então
\[
  \operatorname{Hom}_S(X,Y)=\operatorname{Hom}_{A\text{-alg.}}(C,\Gamma(X,\sh{O}_X))=\operatorname{Hom}_{A_\alpha\text{-alg.}}(C_\alpha,\Gamma(X,\sh{O}_X))
\]
\sref[I]{I.2.2.4} e, analogamente,
\[
  \operatorname{Hom}_{S_\lambda}(X_\lambda,Y_\lambda)=\operatorname{Hom}_{A_\lambda\text{-alg.}}(C_\alpha\otimes_{A_\alpha}A_\lambda,\Gamma(X_\lambda,\sh{O}_{X_\lambda}))=\operatorname{Hom}_{A_\alpha\text{-alg.}}(C_\alpha,\Gamma(X_\lambda,\sh{O}_{X_\lambda})).
\]
Mas, como $X_\alpha$ é quase-compacto e quase-separado, sabe-se por \sref{IV.8.5.4} que $\varinjlim\Gamma(X_\lambda,\sh{O}_{X_\lambda})=\Gamma(X,\sh{O}_X)$; como $C_\alpha$ é uma $A_\alpha$-álgebra de apresentação finita, a bijetividade de \eqref{IV.8.8.1.1} decorre então de \sref{IV.8.8.2.3}.

Antes de passar à demonstração de (ii), observemos os seguintes corolários de (i):
\end{proof}

\begin{corollary}[8.8.2.4]
\label{IV.8.8.2.4}
Suponhamos que $S_0$ seja quase-compacto, que $X_\alpha$ seja de apresentação finita sobre $S_\alpha$ e que $Y_\alpha$ seja de tipo finito sobre $S_\alpha$ e quase-separado sobre $S_\alpha$ (o que ocorre, por exemplo, se $Y_\alpha$ também for
\oldpage[IV-3]{32}
de apresentação finita sobre $S_\alpha$).
Seja $f_\alpha:X_\alpha\to Y_\alpha$ um $S_\alpha$-morfismo.
Para que $f:X\to Y$ seja um isomorfismo, é necessário e suficiente que exista $\lambda\geqslant\alpha$ tal que $f_\lambda:X_\lambda\to Y_\lambda$ seja um isomorfismo.
\end{corollary}

\begin{proof}
A condição é evidentemente suficiente. Para demonstrar que é necessária, observemos primeiro que, como a questão é local em $S_0$ (pois $L$ é filtrante), podemos sempre supor que $S_0$ seja afim e, portanto, quase-separado. Por hipótese, existe um $S$-morfismo $g:Y\to X$ tal que $g\circ f=1_X$ e $f\circ g=1_Y$. Como $X_\alpha$ é de apresentação finita sobre $S_\alpha$, e $Y_\alpha$ é quase-compacto e quase-separado (pois $S_\alpha$ é quase-compacto e quase-separado), decorre de \sref{IV.8.8.2},~(i) que existem $\mu\geqslant\alpha$ e um $S_\mu$-morfismo $g_\mu:Y_\mu\to X_\mu$ tais que $g=g_\mu\times 1_S$. Decorre ainda de \sref{IV.8.8.2},~(i), juntamente com as relações $g\circ f=1_X$ e $f\circ g=1_Y$, que existe $\nu\geqslant\mu$ tal que $g_\nu\circ f_\nu=1_{X_\nu}$ e $f_\nu\circ g_\nu=1_{Y_\nu}$, pois $X_\alpha$ e $Y_\alpha$ são de tipo finito sobre $S_\alpha$ e quase-compactos. Assim, $f_\nu:X_\nu\to Y_\nu$ é um isomorfismo, o que demonstra o corolário.
\end{proof}

\begin{corollary}[8.8.2.5]
\label{IV.8.8.2.5}
Suponhamos que $S_0$ seja quase-compacto e quase-separado, e que $X_\alpha$ e $Y_\alpha$ sejam de apresentação finita sobre $S_\alpha$.
Para que $X$ e $Y$ sejam $S$-isomorfos, é necessário e suficiente que exista $\lambda\geqslant\alpha$ tal que $X_\lambda$ e $Y_\lambda$ sejam $S_\lambda$-isomorfos. Além disso, para todo $S$-isomorfismo $f:X\xrightarrow{\sim}Y$, existem $\mu\geqslant\lambda$ e um $S_\mu$-isomorfismo $f_\mu:X_\mu\xrightarrow{\sim}Y_\mu$ tais que $f=f_\mu\times 1_S$.
\end{corollary}

\begin{proof}
A condição é evidentemente suficiente. Reciprocamente, se existe um $S$-isomorfismo $f:X\to Y$, decorre de \sref{IV.8.8.2},~(i) que $f$ é da forma $f_\mu\times 1_S$ para algum $\mu\geqslant\alpha$ e algum $S_\mu$-morfismo $f_\mu:X_\mu\to Y_\mu$. Como $f$ é um isomorfismo, decorre de \sref{IV.8.8.2.4} que existe $\nu\geqslant\mu$ tal que $f_\nu:X_\nu\to Y_\nu$ é um isomorfismo.
\end{proof}

\begin{proof}[Demonstração de \sref{IV.8.8.2}, (ii)]
Consideremos primeiro o caso em que $S_0=\operatorname{Spec}(A_0)$ e $X=\operatorname{Spec}(B)$ são afins; então a afirmação~(ii) é equivalente ao lema~\sref[I]{I.8.4.2}.

Para demonstrar (ii) no caso geral, observemos que $S$ é quase-compacto e quase-separado. Como $X$ é de apresentação finita sobre $S$, e $S$ é afim sobre $S_0$, existe uma cobertura aberta finita $(U_i)$ de $S_0$ e, se $(W_i)$ denota a cobertura de $S$ por abertos afins formada pelas imagens inversas dos $U_i$ pelo morfismo $S\to S_0$, uma cobertura finita $(X_r)$ de $X$ por abertos afins tal que a imagem de cada $X_r$ por $X\to S$ esteja contida em algum $W_{i(r)}$. O anel $A(X_r)$ é então uma álgebra de apresentação finita sobre $A(W_{i(r)})$ \sref[I]{I.1.4.6}.

Pelo lema~\sref[I]{I.8.4.2} e pelo fato de que $L$ é filtrante, existem um índice $\lambda\in L$ e, para todo $r$, um esquema afim $Z_{r\lambda}$ de apresentação finita sobre a imagem inversa $W_{i(r),\lambda}$ de $U_{i(r)}$ em $S_\lambda$, juntamente com um $S$-isomorfismo
\[
 g_r:Z_{r\lambda}\times_{S_\lambda}S\xrightarrow{\sim}X_r.
\]
Seja $Z_{rs}$ a imagem inversa por $g_r$ do subpré-esquema induzido pelo aberto $X_r\cap X_s$ em $X_r$, que é quase-compacto porque $X$ é quase-separado, e seja
\[
 g'_{rs}:Z_{rs}\xrightarrow{\sim}X_r\cap X_s
\]
a restrição de $g_r$. Por \sref{IV.8.2.11}, existem $\mu\geqslant\lambda$ e, para todo par $(r,s)$, um aberto quase-compacto $Z_{rs\mu}$ de $Z_{r\mu}=v_{\lambda\mu}^{-1}(Z_{r\lambda})$ tal que $Z_{rs}$ é a imagem inversa de $Z_{rs\mu}$. Além disso, como $S_\mu$ é quase-separado e $W_{i(r),\mu}$ é um aberto quase-compacto de $S_\mu$, cada $Z_{rs\mu}$ é de apresentação finita sobre $S_\mu$ \sref[I]{I.1.6.2}.

Para todo par $(r,s)$, consideremos o isomorfismo
\[
 h_{sr}=(g'_{sr})^{-1}\circ g'_{rs}:Z_{rs}\xrightarrow{\sim}Z_{sr}.
\]
Por \sref{IV.8.8.2.4}, existem $\nu\geqslant\mu$ e, para todo par $(r,s)$, um isomorfismo
\[
 h_{sr\nu}:Z_{rs\nu}\xrightarrow{\sim}Z_{sr\nu}
\]
tal que $h_{sr}=h_{sr\nu}\times 1_S$. Finalmente, para toda terna $(r,s,t)$, denotemos por $h'_{sr}$ a restrição de $h_{sr}$ a

\oldpage[IV-3]{33}

$Z_{rs}\cap Z_{rt}$. Decorre imediatamente das definições que $h'_{sr}$ é um isomorfismo de $Z_{rs}\cap Z_{rt}$ sobre $Z_{sr}\cap Z_{st}$ e que
\[
 h'_{ts}\circ h'_{sr}=h'_{tr}.
\]
Por \sref{IV.8.8.2},~(i), existe então $\rho\geqslant\nu$ tal que, para toda terna $(r,s,t)$,
\[
 h'_{ts\rho}\circ h'_{sr\rho}=h'_{tr\rho}.
\]
Consequentemente, os isomorfismos $h_{sr\rho}$ satisfazem a condição de colagem \sref[0\textsubscript{I}]{0.4.1.7} e definem, portanto, um pré-esquema $X_\rho$ tal que $X$ é isomorfo a $X_\rho\times_{S_\rho}S$. Além disso, os $Z_{r\rho}$ são de apresentação finita sobre $S_\rho$ e, identificando-os com abertos de $X_\rho$, suas interseções $Z_{r\rho}\cap Z_{s\rho}$, que são isomorfas aos $Z_{rs\rho}$, são quase-compactas. Assim, $X_\rho$ é de apresentação finita sobre $S_\rho$ \sref[I]{I.1.6.3}.
\end{proof}

\begin{env}[8.8.3]
\label{IV.8.8.3}
O conteúdo essencial de \sref{IV.8.8.2} exprime-se também dizendo que, se $S_0$ é quase-compacto e quase-separado, a categoria dos $S$-pré-esquemas de apresentação finita é determinada, a menos de equivalência, pelos dados das categorias dos $S_\lambda$-pré-esquemas de apresentação finita, dos funtores $\rho_{\mu\lambda}:X_\lambda\mapsto X_\lambda\times_{S_\lambda}S_\mu$ (para $\lambda\leqslant\mu$) entre essas categorias e dos isomorfismos de transitividade $\rho_{\nu\mu}\circ\rho_{\mu\lambda}\xrightarrow{\sim}\rho_{\nu\lambda}$.
De maneira figurada, pode-se dizer que o dado de um $S$-pré-esquema de apresentação finita $X$ é ``funtorialmente'' equivalente ao dado de um $S_\lambda$-pré-esquema de apresentação finita $X'_\lambda$; se um $S_\mu$-pré-esquema de apresentação finita $X''_\mu$ é tal que $X$ é $S$-isomorfo a $X''_\mu\times_{S_\mu}S$, existe $\nu$ tal que $\lambda\leqslant\nu$, $\mu\leqslant\nu$, e $X'_\lambda\times_{S_\lambda}S_\nu$ e $X''_\mu\times_{S_\mu}S_\nu$ são $S_\nu$-isomorfos.
O fato de $L$ ser filtrante implica, além disso, que, se forem dadas uma \emph{família} finita $(X^{(i)})_{i\in I}$ de $S$-pré-esquemas de apresentação finita e uma família finita $(f^{(j)})_{j\in J}$ de $S$-morfismos entre esses pré-esquemas (sendo, portanto, $f^{(j)}$ da forma $X^{(\sigma(j))}\to X^{(\tau(j))}$, onde $\sigma$ e $\tau$ são duas aplicações de $J$ em $I$), então existem um índice $\lambda\in L$, uma família $(X^{(i)}_\lambda)_{i\in I}$ de $S_\lambda$-pré-esquemas de apresentação finita e uma família $(f^{(j)}_\lambda)_{j\in J}$ de $S_\lambda$-morfismos tais que $X^{(i)}$ se identifica com $X^{(i)}_\lambda\times_{S_\lambda}S$ e $f^{(j)}$ com $f^{(j)}_\lambda\times 1_S$ para todos os $i$ e $j$; além disso, se houver uma relação $f^{(j)}=f^{(k)}$, pode-se supor $\lambda$ escolhido de modo que $f^{(j)}_\lambda=f^{(k)}_\lambda$.

Consideremos, em particular, três $S$-pré-esquemas de apresentação finita $X$, $Y$, $Z$ e dois $S$-morfismos $f:X\to Z$, $g:Y\to Z$, de modo que o produto $X\times_Z Y$ relativo a esses morfismos seja também um $S$-pré-esquema de apresentação finita \sref[I]{I.1.6.2}.
Resulta então do que precede e de \sref[I]{I.3.3.11} que se pode escrever $X\times_Z Y=(X_\lambda\times_{Z_\lambda}Y_\lambda)\times_{S_\lambda}S$ para um $\lambda$ conveniente, sendo $X_\lambda$, $Y_\lambda$, $Z_\lambda$ $S_\lambda$-pré-esquemas de apresentação finita; pode-se, portanto, dizer que a determinação dos $S$-pré-esquemas de apresentação finita pelos dados dos $S_\lambda$-pré-esquemas de apresentação finita é ``\emph{compatível com os produtos fibrados}''.
Por outro lado, vimos em \sref{IV.8.6.3} que, se $g$ é uma imersão, pode-se supor o mesmo para $g_\lambda$; identificando $Y$ (resp.\ $Y_\lambda$) com um subpré-esquema de $Z$ (resp.\ $Z_\lambda$), vê-se que se pode escrever, para um $\lambda$ conveniente, $f^{-1}(Y)=f_\lambda^{-1}(Y_\lambda)\times_{S_\lambda}S$ \sref[I]{I.4.4.1};
há, portanto, também ``\emph{compatibilidade com a formação de imagens inversas de subpré-esquemas}''.
Mais particularmente, se $f_1$, $f_2$ são dois $S$-morfismos de $X$ em $Y$, chama-se \emph{núcleo} desse par de morfismos a imagem inversa $N$ do subpré-esquema diagonal de $Y\times_S Y$ pelo $S$-morfismo $(f_1,f_2)_S$; deduz-se do que precede que se terá, para um $\lambda$ conveniente, $N=N_\lambda\times_{S_\lambda}S$, onde $N_\lambda$ é o núcleo do par de $S_\lambda$-morfismos $(f_{1\lambda},f_{2\lambda})$ de $X_\lambda$ em $Y_\lambda$.
Essas observações estendem-se aos produtos finitos e aos ``núcleos'' de sistemas finitos arbitrários de $S$-morfismos de um $S$-pré-esquema de apresentação finita em
\oldpage[IV-3]{34}
outro; conclui-se que, de maneira geral, a formação de $S$-pré-esquemas de apresentação finita pelos dados de $S_\lambda$-pré-esquemas de apresentação finita é ``\emph{compatível com os limites projetivos finitos}'', sendo tal limite, por definição, o núcleo de um sistema finito de morfismos de um $S$-pré-esquema em um produto de $S$-pré-esquemas \sref[T]{T.1.8}.

No que segue, permitir-nos-emos efetuar livremente as traduções implicadas pelas propriedades precedentes (bem como por \sref{IV.8.3.11}, \sref{IV.8.5.2} e \sref{IV.8.6.3}), sem sempre nos empenharmos em justificá-las como acima.
Por exemplo, o dado de um ``pré-esquema em grupos'' $G$ de apresentação finita sobre $S$ equivale ao dado de um pré-esquema em grupos $G_\lambda$ de apresentação finita sobre algum $S_\lambda$, para um $\lambda$ suficientemente grande; com efeito, escrever a condição de associatividade para o $S$-morfismo ``lei de composição'' $G\times_S G\to G$ do pré-esquema em grupos equivale a escrever que o núcleo de dois $S$-morfismos da forma $G\times_S G\times_S G\to G$ é igual a $G\times_S G\times_S G$ \sref[II]{II.8.3.9}, e as demais condições que intervêm na definição de um pré-esquema em grupos interpretam-se de maneira análoga.
\end{env}

\subsection{Primeiras aplicações à eliminação das hipóteses noetherianas}
\label{subsection:IV.8.9}

\begin{proposition}[8.9.1]
\label{IV.8.9.1}
Seja $A$ um anel e seja $X$ um $A$-pré-esquema.
\begin{enumerate}[label=(\roman*)]
\item As condições seguintes são equivalentes:
\begin{enumerate}[label=\emph{\alph*)}]
\item $X$ é de apresentação finita sobre $A$.
\item Existem um anel noetheriano $A_0$, um pré-esquema $X_0$ de tipo finito sobre $A_0$, um homomorfismo de anéis $A_0\to A$ e um $A$-isomorfismo $X_0\otimes_{A_0}A\xrightarrow{\sim}X$.
\item Existem um subanel $A_0$ de $A$, que é uma $\mathbf{Z}$-álgebra de tipo finito, um pré-esquema $X_0$ de tipo finito sobre $A_0$ e um $A$-isomorfismo $X_0\otimes_{A_0}A\xrightarrow{\sim}X$.
\end{enumerate}
\item Se $\sh{F}$ é um $\sh{O}_X$-Módulo quase-coerente de apresentação finita, pode-se supor o subanel $A_0$ escolhido de modo que exista um $\sh{O}_{X_0}$-Módulo coerente $\sh{F}_0$ tal que $\sh{F}$ seja isomorfo a $\sh{F}_0\otimes_{A_0}A$; $\operatorname{Supp}(\sh{F})$ é construtível e fechado em $X$, e existe um subpré-esquema fechado $Z$ de $X$, cujo espaço subjacente é $\operatorname{Supp}(\sh{F})$, tal que a imersão canônica $Z\to X$ seja de apresentação finita.
\item Se $Y$ é outro $A$-pré-esquema de apresentação finita e $f:X\to Y$ é um $A$-morfismo, pode-se supor o subanel $A_0$ de $A$ escolhido de modo que existam um pré-esquema $Y_0$ de tipo finito sobre $A_0$, um $A$-isomorfismo $Y_0\otimes_{A_0}A\xrightarrow{\sim}Y$ e um $A_0$-morfismo $f_0:X_0\to Y_0$ (necessariamente de tipo finito) tal que $f$ se identifique com $f_0\otimes 1_A$.
\end{enumerate}
\end{proposition}

\begin{proof}
\begin{enumerate}[label=(\roman*)]
\item Como $A$ é o limite indutivo de seus subanéis de tipo finito sobre $\mathbf{Z}$, \emph{a)} implica \emph{c)} em virtude de \sref{IV.8.8.2},~(ii); \emph{c)} implica \emph{b)}, pois uma $\mathbf{Z}$-álgebra de tipo finito é um anel noetheriano; por fim, se $A_0$ é noetheriano, todo $A_0$-pré-esquema de tipo finito é de apresentação finita \sref[I]{I.1.6.1}, de modo que \emph{b)} implica \emph{a)} em virtude de \sref[I]{I.1.6.2},~(iii).
\end{enumerate}
A asserção~(ii) decorre imediatamente de \sref{IV.8.5.2},~(ii), \sref{IV.8.3.11} e \sref[I]{I.1.8.2}, e a asserção~(iii), de \sref{IV.8.8.2},~(i).
\end{proof}

\begin{env}[8.9.2]
\label{IV.8.9.2}
A proposição \sref{IV.8.9.1} e os resultados de \sref{IV.8.5.2} e \sref{IV.8.8.2} permitem estender, em numerosos casos, aos morfismos de apresentação finita $X\to Y$ os resultados demonstrados por meio de técnicas noetherianas, supondo $Y$ \emph{localmente noetheriano}.
\oldpage[IV-3]{35}
Veremos numerosos exemplos no que segue e limitar-nos-emos aqui a apresentar alguns resultados desse tipo.
\end{env}

\begin{proposition}[8.9.3]
\label{IV.8.9.3}
Seja $A$ um anel e seja $M$ um $A$-módulo de apresentação finita.
Então todo $A$-endomorfismo sobrejetivo de $M$ é bijetivo.
\end{proposition}

\begin{proof}
Com efeito, consideremos $A$ como limite indutivo de suas sub-$\mathbf{Z}$-álgebras de tipo finito.
Segue de \sref{IV.8.5.2},~(ii) que existem uma dessas subálgebras $A_0$ e um $A_0$-módulo $M_0$ de apresentação finita tais que $M$ é $A$-isomorfo a $M_0\otimes_{A_0}A$; além disso, se $f:M\to M$ é um $A$-endomorfismo sobrejetivo, pode-se supor \sref{IV.8.5.2},~(i) que existe um $A_0$-endomorfismo $f_0:M_0\to M_0$ tal que $f=f_0\otimes 1_A$; por fim, em virtude de \sref{IV.8.5.7}, pode-se supor que $f_0$ é \emph{sobrejetivo}.
Mas, como $A_0$ é noetheriano e $M_0$ é um $A_0$-módulo de tipo finito, $M_0$ é um módulo noetheriano; portanto (Bourbaki, \emph{Alg.}, cap.~VIII, \textsection\,2, n\textsuperscript{o}\,2, lema~3), $f_0$ é bijetivo, e o mesmo ocorre com $f$ por mudança de base.
\end{proof}

\begin{proposition}[8.9.4]
\label{IV.8.9.4}
(``teorema da planicidade genérica''). --- Seja $Y$ um pré-esquema íntegro, seja $u:X\to Y$ um morfismo de tipo finito e localmente de apresentação finita, e seja $\sh{F}$ um $\sh{O}_X$-Módulo quase-coerente de apresentação finita.
Então existe um aberto não vazio $U$ de $Y$ tal que $\sh{F}|u^{-1}(U)$ é plano sobre $U$.
\end{proposition}

\begin{proof}
O raciocínio do início de \sref{IV.6.9.1} conduz a demonstrar o

\begin{lemma}[8.9.4.1]
\label{IV.8.9.4.1}
Seja $A$ um anel íntegro, seja $B$ uma $A$-álgebra de apresentação finita e seja $M$ um $B$-módulo de apresentação finita.
Então existe $f\neq 0$ em $A$ tal que $M_f$ é um $A_f$-módulo livre.
\end{lemma}

\begin{proof}
Com efeito, por \sref{IV.8.9.1} existem uma $\mathbf{Z}$-subálgebra $A_0$ de $A$ de tipo finito, uma $A_0$-álgebra $B_0$ de tipo finito e um $B_0$-módulo $M_0$ de tipo finito tais que $B$ é $A$-isomorfo a $B_0\otimes_{A_0}A$ e $M$ é $B$-isomorfo a $M_0\otimes_{A_0}A$; por \sref{IV.6.9.2}, existe $f_0\neq 0$ em $A_0$ tal que $(M_0)_{f_0}$ é um $(A_0)_{f_0}$-módulo livre.
Como $M_{f_0}=(M_0)_{f_0}\otimes_{A_0}A$ e $A_{f_0}=(A_0)_{f_0}\otimes_{A_0}A$, $M_{f_0}$ é um $A_{f_0}$-módulo livre.
\end{proof}
\end{proof}

\begin{corollary}[8.9.5]
\label{IV.8.9.5}
Seja $S$ um pré-esquema quase-compacto e quase-separado, seja $u:X\to S$ um morfismo de apresentação finita e seja $\sh{F}$ um $\sh{O}_X$-Módulo quase-coerente de apresentação finita.
Então existe uma partição $(S_i)_{1\leqslant i\leqslant n}$ de $S$ em um número finito de subconjuntos localmente fechados de $S$ tal que, para $1\leqslant i\leqslant n$, existe um subpré-esquema $S'_i$ de $S$, com espaço subjacente $S_i$, que é de apresentação finita sobre $S$, e tal que, pondo $X_i=X\times_S S'_i$, o $\sh{O}_{X_i}$-Módulo $\sh{F}_i=\sh{F}\otimes_{\sh{O}_X}\sh{O}_{X_i}$ é plano sobre $S'_i$.
\end{corollary}

\begin{proof}
Consideremos uma cobertura finita $(U_j)_{1\leqslant j\leqslant m}$ de $S$ por abertos afins e definamos, por indução, $T_1=U_1$ e
\[
T_j=U_j\setminus\bigcup_{h<j}(U_j\cap U_h)\qquad(j\geqslant 2).
\]
Cada $T_j$ é fechado no aberto afim $U_j$, e os $T_j$ são disjuntos dois a dois.
Além disso, os $U_j\cap U_h$ são quase-compactos porque $S$ é quase-separado e, portanto, como $S$ também é quase-compacto, são construtíveis em $S$ \sref{IV.1.8.1}; consequentemente, cada $T_j$ é igualmente construtível.
Bastará, portanto, demonstrar o corolário para um subpré-esquema adequado $T'_j$ de $S$, com espaço subjacente $T_j$ e de apresentação finita sobre $S$, e para o morfismo e o Módulo deduzidos, respectivamente, de $u$ e de $\sh{F}$ pela mudança de base $T'_j\to S$.
Para isso, observemos que, se dotarmos $U_j$ da estrutura de pré-esquema induzida pela de $S$, a imersão aberta $U_j\to S$ é quase-compacta porque $S$ é quase-separado \sref{IV.1.2.7} e, portanto, é de
\oldpage[IV-3]{36}
apresentação finita \sref{IV.1.6.2},~(i).
Basta, portanto, que $T'_j$ seja de apresentação finita sobre $U_j$; em outras palavras, podemos desde já reduzir-nos ao caso em que $U_j=S$ e $T_j=T$ é fechado e construtível em $S$.

Escrevamos $S=\operatorname{Spec}(A)$ e consideremos $A$ como limite indutivo de suas $\mathbf{Z}$-subálgebras de tipo finito, de modo que $S=\varprojlim S_\lambda$, onde os $S_\lambda$ são afins e noetherianos.
Por \sref{IV.8.3.11}, existem $\lambda$ e um subconjunto fechado (necessariamente construtível) $T_\lambda$ de $S_\lambda$ tais que $T=u_\lambda^{-1}(T_\lambda)$, onde $u_\lambda:S\to S_\lambda$ é o morfismo canônico.
Dotemos $T_\lambda$ de uma estrutura de subpré-esquema e tomemos $T'=T_\lambda\times_{S_\lambda}S$.
Como a imersão canônica $T_\lambda\to S_\lambda$ é de apresentação finita \sref{IV.1.6.1}, o mesmo vale para a imersão $T'\to S$.
Como $T'$ é afim, vemos finalmente que basta tratar o caso em que $T'=S$ é afim.

Então, com as notações de \sref{IV.8.9.1}, existem $\lambda$, um morfismo de tipo finito $u_\lambda:X_\lambda\to S_\lambda$ e um $\sh{O}_{X_\lambda}$-Módulo coerente $\sh{F}_\lambda$ tais que $X$ é isomorfo a $X_\lambda\times_{S_\lambda}S$ e $\sh{F}$ a $\sh{F}_\lambda\otimes_{\sh{O}_{X_\lambda}}\sh{O}_X$.
Aplicando \sref{IV.6.9.3} a $S_\lambda$, $X_\lambda$ e $\sh{F}_\lambda$, obtemos subpré-esquemas $S'_{\lambda i}$ de $S_\lambda$ cujos conjuntos subjacentes formam uma partição de $S_\lambda$ e tais que, pondo $X_{\lambda i}=X_\lambda\times_{S_\lambda}S'_{\lambda i}$, o Módulo
\[
\sh{F}_{\lambda i}=\sh{F}_\lambda\otimes_{\sh{O}_{X_\lambda}}\sh{O}_{X_{\lambda i}}
\]
é plano sobre $S'_{\lambda i}$.
Os subpré-esquemas $S'_i=S'_{\lambda i}\times_{S_\lambda}S$ de $S$ possuem então as propriedades requeridas em virtude de \sref{IV.2.1.4}.
\end{proof}

\subsection{Propriedades de permanência dos morfismos por passagem ao limite projetivo}
\label{subsection:IV.8.10}

Neste parágrafo, conservamos as hipóteses gerais e as notações de \sref{IV.8.8.1}.

\begin{proposition}[8.10.1]
\label{IV.8.10.1}
Se existe $\lambda$ tal que, para $\mu\geqslant\lambda$, $f_\mu$ é um morfismo aberto (resp.\ universalmente aberto), então $f$ é um morfismo aberto (resp.\ universalmente aberto).
\end{proposition}

\begin{proof}
Como $X=X_\lambda\times_{Y_\lambda}Y$, a asserção relativa aos morfismos universalmente abertos é um caso particular de \sref{IV.2.4.3},~(iii).
Suponhamos, pois, que $f_\mu$ seja aberto para $\mu\geqslant\lambda$; basta ver que, para todo aberto \emph{quase-compacto} $U$ de $X$, $f(U)$ é aberto em $Y$.
Ora, existe $\mu\geqslant\lambda$ tal que $U=v_\mu^{-1}(U_\mu)$, onde $U_\mu$ é um aberto quase-compacto de $X_\mu$ \sref{IV.2.3.11}; tem-se então $f(v_\mu^{-1}(U_\mu))=w_\mu^{-1}(f_\mu(U_\mu))$ \sref[I]{I.3.4.8}, e, portanto, $f(U)$ é aberto em $Y$.
\end{proof}

\begin{corollary}[8.10.2]
\label{IV.8.10.2}
Seja $f:X\to Y$ um morfismo.
Para que $f$ seja universalmente aberto, basta que, para todo inteiro $n>0$, pondo $Y_n=Y\otimes_\mathbf{Z}\mathbf{Z}[T_1,\ldots,T_n](=\mathbf{V}_Y^n)$, $X_n=X\times_Y Y_n$ e $f_n=f\times 1_{Y_n}:X_n\to Y_n$, a projeção $f_n$ seja um morfismo aberto.
\end{corollary}

\begin{proof}
Para demonstrar que $f$ é universalmente aberto, basta demonstrar que continua a sê-lo sobre todo aberto afim $U$ de $Y$; por hipótese, se $U_n=U\otimes_\mathbf{Z}\mathbf{Z}[T_1,\ldots,T_n]$ é a imagem inversa de $U$ em $Y_n$, o morfismo $f_n^{-1}(U_n)\to U_n$, restrição de $f_n$, é aberto, e vê-se que podemos reduzir-nos ao caso em que $Y=\operatorname{Spec}(A)$ é afim.
Além disso, basta evidentemente demonstrar que, para todo morfismo $Y'\to Y$, onde $Y'=\operatorname{Spec}(A')$ é também \emph{afim}, $f'=f_{|Y'}$ é aberto.
Suponhamos primeiro que $A'$ seja uma $A$-álgebra de tipo finito e, portanto, um \emph{quociente} de uma álgebra de polinômios $A[T_1,\ldots,T_n]$; então $Y'$ é um subpré-esquema fechado de $Y_n$ e $f'$ é a restrição
\oldpage[IV-3]{37}
de $f_n$ a $f_n^{-1}(Y')$; mas, para todo aberto $V$ de $X_n$, tem-se $f_n(V\cap f_n^{-1}(Y'))=f_n(V)\cap Y'$, e, como por hipótese $f_n(V)$ é aberto em $Y_n$, isso demonstra que $f'$ também é um morfismo aberto.
Quando $A'$ é arbitrária, pode ser considerada como limite indutivo de suas sub-$A$-álgebras de tipo finito $A'_\lambda$, e o fato de $f'$ ser um morfismo aberto decorre do que precede e de \sref{IV.8.10.1}.
\end{proof}

\begin{proposition}[8.10.3]
\label{IV.8.10.3}
Suponhamos que exista $\alpha$ tal que: $1^\circ$ $S_\alpha$ seja quase-compacto; $2^\circ$ os morfismos $X_\alpha\to S_\alpha$, $Y_\alpha\to S_\alpha$ sejam quase-compactos e o morfismo $Y_\alpha\to S_\alpha$ seja quase-separado; $3^\circ$ para $\alpha\leqslant\lambda\leqslant\mu$, os morfismos $u_{\lambda\mu}:S_\mu\to S_\lambda$ sejam planos; $4^\circ$ $\overline{f_\alpha(X_\alpha)}$ seja construtível em $Y_\alpha$.
Então, para que $f$ seja dominante, é necessário e suficiente que exista $\lambda\geqslant\alpha$ tal que $f_\lambda$ seja dominante.
\end{proposition}

\begin{proof}
As hipóteses implicam que $Y_\alpha$ é quase-compacto e que o morfismo $f_\alpha$ é quase-compacto \sref{IV.1.2.4}; portanto, $\overline{f_\alpha(X_\alpha)}=Z_\alpha$ é pró-construtível \sref{IV.1.9.5},~(vii) em $Y_\alpha$.
Se pusermos $Z_\lambda=w_{\alpha\lambda}^{-1}(Z_\alpha)$ para $\lambda\geqslant\alpha$ e $Z=w_\alpha^{-1}(Z_\alpha)$, teremos $Z_\lambda=\overline{f_\lambda(X_\lambda)}$ e $Z=\overline{f(X)}$ \sref[I]{I.3.4.8}, e $Z_\lambda$ será pró-construtível em $Y_\lambda$ \sref{IV.1.9.5},~(vii).
Basta então aplicar \sref{IV.8.3.13}, substituindo $S_\lambda$, $Z'_\lambda$ e $Z''_\lambda$ por $Y_\lambda$, $Y_\lambda$ e $Z_\lambda$, respectivamente.
\end{proof}

\begin{proposition}[8.10.4]
\label{IV.8.10.4}
Suponhamos que exista $\alpha$ tal que $Y_\alpha$ seja quase-compacto e $f_\alpha$ seja de tipo finito e quase-separado.
Para que o morfismo $f$ seja separado, é necessário e suficiente que exista $\lambda\geqslant\alpha$ tal que $f_\lambda$ seja separado.
\end{proposition}

\begin{proof}
Como a questão é local sobre $Y_\alpha$ (pois $Y_\alpha$ é quase-compacto e $L$ é filtrante), podemos reduzir-nos ao caso em que $Y_\alpha$ é afim e, portanto, quase-separado; a hipótese implica então que $X_\alpha$ (e, por conseguinte, os $X_\lambda$ e $X$) é quase-compacto e quase-separado.
Ponhamos $X'_\lambda=X_\lambda\times_{Y_\lambda}X_\lambda$ para $\lambda\geqslant\alpha$ e $X'=X\times_Y X$; tem-se $X'_\lambda=X'_\alpha\times_{Y_\alpha}Y_\lambda$ e $X'=X'_\alpha\times_{Y_\alpha}Y$. A primeira projeção $X'_\alpha\to X_\alpha$ é quase-compacta e quase-separada por hipótese \sref{IV.1.2.3},~(iii) e, portanto, $X'_\lambda$ é quase-compacto e quase-separado.
Observemos agora que, se denotarmos por $\Delta_\lambda$ (resp.\ $\Delta$) a diagonal de $X_\lambda\times_{Y_\lambda}X_\lambda$ (resp.\ $X\times_Y X$), segue de \sref[I]{I.5.3.4} e \sref[I]{I.3.4.8} que $\Delta_\mu$ (resp.\ $\Delta$) é a imagem inversa de $\Delta_\lambda$ pela aplicação $v'_{\lambda\mu}:X'_\mu\to X'_\lambda$ (resp.\ $v'_\lambda:X'\to X'_\lambda$).
Por outro lado, $\Delta_\lambda$ é \emph{construtível} em $X'_\lambda$: com efeito, como $f_\lambda$ é quase-separado, a imersão diagonal $X_\lambda\to X'_\lambda$ é quase-compacta e localmente de apresentação finita porque $f_\lambda$ é de tipo finito \sref[I]{I.1.4.3} e \sref[I]{I.1.5.4},~(iii); segue então de \sref{IV.1.8.4.1} que $\Delta_\lambda$ é localmente construtível e, portanto, construtível, já que $X'_\lambda$ é quase-compacto.
Podemos agora aplicar \sref{IV.8.3.12}, substituindo $S_\lambda$ e $Z_\lambda$ por $X'_\lambda$ e $\Delta_\lambda$, respectivamente.
\end{proof}

\begin{theorem}[8.10.5]
\label{IV.8.10.5}
Suponhamos $S_0$ quase-compacto, $X_\alpha$ e $Y_\alpha$ de apresentação finita sobre $S_\alpha$, e seja $f_\alpha:X_\alpha\to Y_\alpha$ um $S_\alpha$-morfismo.
Consideremos, para um morfismo, a propriedade de ser:
\begin{enumerate}[label=(\roman*)]
\item um isomorfismo;
\item[\rm(i\,\emph{bis})] um monomorfismo;
\item uma imersão;
\item uma imersão aberta;
\item uma imersão fechada;
\item separado;
\item sobrejetivo;
\oldpage[IV-3]{38}
\item radicial;
\item afim;
\item quase-afim;
\item finito;
\item quase-finito;
\item próprio.
\end{enumerate}
Então, se $\mathbf{P}$ denota uma das propriedades precedentes, para que $f$ possua a propriedade $\mathbf{P}$, é necessário e suficiente que exista $\lambda\geqslant\alpha$ tal que $f_\lambda$ possua a propriedade $\mathbf{P}$ (caso em que $f_\mu$ também a possui para $\mu\geqslant\lambda$).

Se, além disso, supusermos $S_0$ quase-separado, a mesma conclusão vale quando $\mathbf{P}$ é a propriedade de ser:
\begin{enumerate}[label=(\roman*)]
\setcounter{enumi}{12}
\item projetivo;
\item quase-projetivo.
\end{enumerate}
\end{theorem}

\begin{proof}
O caso em que $\mathbf{P}$ é uma das propriedades~(i) ou~(v) foi incluído no enunciado apenas a título de recordação; nesses casos, o teorema decorre do que foi demonstrado, respectivamente, em \sref{IV.8.8.2.4} e \sref{IV.8.10.4}.
Além disso, levando em conta \sref[I]{I.5.4.1} e \sref[I]{I.5.3.4}, (v) também decorre de (iv).
O caso (i\,\emph{bis}) decorre imediatamente de (i), utilizando \sref[I]{I.5.3.8} e observando (como já utilizamos em \sref{IV.8.10.4}) que a diagonal $\Delta_f$ é deduzida de $\Delta_{f_\lambda}$ pela mudança de base $S\to S_\lambda$.

Observaremos, por outro lado, que (vi), (vii) e (xi) são, na realidade, condições sobre as \emph{fibras} $f^{-1}(y)$ dos morfismos considerados, levando em conta a transitividade das fibras por mudança de base \sref[I]{I.3.6.4}: a condição~(vi) significa, com efeito, que todas as fibras devem ser não vazias; a condição~(vii), que devem ser radiciais \sref[I]{I.3.5.8}; e a condição~(xi), que devem ser finitas \sref[III]{III.6.2.2}, \sref[III]{III.6.2.3} e \sref[I]{I.6.4.4}, levando em conta que $f$ e os $f_\lambda$ são morfismos de tipo finito por \sref{IV.1.5.4},~(v).
O teorema nesses três casos será ainda consequência de um resultado geral sobre esse tipo de propriedades, que será estabelecido em \sref{IV.9.3.3}; adiamos, portanto, até então a demonstração do teorema no caso~(xi) (o leitor poderá verificar, naturalmente, que, exceto nos n\textsuperscript{os}~{\hyperlink{ega4.c04.subsec.8-11}{8.11}} e {\hyperlink{ega4.c04.subsec.8-12}{8.12}}, não utilizaremos o teorema nesses casos antes de \sref{IV.9.3.3}, e que \sref{IV.8.11} e \sref{IV.8.12} não serão utilizados antes de \sref{IV.9.3.3}).

Para os casos que restam demonstrar, podemos limitar-nos a provar que a condição do enunciado é necessária, pois todas as propriedades $\mathbf{P}$ consideradas são invariantes por mudança de base (vejam-se os caps.~I\textsuperscript{o} e II para as referências relativas a cada uma dessas propriedades).
Podemos supor, além disso, que $S_\alpha=S_0$ e que $Y_\alpha=S_\alpha$ e, portanto, $Y_\lambda=S_\lambda$ para todo $\lambda\geqslant\alpha$.
Por fim, as propriedades~(i) a~(xii) são locais sobre $S_0$; assim, como $S_0$ é união finita de abertos afins e $L$ é filtrante, podemos limitar-nos a demonstrá-las para $S_0=\operatorname{Spec}(A_0)$ \emph{afim} (e, portanto, quase-separado).
Denotaremos por $A_\lambda$ (resp.\ $A$) o anel de $S_\lambda$ (resp.\ $S$).

\emph{Casos} (ii), (iii), (iv): Suponhamos que $f$ seja uma imersão (resp.\ uma imersão aberta, uma imersão fechada), e seja $X'$ o subpré-esquema (resp.\ induzido sobre um aberto, resp.\ fechado) de $S$ associado a $f$, que é, portanto, um $S$-pré-esquema de apresentação finita.
\oldpage[IV-3]{39}
Em virtude de \sref{IV.8.6.3}, existem então $\lambda\geqslant\alpha$ e um subpré-esquema $X'_\lambda$ de $S_\lambda$, de apresentação finita sobre $S_\lambda$, tais que $X'$ é isomorfo a $X'_\lambda\times_{S_\lambda}S$.
Para todo $\mu\geqslant\lambda$, $X'_\mu=X'_\lambda\times_{S_\lambda}S_\mu$ é, portanto, um subpré-esquema (resp.\ induzido sobre um aberto, resp.\ fechado) de $S_\mu$, e segue então de \sref{IV.8.8.2.4} e \sref{IV.8.8.2.5} que existem $\mu\geqslant\lambda$ e um isomorfismo $g_\mu:X_\mu\to X'_\mu$ tais que $g=g_\mu\times 1_S$ é o isomorfismo $X\to X'$ associado a $f$; daí a conclusão.

\emph{Casos} (vi) \emph{e} (vii): Sabemos por \sref{IV.1.8.4.1} que $Z_\alpha=f_\alpha(X_\alpha)$ é construtível em $S_\alpha$; se pusermos $Z_\lambda=u_{\alpha\lambda}^{-1}(Z_\alpha)$ para $\lambda\geqslant\alpha$ e $Z=u_\alpha^{-1}(Z_\alpha)$, teremos $Z_\lambda=f_\lambda(X_\lambda)$ e $Z=f(X)$ \sref[I]{I.3.4.8}.
Como $Z_\alpha$ é pró-construtível em $S_\alpha$, e como, em virtude de \sref{IV.8.3.11}, a aplicação canônica $\varinjlim\mathfrak{C}(S_\lambda)\to\mathfrak{C}(S)$ é injetiva, a relação $f(X)=S$ implica a existência de $\lambda\geqslant\alpha$ tal que $f_\lambda(X_\lambda)=S_\lambda$, o que demonstra o teorema no caso~(vi).
Para demonstrá-lo no caso~(vii), observemos que o morfismo estrutural $X_\alpha\times_{S_\alpha}X_\alpha\to S_\alpha$ é de apresentação finita, pois a primeira projeção $X_\alpha\times_{S_\alpha}X_\alpha\to X_\alpha$ o é \sref{IV.1.6.2}. Basta, portanto, em virtude de \sref{IV.1.8.7.1}, aplicar o caso~(vi) ao morfismo diagonal
\[
  \Delta_{f_\alpha}:X_\alpha\longrightarrow X_\alpha\times_{S_\alpha}X_\alpha,
\]
observando que $\Delta_{f_\lambda}=\Delta_{f_\alpha}\times1_{S_\lambda}$ e $\Delta_f=\Delta_{f_\alpha}\times1_S$ \sref[I]{I.5.3.4} e \sref[I]{I.3.3.11}.

\emph{Casos} (viii) \emph{e} (ix): Como $S=\operatorname{Spec}(A)$ é afim, dizer que $f:X\to S$ é afim (resp.\ quase-afim) significa que existem um inteiro $r$ e uma imersão fechada (resp.\ uma imersão) $j:X\to\mathbf{V}^r_S=\operatorname{Spec}(A[T_1,\ldots,T_r])$ de $S$-pré-esquemas, pois $f$ é de tipo finito e $S$ é quase-compacto \sref[II]{II.5.1.9}.
Como $\mathbf{V}^r_S=\mathbf{V}^r_{S_0}\times_{S_0}S$, e $\mathbf{V}^r_{S_0}$ é um $S_0$-pré-esquema de apresentação finita, segue de \sref{IV.8.8.2},~(i), aplicado ao $S$-morfismo $j$, que existem $\lambda$ e um $S_\lambda$-morfismo $j_\lambda:X_\lambda\to\mathbf{V}^r_{S_\lambda}$ tais que $j=j_\lambda\times 1_S$; aplicando (iv) (resp.\ (ii)) a $j$, deduz-se que existe $\mu\geqslant\lambda$ tal que $j_\mu$ é uma imersão fechada (resp.\ uma imersão); portanto, $f_\mu$ é afim (resp.\ quase-afim).

\emph{Caso} (x): Por hipótese, $X=\operatorname{Spec}(B)$, onde $B$ é uma $A$-álgebra que, como $A$-módulo, é de apresentação finita \sref{IV.1.4.7}; segue, portanto, de \sref{IV.8.5.2},~(ii), que existem $\lambda$ e um $A_\lambda$-módulo de apresentação finita $B_\lambda$ tais que o $A$-módulo $B$ é isomorfo a $B_\lambda\otimes_{A_\lambda}A$.
A estrutura de $A$-álgebra de $B$ é definida por um $A$-homomorfismo $m:B\otimes_A B\to B$; como $B\otimes_A B=(B_\lambda\otimes_{A_\lambda}B_\lambda)\otimes_{A_\lambda} A$, existem, por \sref{IV.8.5.2},~(i), $\mu\geqslant\lambda$ e um $A_\mu$-homomorfismo $m_\mu:B_\mu\otimes_{A_\mu}B_\mu\to B_\mu$ tais que $m=m_\mu\otimes 1$.
Considerando os diagramas usuais que exprimem a associatividade e a comutatividade de $m$, vê-se, aplicando novamente \sref{IV.8.5.2},~(i), que podemos supor $\nu$ suficientemente grande para que $m_\nu$ defina sobre $B_\nu$ uma multiplicação associativa e comutativa.
Do mesmo modo, podemos supor $\nu$ suficientemente grande para que $B_\nu$, assim definido, admita um elemento unidade.
Se $X'_\nu=\operatorname{Spec}(B_\nu)$, fica claro então que $X$ é $S$-isomorfo a $X'_\nu\times_{S_\nu}S$; portanto, em virtude de (i), existe $\rho\geqslant\nu$ tal que $X_\rho$ e $X'_\rho$ são $S_\rho$-isomorfos, o que conclui a demonstração neste caso.

Para demonstrar o teorema no caso~(xii), demonstraremos primeiro a proposição seguinte:
\end{proof}

\begin{proposition}[8.10.5.1]
\label{IV.8.10.5.1}
(Lema de Chow para morfismos de apresentação finita). --- Seja $\mathrm{A}$ um anel, sejam $X$ e $Y$ dois $\mathrm{A}$-pré-esquemas de apresentação finita e seja $f:X\to Y$ um $\mathrm{A}$-morfismo separado. Então existem dois $\mathrm{A}$-pré-esquemas $X'$ e $P$ de apresentação finita, e $\mathrm{A}$-morfismos $p:P\to Y$, $j:X'\to P$ e $g:X'\to X$, tais que o diagrama
\oldpage[IV-3]{40}
\[
\begin{tikzcd}[column sep=large,row sep=large]
X' \arrow[r,"j"] \arrow[d,"g"'] & P \arrow[d,"p"] \\
X \arrow[r,"f"'] & Y
\end{tikzcd}
\]
é comutativo e tais que: $1^\circ$ $p$ é projetivo; $2^\circ$ $g$ é projetivo e sobrejetivo; $3^\circ$ $j$ é uma imersão aberta.
\end{proposition}

\begin{proof}
Escolhamos $\mathrm{A}_0\subset\mathrm{A}$, $X_0$, $Y_0$ e $f_0$ como em \sref{IV.8.9.1}, com $Y_0$ noetheriano e $f_0$ de tipo finito; por \sref{IV.8.10.4}, podemos supor, além disso, que $f_0$ é separado. O lema de Chow \sref[II]{II.5.6.1} fornece então morfismos de tipo finito $p_0:P_0\to Y_0$, $g_0:X'_0\to X_0$ e $j_0:X'_0\to P_0$ tais que
\[
\begin{tikzcd}[column sep=large,row sep=large]
X'_0 \arrow[r,"j_0"] \arrow[d,"g_0"'] & P_0 \arrow[d,"p_0"] \\
X_0 \arrow[r,"f_0"'] & Y_0
\end{tikzcd}
\]
é comutativo, $p_0$ é projetivo, $g_0$ é projetivo e sobrejetivo, e $j_0$ é uma imersão aberta. As asserções decorrem da invariância dessas propriedades por mudança de base \sref[II]{II.5.5.5},~(iii), \sref[I]{I.3.5.2} e \sref[I]{I.4.3.2}.
\end{proof}

\begin{proof}[Continuação da demonstração de \sref{IV.8.10.5}]
\emph{Caso (xii).} Apliquemos \sref{IV.8.10.5.1} a $f_0:X_0\to S_0$. Obtemos um diagrama comutativo
\[
\begin{tikzcd}[column sep=large,row sep=large]
X'_0 \arrow[r,"j_0"] \arrow[d,"g_0"'] & P_0 \arrow[d,"p_0"] \\
X_0 \arrow[r,"f_0"'] & S_0
\end{tikzcd}
\]
onde $p_0$ é projetivo, $g_0$ é projetivo e sobrejetivo, e $j_0$ é uma imersão aberta. Para cada $\lambda$, obtém-se o diagrama análogo, no qual $p_\lambda=p_0\times1_{S_\lambda}$, $g_\lambda=g_0\times1_{S_\lambda}$ e $j_\lambda=j_0\times1_{S_\lambda}$ possuem as mesmas propriedades respectivas, e o mesmo ocorre com os morfismos obtidos por passagem ao limite projetivo $p=p_0\times1_S$, $g=g_0\times1_S$ e $j=j_0\times1_S$. Como $g$ é próprio \sref[II]{II.5.5.3}, $p\circ j=f\circ g$ também o é \sref[II]{II.5.4.2}; e, como $p$ é separado, $j$ é próprio e, portanto, uma imersão fechada. Apliquemos o caso~(iv) a $j$, observando que $X'_0$ e $P_0$ são $S_0$-pré-esquemas de apresentação finita por \sref{IV.8.10.5.1} e \sref{IV.1.6.2}. Existe $\lambda$ tal que $j_\lambda$ é uma imersão fechada e, portanto, própria \sref[II]{II.5.4.2}. Assim,
\[
  f_\lambda\circ g_\lambda=p_\lambda\circ j_\lambda
\]
é próprio \sref[II]{II.5.5.3} e \sref[II]{II.5.4.2}. Como $g_\lambda$ é sobrejetivo e, pela hipótese sobre $f$ e pelo caso~(v), pode-se supor que $f_\lambda$ é separado, \sref[II]{II.5.4.3} demonstra que $f_\lambda$ é próprio.

\emph{Casos (xiii) e (xiv).} Pelo caso~(xii) e \sref[II]{II.5.5.3}, que se aplica porque os $S_\lambda$ são quase-compactos e quase-separados por \sref{IV.1.7.19}, basta considerar o caso em que $f$ é quase-projetivo. Suponhamos que exista um $\sh{O}_X$-Módulo invertível $\sh{L}$ que seja amplo relativamente a $f$. Como $S_0$ é quase-compacto e quase-separado, o mesmo vale para $X_0$ \sref{IV.1.2.3}; existem, portanto, $\lambda$ e um $\sh{O}_{X_\lambda}$-Módulo quase-coerente de apresentação finita $\sh{L}_\lambda$ tais que $\sh{L}=v_\lambda^*(\sh{L}_\lambda)$ \sref{IV.8.5.2},~(ii). Por \sref{IV.8.5.5}, podemos supor, além disso, que $\sh{L}_\lambda$ é invertível. A asserção decorre agora do lema mais preciso a seguir.

\begin{lemma}[8.10.5.2]
\label{IV.8.10.5.2}
\oldpage[IV-3]{41}
Suponhamos que $S_0$ seja quase-compacto, e seja $\sh{L}_\lambda$ um $\sh{O}_{X_\lambda}$-Módulo invertível. Para que $\sh{L}$ seja amplo relativamente a $f$ (resp.\ muito amplo relativamente a $f$), é necessário e suficiente que exista $\mu\geqslant\lambda$ tal que $\sh{L}_\mu$ seja amplo relativamente a $f_\mu$ (resp.\ muito amplo relativamente a $f_\mu$).
\end{lemma}

\begin{proof}
A condição é evidentemente suficiente \sref[II]{II.4.4.10} e \sref[II]{II.4.6.13}. Reciprocamente, como os $S_\lambda$ são quase-compactos e os $f_\lambda$ são de tipo finito, após substituir $\sh{L}$ por uma potência tensorial adequada, podemos nos reduzir ao caso em que $\sh{L}$ é muito amplo \sref[II]{II.4.6.11}. A questão é local sobre $S_0$ por \sref[II]{II.4.4.5} e pelo fato de que $L$ é filtrante, de modo que podemos supor que $S_0$ e, portanto, $S$ sejam afins. Por \sref[II]{II.4.4.1},~(ii) e \sref[II]{II.4.1.2}, existe uma $S$-imersão
\[
  j:X\longrightarrow\mathbf{P}^r_S=P
\]
tal que $\sh{L}\simeq j^*(\sh{O}_P(1))$. Por \sref{IV.8.8.2},~(i), o caso~(ii) e \sref[II]{II.4.1.3}, existem $\mu\geqslant\lambda$ e uma imersão $j_\mu:X_\mu\to\mathbf{P}^r_{S_\mu}=P_\mu$ tais que $j=j_\mu\times1_S$. Então, \sref[II]{II.4.1.3.2} e \sref{IV.8.5.2.5} fornecem $\nu\geqslant\mu$ tal que
\[
  \sh{L}_\nu\simeq j_\nu^*(\sh{O}_{P_\nu}(1)),
\]
o que demonstra que $\sh{L}_\nu$ é muito amplo relativamente a $f_\nu$ \sref[II]{II.4.4.2}.
\end{proof}
\end{proof}

\subsection{Aplicação aos morfismos quase-finitos}
\label{subsection:IV.8.11}
\label{IV.8.11}

Propomo-nos demonstrar neste parágrafo os dois teoremas seguintes:

\begin{theorem}[8.11.1]
\label{IV.8.11.1}
Seja $f:X\to Y$ um morfismo próprio, localmente de apresentação finita e quase-finito.
Então o morfismo $f$ é finito.
\end{theorem}

\begin{theorem}[8.11.2]
\label{IV.8.11.2}
Seja $f:X\to Y$ um morfismo localmente de apresentação finita, quase-finito e separado.
Então o morfismo $f$ é quase-afim e, \emph{a fortiori}, quase-projetivo.
\end{theorem}

\begin{remark}[8.11.3]
\label{IV.8.11.3}
Veremos adiante que, para a demonstração de \sref{IV.8.11.1} e \sref{IV.8.11.2}, podemos nos reduzir ao caso em que $Y$ é localmente noetheriano; observe-se que, nesse caso, obtemos assim outra demonstração do teorema de Chevalley \sref[III]{III.4.4.2}.
\end{remark}

\begin{env}[8.11.4]
\label{IV.8.11.4}
As hipóteses e as conclusões de \sref{IV.8.11.1} e \sref{IV.8.11.2} são todas locais sobre $Y$.
Com efeito, isto decorre de \sref{IV.1.6.1}, \sref{IV.1.2.6}, \sref[II]{II.5.1.1}, \sref[II]{II.5.4.1} e \sref[II]{II.6.2.2}.
Podemos, portanto, supor $Y=\operatorname{Spec}(A)$ afim.
Sabemos então que existem um subanel $A_0$ de $A$, que é uma $\bb{Z}$-álgebra de tipo finito, e um morfismo de tipo finito $f_0:X_0\to\operatorname{Spec}(A_0)$ tais que $X$ se identifica com $X_0\otimes_{A_0}A$ e $f$ com $f_0\times 1$ \sref{IV.8.9.1}.
Além disso, $A$ pode ser considerado como limite indutivo de seus subanéis que contêm $A_0$ e que são $\bb{Z}$-álgebras de tipo finito; utilizando o método de \sref{IV.8.1.2},~$c$)) e \sref{IV.8.10.5},~(v),~(xi) e (xii)), vê-se que basta demonstrar os teoremas \sref{IV.8.11.1} e \sref{IV.8.11.2} para $f_0$.
Suponhamos doravante $Y$ noetheriano; utilizando agora o método de \sref{IV.8.1.2},~$a$)) e \sref{IV.8.10.5},~(v),~(ix),~(x),~(xi) e (xii)), pode-se substituir $Y$ por $\operatorname{Spec}(\sh{O}_y)$, onde $y$ é um
\oldpage[IV-3]{42}
ponto arbitrário de $Y$, de modo que podemos finalmente supor que $Y=\operatorname{Spec}(A)$, onde $A$ é um \emph{anel local noetheriano}.
Seja $\mathfrak{m}$ o ideal maximal de $A$, e seja $\hat{A}$ o completamento de $A$ para a topologia $\mathfrak{m}$-ádica. O anel $\hat{A}$ é local noetheriano e $\operatorname{Spec}(\hat{A})\to\operatorname{Spec}(A)$ é fielmente plano e quase-compacto \sref[0]{0.7.3.5}; aplicando \sref{IV.2.7.1},~(i), (vii), (xiv), (xv) e (xvi), podemos, portanto, nos reduzir ao caso em que $A$ é \emph{completo}.
Resulta então de \sref[II]{II.6.2.6} que $X=X'\sqcup X''$, onde $X'$ é um $Y$-esquema finito e $X''$ um $Y$-esquema quase-finito tal que $X''\cap f^{-1}(y)=\varnothing$.
\end{env}

Coloquemo-nos primeiro nas hipóteses de \sref{IV.8.11.1}; como $f$ é próprio, $X''$, que é fechado em $X$, é próprio sobre $Y$ \sref[II]{II.5.4.10}, de modo que $f(X'')$ é fechado em $Y$; mas $y$ não pertence a $f(X'')$ e, além disso, $y$ é aderente a todo ponto de $Y$, logo $f(X'')=\varnothing$; por conseguinte, $X''$ é vazio e $f$ é \emph{finito}.

Coloquemo-nos agora nas hipóteses de \sref{IV.8.11.2} e, restringindo-nos ainda (como permite o que precede) ao caso em que $Y=\operatorname{Spec}(A)$ é afim e noetheriano de dimensão finita, raciocinemos também por indução sobre a dimensão de $Y$.
Reduzindo-nos como acima ao caso em que $A$ é, além disso, local e completo, temos $\dim(\sh{O}_y)=\dim(A)=\dim(Y)$ e, para todo $z\neq y$, $\dim(\sh{O}_z)<\dim(\sh{O}_y)$, de modo que $\dim(Y-\{y\})<\dim(Y)$.
Ora, por hipótese, temos $f(X'')\subset Y-\{y\}$, e a restrição de $f$ a $X''$ é evidentemente um morfismo quase-finito e separado; aplicando a hipótese de indução a $Y-\{y\}$ e a $X''$, vemos que $X''$ é quase-afim sobre $Y-\{y\}$; mas, como o aberto $Y-\{y\}$ é quase-afim sobre $Y$ porque $Y$ é noetheriano, $X''$ também é quase-afim sobre $Y$ \sref[II]{II.5.1.10},~(ii)); e, como além disso $X'$ é finito (e \emph{a fortiori} afim) sobre $Y$, $X$ é quase-afim sobre $Y$ \sref[II]{II.4.6.17} e \sref[II]{II.5.1.2},~$c'$)).

\begin{proposition}[8.11.5]
\label{IV.8.11.5}
Seja $f:X\to Y$ um morfismo de apresentação finita.
As propriedades seguintes são equivalentes:
\begin{enumerate}[label=\emph{\alph*)}]
  \item $f$ é uma imersão fechada.
  \item $f$ é um monomorfismo próprio.
  \item $f$ é próprio e, para todo $y\in Y$, $f^{-1}(y)$ é radicial e geometricamente reduzido sobre $k(y)$ (isto é, vazio ou $k(y)$-isomorfo a $\operatorname{Spec}(k(y))$).
\end{enumerate}
\end{proposition}

É claro que \emph{a)} implica \emph{b)}.
Para ver que \emph{b)} implica \emph{c)}, observemos \sref[I]{I.3.3.12} que, para todo $y\in Y$, o morfismo $f^{-1}(y)\to\operatorname{Spec}(k(y))$ deduzido de $f$ por mudança de base é um monomorfismo e, portanto, injetivo; por conseguinte, $f^{-1}(y)$ é vazio ou se reduz a um ponto e, em qualquer caso, é afim.
Além disso, se $A$ é o anel de $f^{-1}(y)$, o homomorfismo canônico $A\otimes_k A\to A$ é bijetivo \sref[I]{I.5.3.8}.
Isso implica evidentemente que $A=k$, pois, caso contrário, haveria um elemento $a\in A$ que não pertenceria a $k$, e as imagens de $a\otimes 1$ e de $1\otimes a$ em $A$ seriam ambas iguais a $a$, ao passo que $a\otimes 1\neq1\otimes a$, pois $1$ e $a$ são linearmente independentes sobre $k$.

Resta demonstrar que \emph{c)} implica \emph{a)}.
Decorre, em primeiro lugar, de \sref{IV.8.11.1} que $f$ é um morfismo \emph{finito}, de modo que $X=\operatorname{Spec}(\sh{A})$, onde $\sh{A}$ é uma $\sh{O}_Y$-álgebra finita.
Basta, portanto, demonstrar que o homomorfismo canônico $\sh{O}_Y\to\sh{A}$ é \emph{sobrejetivo} \sref[II]{II.1.4.10}, ou, equivalentemente, que, para todo $y\in Y$, o homomorfismo $\sh{O}_y\to\sh{A}_y$ é sobrejetivo.
Mas, por hipótese,
\oldpage[IV-3]{43}
$f^{-1}(y)=\operatorname{Spec}(\sh{A}_y/\mathfrak{m}_y\sh{A}_y)$ \sref[II]{II.1.5.5} é tal que o homomorfismo correspondente $k(y)=\sh{O}_y/\mathfrak{m}_y\to\sh{A}_y/\mathfrak{m}_y\sh{A}_y$ é sobrejetivo; como $\sh{A}_y$ é um $\sh{O}_y$-módulo de tipo finito, o lema de Nakayama mostra que $\sh{O}_y\to\sh{A}_y$ é sobrejetivo, o que conclui a demonstração.

\begin{remark}[8.11.5.1]
\label{IV.8.11.5.1}
Observe-se que o raciocínio anterior demonstra que, se $f$ é um \emph{monomorfismo}, então, para todo $y\in Y$, $f^{-1}(y)$ é vazio ou $k(y)$-isomorfo a $\operatorname{Spec}(k(y))$.
\end{remark}

\begin{proposition}[8.11.6]
\label{IV.8.11.6}
Se um morfismo $f:X\to Y$ de apresentação finita é um homeomorfismo universal, então ele é finito, sobrejetivo e radicial (sendo a recíproca verdadeira por \sref{IV.2.4.5},~(i)).
\end{proposition}

\begin{proof}
Com efeito, sendo $f$ de tipo finito, universalmente fechado e separado em virtude de \sref{IV.2.4.4}, ele é próprio por definição \sref[II]{II.5.4.1}.
Como é evidentemente quase-finito \sref[II]{II.6.2.3}, ele é finito por \sref{IV.8.11.1}.
Sabemos, além disso, que ele é radicial \sref{IV.2.4.4} e evidentemente sobrejetivo.
\end{proof}

\subsection{Nova demonstração e generalização do ``Teorema principal'' de Zariski}
\label{subsection:IV.8.12}
\label{IV.8.12}

\begin{lemma}[8.12.1]
\label{IV.8.12.1}
Seja $f:X\to Y$ um morfismo quase-compacto e quase-separado, seja $\sh{C}$ uma $\sh{O}_Y$-Álgebra quase-coerente e seja $Z=\operatorname{Spec}(\sh{C})$, que é um $Y$-pré-esquema afim sobre $Y$.
Seja $g:X\to Z$ um $Y$-morfismo e seja $\varphi=\sh{A}(g):\sh{C}\to f_*(\sh{O}_X)$ o $\sh{O}_Y$-homomorfismo correspondente de $\sh{O}_Y$-Álgebras \sref[II]{II.1.2.7}.
Suponhamos que $g$ seja uma imersão.
Então, para que a imagem fechada de $X$ por $g$ \sref[I]{I.9.5.3} seja igual a $Z$, é necessário e suficiente que $\varphi$ seja injetivo; $g$ é então uma imersão aberta.
\end{lemma}

\begin{proof}
A hipótese implica que $f_*(\sh{O}_X)$ é uma $\sh{O}_Y$-Álgebra quase-coerente \sref{IV.1.7.5}; além disso, como o morfismo canônico $h:Z\to Y$ é afim, logo quase-compacto e separado, $g$ é um morfismo quase-compacto e quase-separado \sref{IV.1.2.2} e \sref{IV.1.1.2}, logo $g_*(\sh{O}_X)$ é uma $\sh{O}_Z$-Álgebra quase-coerente \sref{IV.1.7.5}.
Nessas condições, dizer que a imagem fechada de $X$ por $g$ é igual a $Z$ significa \sref[I]{I.9.5.2} que o homomorfismo canônico $\theta:\sh{O}_Z\to g_*(\sh{O}_X)$ é injetivo.
Mas temos $h_*(\sh{O}_Z)=\sh{C}$ por definição de $Z$ \sref[II]{II.1.3.1}, e $h_*(g_*(\sh{O}_X))=f_*(\sh{O}_X)$.
Como $Z$ é afim sobre $Y$, é equivalente dizer que o homomorfismo $\theta:\sh{O}_Z\to g_*(\sh{O}_X)$ é injetivo ou que o homomorfismo correspondente $\varphi=h_*(\theta):\sh{C}\to f_*(\sh{O}_X)$ é injetivo \sref[I]{I.1.3.9}.
O fato de $g$ ser então uma imersão aberta decorre de \sref[I]{I.9.5.10} e da hipótese.
\end{proof}

\begin{lemma}[8.12.2]
\label{IV.8.12.2}
Seja $Y$ um pré-esquema quase-compacto e quase-separado, seja $f:X\to Y$ um morfismo quase-separado de tipo finito, seja $\sh{C}$ uma $\sh{O}_Y$-Álgebra quase-coerente e seja $Z=\operatorname{Spec}(\sh{C})$.
Seja $g:X\to Z$ um $Y$-morfismo e seja $\varphi:\sh{C}\to f_*(\sh{O}_X)$ o $\sh{O}_Y$-homomorfismo correspondente de $\sh{O}_Y$-Álgebras.
Seja $(\sh{C}_\lambda)$ a família crescente filtrante das sub-$\sh{O}_Y$-Álgebras quase-coerentes de tipo finito de $\sh{C}$ (cuja união é $\sh{C}$ \sref[I]{I.9.6.6} e \sref{IV.1.7.9}); ponhamos $Z_\lambda=\operatorname{Spec}(\sh{C}_\lambda)$ e seja $g_\lambda$ o morfismo composto $X\xrightarrow{g}Z\to Z_\lambda$.
Então as condições seguintes são equivalentes:
\begin{enumerate}[label=\emph{\alph*)}]
  \item $g$ é uma imersão.
  \item Existe $\lambda$ tal que $g_\lambda$ é uma imersão.
\end{enumerate}
\oldpage[IV-3]{44}
Além disso, quando $g_\lambda$ é uma imersão, o mesmo ocorre com $g_\mu$ para $\mu\geqslant\lambda$.

Basta aplicar \sref[II]{II.3.8.4}, substituindo $\sh{L}$ por $\sh{O}_X$ e $\sh{S}$ por $\sh{C}[\mathrm{T}]$, e levando em conta \sref[II]{II.3.1.7}.
\end{lemma}

\begin{proposition}[8.12.3]
\label{IV.8.12.3}
Seja $Y$ um pré-esquema quase-compacto e quase-separado, e seja $f:X\to Y$ um morfismo separado de tipo finito.
Seja $\sh{B}=f_*(\sh{O}_X)$, que é uma $\sh{O}_Y$-Álgebra quase-coerente \sref[I]{I.9.2.2}; seja $\sh{C}$ o fecho integral de $\sh{O}_Y$ em $\sh{B}$, que é uma $\sh{O}_Y$-Álgebra quase-coerente \sref[II]{II.6.3.4}; ponhamos $Z=\operatorname{Spec}(\sh{C})$ e seja $g:X\to Z$ o $Y$-morfismo correspondente à injeção canônica $\varphi:\sh{C}\to\sh{B}=f_*(\sh{O}_X)$ \sref[II]{II.1.2.7}. Seja $(\sh{C}_\lambda)$ a família filtrante crescente das sub-$\sh{O}_Y$-Álgebras quase-coerentes de tipo finito de $\sh{C}$ (cuja reunião é $\sh{C}$ \sref[I]{I.9.6.6} e \sref{IV.1.7.9}) e, para todo $\lambda$, ponhamos $Z_\lambda=\operatorname{Spec}(\sh{C}_\lambda)$ e seja $g_\lambda:X\to Z_\lambda$ o $Y$-morfismo correspondente à injeção $\sh{C}_\lambda\to\sh{B}=f_*(\sh{O}_X)$.
Então as condições seguintes são equivalentes:
\begin{enumerate}[label=\emph{\alph*)}]
  \item[\emph{a)}] Existe uma fatoração de $f$ da forma
\[
X\xrightarrow{f'}Y'\xrightarrow{u}Y
\]
onde $f'$ é uma imersão e $u$ é um morfismo finito.
  \item[\emph{a$'$)}] Existe uma fatoração $X\xrightarrow{f'}Y'\xrightarrow{u}Y$ de $f$, onde $f'$ é uma imersão aberta e $u$ é um morfismo finito.
  \item[\emph{b)}] O morfismo $g:X\to Z$ é uma imersão.
  \item[\emph{c)}] Existe $\lambda$ tal que $g_\lambda:X\to Z_\lambda$ é uma imersão.
\end{enumerate}
Além disso, quando $g_\lambda$ é uma imersão, $g$ é uma imersão aberta, $g(X)$ é denso em $Z$, e existe $\lambda$ tal que, para $\mu\geqslant\lambda$, $g_\mu$ é uma imersão aberta.
\end{proposition}

Como o homomorfismo $\varphi:\sh{C}\to f_*(\sh{O}_X)$ é injetivo, decorre de \sref{IV.8.12.1} que, se $g$ é uma imersão, ele é uma imersão aberta e $g(X)$ é denso em $Z$; o mesmo vale para $g_\lambda$.
O fato de \emph{a)} implicar \emph{a$'$)} decorre também de \sref{IV.8.12.1}, aplicado substituindo $Z$ por $Y'$ e $g$ por $f'$ (pois $Y'$ é finito sobre $Y$, portanto afim sobre $Y$): com efeito, se $Y''$ é a imagem fechada de $X$ por $f'$, $Y''$ é finito sobre $Y$ e $f'$ se fatora como $X\xrightarrow{f''}Y''\xrightarrow{j}Y'$, onde $j$ é a injeção canônica e $f''$ é uma imersão \sref[I]{I.4.1.10}; decorre, portanto, de \sref{IV.8.12.1} que $f''$ é uma imersão aberta.

A equivalência de \emph{b)} e \emph{c)} decorre de \sref{IV.8.12.2}, assim como o fato de que $g_\lambda$ é então uma imersão para $\lambda$ suficientemente grande.
É claro que \emph{c)} implica \emph{a)}, pois $Z_\lambda$ é finito sobre $Y$ \sref[II]{II.6.3.4} e \sref[II]{II.6.1.2}.
Demonstremos, por fim, que \emph{a)} implica \emph{c)}.
Vimos acima que podemos supor que $Y'$ seja a imagem fechada de $X$ por $f'$ e decorre então de \sref{IV.8.12.1} que, pondo $\sh{B}'=u_*(\sh{O}_{Y'})$, de modo que $Y'$ se identifica com $\operatorname{Spec}(\sh{B}')$, o homomorfismo $\varphi':\sh{B}'\to\sh{B}=f_*(\sh{O}_X)$ é \emph{injetivo}.
Mas, como por hipótese $\sh{B}'$ é uma $\sh{O}_Y$-álgebra finita, ela se identifica, pela definição de $\sh{C}$, com uma das $\sh{O}_Y$-subálgebras $\sh{C}_\lambda$, o que demonstra \emph{c)}.

Diremos que um morfismo $f:X\to Y$, onde $Y$ é quase-compacto e quase-separado, é \emph{pseudofinito} se é de tipo finito e satisfaz a condição \emph{a)} de \sref{IV.8.12.3} (caso em que é necessariamente separado).

\begin{corollary}[8.12.4]
\label{IV.8.12.4}
Seja $Y$ um pré-esquema quase-compacto e quase-separado, e seja $f:X\to Y$ um morfismo.
\oldpage[IV-3]{45}
\begin{enumerate}[label=(\roman*)]
  \item Suponhamos que $f$ seja pseudofinito.
Então, para todo morfismo $Y'\to Y$, onde $Y'$ é quase-compacto e quase-separado, $f_{(Y')}:X'=X_{(Y')}\to Y'$ é pseudofinito.
  \item Seja $(U_\lambda)$ uma cobertura de $Y$ por abertos quase-compactos.
Para que $f$ seja pseudofinito, é necessário e suficiente que, para todo $\lambda$, a restrição $f_\lambda:f^{-1}(U_\lambda)\to U_\lambda$ de $f$ seja um morfismo pseudofinito.
  \item Suponhamos, além disso, que $Y$ seja noetheriano e $f$ de tipo finito.
Então, para que $f$ seja pseudofinito, é necessário e suficiente que, para todo $y\in Y$, o morfismo $f_y=f\times 1:X_y=X\times_Y\operatorname{Spec}(\sh{O}_y)\to\operatorname{Spec}(\sh{O}_y)$ seja pseudofinito.
\end{enumerate}
\end{corollary}

\begin{proof}
(i) É claro que $f_{(Y')}$ é de tipo finito \sref{IV.1.5.4}. Além disso, uma fatoração $X\xrightarrow{g}Z\xrightarrow{u}Y$, onde $g$ é uma imersão e $u$ é finito, induz uma fatoração $X'\xrightarrow{g'}Z'\xrightarrow{u'}Y'$ de $f_{(Y')}$, onde $Z'=Z_{(Y')}$, $g'=g_{(Y')}$ e $u'=u_{(Y')}$. O morfismo $g'$ é uma imersão \sref[I]{I.4.3.2} e $u'$ é finito \sref[II]{II.6.1.5}; portanto, $f_{(Y')}$ é pseudofinito.

(ii) A condição é necessária em virtude de (i), pois os $U_\lambda$ são quase-separados, uma vez que $Y$ o é.
Para ver que ela é suficiente, observemos (com as notações de \sref{IV.8.12.3}) que, se pusermos $X_\lambda=f^{-1}(U_\lambda)$, teremos $\sh{B}|U_\lambda=(f_\lambda)_*(\sh{O}_{X_\lambda})$, e $\sh{C}|U_\lambda$ será o fecho integral de $\sh{O}_{U_\lambda}$ em $\sh{B}|U_\lambda$. Por conseguinte, se $h:Z\to Y$ for o morfismo canônico, $Z_\lambda=\operatorname{Spec}(\sh{C}|U_\lambda)$ se identificará com $h^{-1}(U_\lambda)$. Para que $g:X\to Z$ seja uma imersão, é necessário e suficiente que, para todo $\lambda$, sua restrição $g_\lambda:f^{-1}(U_\lambda)\to h^{-1}(U_\lambda)$ seja uma imersão \sref[I]{I.4.2.4}. A conclusão decorre agora de \sref{IV.8.12.3}.

(iii) Basta, em virtude de (ii), demonstrar que $y$ admite uma vizinhança $U$ tal que a restrição $f^{-1}(U)\to U$ de $f$ seja um morfismo pseudofinito.
Designemos por $(U_\lambda)$ o sistema projetivo filtrante decrescente das vizinhanças abertas afins de $y$ e apliquemos o método de \sref{IV.8.1.2},~$a$)).
Como $Y$ é noetheriano, as restrições $f_\lambda:f^{-1}(U_\lambda)\to U_\lambda$ de $f$ são de apresentação finita, e o mesmo vale para $f_y$.
Por hipótese, $f_y$ se fatora como $X_y\xrightarrow{g_y}Z_y\xrightarrow{u_y}\operatorname{Spec}(\sh{O}_y)$, onde $u_y$ é finito e $g_y$ é uma imersão.
Como $\sh{O}_y$ é noetheriano, o mesmo vale para $Z_y$; e, como $Z_y$ é de apresentação finita sobre $\operatorname{Spec}(\sh{O}_y)$, existem $\lambda$ e um morfismo de apresentação finita $u_\lambda:Z_\lambda\to U_\lambda$ tais que $Z_y$ se identifica com $Z_\lambda\times_{U_\lambda}\operatorname{Spec}(\sh{O}_y)$ e $u_y$ com $u_\lambda\times 1$ \sref{IV.8.8.2},~(ii)). Além disso, existe um morfismo $g_\lambda:X_\lambda\to Z_\lambda$ tal que $g_y=g_\lambda\times 1$ e $f_\lambda=u_\lambda\circ g_\lambda$ \sref{IV.8.8.2},~(i)). Pode-se, ademais, tomar $\lambda$ de modo que $g_\lambda$ seja uma imersão e $u_\lambda$ um morfismo finito \sref{IV.8.10.5},~(ii) e (x)), o que demonstra que $f_\lambda$ é pseudofinito.
\end{proof}

\begin{env}[8.12.5]
\label{IV.8.12.5}
Podemos agora dar uma demonstração do ``Teorema principal'' de Zariski \sref[III]{III.4.4.3}) que não utiliza os resultados cohomológicos de natureza ``global'' do cap.~III, mas recorre, em contrapartida, às propriedades mais finas dos anéis locais noetherianos; além disso, generalizaremos o enunciado do teorema eliminando as hipóteses noetherianas:
\end{env}

\begin{theorem}[8.12.6]
\label{IV.8.12.6}
(``Teorema principal'' de Zariski). --- Seja $Y$ um pré-esquema quase-compacto e quase-separado.
Se um morfismo $f:X\to Y$ é quase-finito, separado e de apresentação finita, existe uma fatoração de $f$
\begin{equation}
\label{IV.8.12.6.1}
X\xrightarrow{f'}Y'\xrightarrow{u}Y \tag{8.12.6.1}
\end{equation}
onde $f'$ é uma imersão aberta e $u$ é um morfismo finito.
\end{theorem}

\oldpage[IV-3]{46}

\begin{proof}
Em virtude de \sref{IV.8.12.4},~(ii) e do caráter local (sobre $Y$) das noções de morfismo quase-finito, de morfismo separado e de morfismo de apresentação finita, podemos reduzir-nos ao caso em que $Y=\operatorname{Spec}(A)$ é afim.
Aplicando \sref{IV.8.9.1}, podemos supor que existem um subanel $A_0$ de $A$, que é uma $\bb{Z}$-álgebra de tipo finito, e um $A$-isomorfismo $X_0\otimes_{A_0}A\xrightarrow{\sim}X$, identificando-se $f$, por meio desse isomorfismo, com $f_0\times 1$, onde $f_0:X_0\to\operatorname{Spec}(A_0)$ é um morfismo \emph{de tipo finito}; além disso, por \sref{IV.8.10.5},~(v) e (xi), podemos supor que $f_0$ seja separado e quase-finito; se demonstrarmos que $f_0$ é pseudofinito, o mesmo valerá para $f$ por \sref{IV.8.12.4},~(i).
Como $A_0$ é então noetheriano e as noções de morfismo de tipo finito, separado e quase-finito são preservadas por mudança de base, decorre de \sref{IV.8.12.4},~(iii) que podemos até supor que $A$ seja um anel local, \emph{essencialmente de tipo finito sobre} $\bb{Z}$ \sref{IV.1.3.8}.
Ponhamos $n=\dim(A)$ e procedamos por indução sobre $n$; para $n=0$, o teorema é evidente, pois $A$ é um corpo e o morfismo $f$ já é finito \sref[II]{II.6.2.2}.
Ponhamos $B=\Gamma(X,\sh{O}_X)$; denotemos por $C$ o fecho integral de $A$ em $B$, ponhamos $Z=\operatorname{Spec}(C)$ e seja $g:X\to Z$ o $Y$-morfismo correspondente ao $A$-homomorfismo canônico injetivo $C\to B$; em virtude de \sref{IV.8.12.3}, basta demonstrar que $g$ é uma imersão aberta.
Seja $a$ o ponto fechado de $Y$, e seja $U=Y-\{a\}$; $U$ é noetheriano, e todos os seus anéis locais são essencialmente de tipo finito sobre $\bb{Z}$ e de dimensão $<n$; tendo em conta a hipótese de indução e \sref{IV.8.12.4},~(iii), vemos que a restrição $f^{-1}(U)\to U$ de $f$ é um morfismo pseudofinito.
Conclui-se de \sref{IV.8.12.3} que, se $h:Z\to Y$ é o morfismo estrutural, a restrição $f^{-1}(U)\to h^{-1}(U)$ de $g$ é uma imersão aberta.
Ponhamos $A'=\hat{A}$, $Y'=\operatorname{Spec}(A')$, $X'=X_{(Y')}$ e $f'=f_{(Y')}:X'\to Y'$.
Como o morfismo canônico $u:Y'\to Y$ é plano, decorre de \sref{IV.2.3.1} que $B'=\Gamma(X',\sh{O}_{X'})$ se identifica com a $A'$-álgebra $B\otimes_A A'$.
Por outro lado, como $A$ é um anel local excelente \sref{IV.7.8.3}, o morfismo $u:Y'\to Y$ é regular e, \emph{a fortiori}, normal; por conseguinte, \sref{IV.6.14.4}, \emph{o fecho integral de $A'$ em $B'$ é igual a} $C'=C\otimes_A A'$.
Vê-se, portanto, que $Z'=\operatorname{Spec}(C')$ é igual a $Z_{(Y')}$ e que o morfismo $g':X'\to Z'$, proveniente da injeção canônica $C'\to B'$, é igual a $g_{(Y')}$.
Como $u:Y'\to Y$ é fielmente plano e quase-compacto, para demonstrar que $g$ é uma imersão aberta basta demonstrar que $g'$ é uma imersão aberta \sref{IV.2.7.1},~(x).
Ora, $u^{-1}(a)$ reduz-se ao ponto fechado $a'$ de $Y'$ e, por conseguinte, $U'=Y'-\{a'\}=u^{-1}(U)$.
Se $h':Z'\to Y'$ é o morfismo canônico, o fato de a restrição $f^{-1}(U)\to h^{-1}(U)$ de $g$ ser uma imersão aberta acarreta que o mesmo se verifica para a restrição $f'^{-1}(U')\to h'^{-1}(U')$ de $g'$.
O morfismo $f'$ é separado e quase-finito \sref[II]{II.6.2.4}; como $A'$ é \emph{completo}, deduz-se de \sref[II]{II.6.2.6} que $X'$ é $Y'$-isomorfo a uma soma $X'_1\sqcup X'_2$, em que a restrição $f'|X'_1=f'_1:X'_1\to Y'$ é um morfismo \emph{finito}, e $X'_2\subset f'^{-1}(U')$.
Daí resulta que $B'$ é o produto direto das duas $A'$-álgebras $\Gamma(X'_1,\sh{O}_{X'_1})=B'_1$ e $\Gamma(X'_2,\sh{O}_{X'_2})=B'_2$; conclui-se imediatamente que o fecho integral $C'$ de $A'$ em $B'$ é o produto direto dos fechos integrais $C'_1$, $C'_2$ de $A'$ em $B'_1$, $B'_2$, respectivamente, donde $Z'=Z'_1\sqcup Z'_2$, em que $Z'_i=\operatorname{Spec}(C'_i)$ ($i=1,2$); e o morfismo canônico $g':X'\to Z'$ é tal que $g'|X'_i$ é o morfismo canônico $g'_i:X'_i\to Z'_i$ ($i=1,2$).
Mas, como $B'_1$ já é uma $A'$-álgebra finita, tem-se $C'_1=B'_1$, e portanto $g'_1$ é um isomorfismo.
Por outro lado, como $X'_2\subset f'^{-1}(U')$ e é aberto em $f'^{-1}(U')$, já
\oldpage[IV-3]{47}
se sabe que $g'_2$ é uma imersão aberta.
Conclui-se, portanto, que $g'$ é de fato uma imersão aberta.
\end{proof}

\begin{remark}[8.12.7]
\label{IV.8.12.7}
Quando, em \sref{IV.8.12.6}, se supõe que $Y$ seja um \emph{esquema afim}, a demonstração por redução ao caso noetheriano mostra que, na fatoração \sref{IV.8.12.6.1}, os morfismos $f'$ e $u$ são também morfismos \emph{de apresentação finita} \sref{IV.1.6.2}.
\end{remark}

\begin{corollary}[8.12.8]
\label{IV.8.12.8}
Seja $Y$ um esquema quase-compacto tal que exista um $\sh{O}_Y$-Módulo amplo \sref[II]{II.4.5.3}, e seja $f:X\to Y$ um morfismo quase-finito e quase-projetivo.
Então existe uma fatoração de $f$ da forma
\[
X\xrightarrow{f'}Y'\xrightarrow{u}Y
\]
em que $f'$ é uma imersão aberta e $u$ um morfismo finito.
\end{corollary}

\begin{proof}
A hipótese implica que $X$ se identifica com um sub-$Y$-esquema quase-compacto de um $Y$-esquema da forma $Z=\mathbf{P}^r_Y$ \sref[II]{II.5.3.3}.
Existe, portanto, uma vizinhança aberta quase-compacta $U$ de $X$ em $Z$ tal que $X$ seja fechado em $U$; como $Z$ é um esquema, a injeção canônica $U\to Z$ é um morfismo de apresentação finita (\sref{IV.1.2.7} e \sref{IV.1.6.2}), e portanto o morfismo composto $g:U\to Z\to Y$ é também um morfismo de apresentação finita (o fato de $\mathbf{P}^r_Y$ ser de apresentação finita sobre $Y$ decorre imediatamente da definição \sref[II]{II.4.1.1}).
Seja $\sh{J}$ o Ideal quase-coerente de $\sh{O}_U$ que define o subpré-esquema fechado $X$; como $U$ é um esquema quase-compacto, $\sh{J}$ é o limite indutivo filtrante de seus subideais quase-coerentes de tipo finito $\sh{J}_\lambda$ \sref[I]{I.9.4.9}.
Se $X_\lambda$ é o subpré-esquema fechado de $U$ definido por $\sh{J}_\lambda$, tem-se, portanto, $X=\bigcap_\lambda X_\lambda$.
Para todo $y\in Y$, tem-se, portanto, $f^{-1}(y)=\bigcap_\lambda(X_\lambda\cap g^{-1}(y))$ e, como os conjuntos $X_\lambda\cap g^{-1}(y)$ são fechados no espaço noetheriano $g^{-1}(y)$, existe, para todo $y$, um índice $\lambda(y)$ tal que $f^{-1}(y)=X_{\lambda(y)}\cap g^{-1}(y)$.
Denotemos por $E_\lambda$ o conjunto dos $y\in Y$ tais que a fibra $X_\lambda\cap g^{-1}(y)$ da restrição de $g$ a $X_\lambda$ seja um $k(y)$-pré-esquema finito.
A hipótese de que $f$ é quase-finito implica, em virtude do que precede, que $Y=\bigcup_\lambda E_\lambda$.
Ora, cada um dos $X_\lambda$ é, por definição, \emph{de apresentação finita} sobre $Y$; segue-se, portanto, de \sref{IV.9.2.3} e \sref{IV.9.2.6}~($^1$) que os $E_\lambda$ são conjuntos \emph{construtíveis} no esquema $Y$; como formam uma família filtrante crescente, existe um índice $\lambda$ tal que $E_\lambda=Y$ \sref{IV.1.9.9}, e, para esse $\lambda$, o morfismo $f_\lambda:X_\lambda\to Y$, restrição de $g$ a $X_\lambda$, é quase-finito. Como é de apresentação finita e separado, pode-se aplicar-lhe \sref{IV.8.12.6}, e $f_\lambda$ se fatora como
\[
X_\lambda\xrightarrow{j_\lambda}Y'_\lambda\xrightarrow{u_\lambda}Y
\]
em que $j_\lambda$ é uma imersão aberta e $u_\lambda$ um morfismo finito.
Como $X$ é um subpré-esquema fechado de $X_\lambda$, demonstramos assim que $f$ possui a propriedade \sref{IV.8.12.3}[$a$], donde resulta o corolário em virtude da equivalência entre \sref{IV.8.12.3}[$a$] e \sref{IV.8.12.3}[$a'$].
\end{proof}

\footnotetext{($^1$) O leitor verificará que os corolários \sref{IV.8.12.8} a \sref{IV.8.12.11} não são utilizados na \textsection~9.}

\begin{corollary}[8.12.9]
\label{IV.8.12.9}
\oldpage[IV-3]{48}
Seja $f:X\to Y$ um morfismo localmente quase-finito $(\mathbf{Err}_{\text{III}},~20)$.
Para todo $x\in X$, existem uma vizinhança aberta $U$ de $x$ em $X$ e uma vizinhança aberta $V$ de $y=f(x)$ em $Y$, tais que $f(U)\subset V$, e uma fatoração
\[
U\xrightarrow{f'}V'\xrightarrow{u}V
\]
da restrição de $f$ a $U$, em que $f'$ é uma imersão aberta e $u$ um morfismo finito.
\end{corollary}

\begin{proof}
Basta evidentemente tomar para $V$ uma vizinhança afim de $y$ em $Y$, e para $U$ uma vizinhança afim de $x$ em $X$, contida em $f^{-1}(V)$ e tal que $f|U$ seja quase-finito.
Como o morfismo $U\to V$, restrição de $f$, é então afim (e, portanto, quase-projetivo), pode-se aplicar \sref{IV.8.12.8}.
\end{proof}

\begin{corollary}[8.12.10]
\label{IV.8.12.10}
Sejam $Y$ um pré-esquema integral e normal, $X$ um pré-esquema integral e $f:X\to Y$ um morfismo birracional e localmente quase-finito $(\mathbf{Err}_{\text{III}},~20)$.
Então $f$ é um isomorfismo local; para que $f$ seja uma imersão aberta, é necessário e suficiente que $f$ seja, além disso, separado.
\end{corollary}

\begin{proof}
A segunda afirmação decorre imediatamente da primeira e de \sref[I]{I.8.2.8}.
Para demonstrar a primeira afirmação, podemos supor $X$ e $Y$ afins e $f$ quase-finito; consideremos a fatoração $f=u\circ f'$ de \sref{IV.8.12.8}, que identifica $X$, por meio de $f'$, com o subpré-esquema induzido sobre um aberto de $Y'$.
Como $X$ é integral, podemos, em virtude de \sref[I]{I.5.2.3}, substituir $Y'$ pelo subpré-esquema reduzido de $Y'$ cujo espaço subjacente é $\overline{X}$; podemos, portanto, supor também que $Y'$ é \emph{integral}.
Além disso, como $f$ é birracional, $u$ também o é.
A conclusão decorre então do lema seguinte:
\end{proof}

\begin{lemma}[8.12.10.1]
\label{IV.8.12.10.1}
Sejam $Y'$ um pré-esquema integral e $Y$ um pré-esquema integral e normal; então todo morfismo $u:Y'\to Y$ finito e birracional é um isomorfismo.
\end{lemma}

\begin{proof}
Ponhamos $\sh{A}=u_*(\sh{O}_{Y'})$; então $\sh{A}$ é uma $\sh{O}_Y$-álgebra finita e $Y'$ se identifica com $\operatorname{Spec}(\sh{A})$ \sref[II]{II.1.3.6}.
Se $R(Y)$ é o corpo de funções racionais de $Y$, tem-se, portanto, para todo $y\in Y$, $\sh{O}_{Y,y}\subset\sh{A}_y\subset R(Y)$; mas, como o anel $\sh{O}_{Y,y}$ é, por hipótese, integralmente fechado e tem $R(Y)$ como corpo de frações, tem-se necessariamente $\sh{A}_y=\sh{O}_{Y,y}$, donde o lema.
\end{proof}

\begin{corollary}[8.12.11]
\label{IV.8.12.11}
Sejam $Y$ um pré-esquema integral, $X$ um pré-esquema integral e normal, e $f:X\to Y$ um morfismo dominante e localmente quase-finito.
Sejam $K$ e $L$ (uma extensão de $K$) os corpos de funções racionais de $Y$ e $X$, respectivamente; e seja $Y'$ o fecho integral de $Y$ relativamente a $L$ \sref[II]{II.6.3.4}; então $f$ se fatora de uma única maneira como $f:X\xrightarrow{f'}Y'\xrightarrow{u}Y$, onde $f'$ é birracional e corresponde ao automorfismo identidade de $L$; $f'$ é então um isomorfismo local e, para que $f'$ seja uma imersão aberta, é necessário e suficiente que $f$ seja separado.
\end{corollary}

\begin{proof}
A existência e a unicidade da fatoração de $f$ decorrem de \sref[II]{II.6.3.9}.
Além disso, decorre de \sref[II]{II.6.2.4},~(v), reduzindo-nos ao caso afim, que $f'$ é localmente quase-finito; por outro lado, decorre de \sref[I]{I.5.5.1} que, para que $f'$ seja separado, é necessário e suficiente que $f$ o seja, pois $u$ é afim e, portanto, separado; as duas últimas afirmações são, pois, consequências de \sref{IV.8.12.10} aplicado a $f'$.
\end{proof}

\oldpage[IV-3]{49}

\subsection{Tradução em termos de pro-objetos}
\label{subsection:IV.8.13}

A proposição seguinte é essencialmente equivalente a \sref{IV.8.8.2},~(i):

\begin{proposition}[8.13.1]
\label{IV.8.13.1}
Sejam $S$ um pré-esquema e $(X_\lambda,v_{\lambda\mu})$ um sistema projetivo filtrante de $S$-pré-esquemas; suponhamos que exista $\alpha$ tal que $v_{\alpha\lambda}$ seja um morfismo afim para todo $\lambda\geqslant\alpha$ (o que implica, por \sref[II]{II.1.6.2}, que $v_{\lambda\mu}$ seja afim para $\alpha\leqslant\lambda\leqslant\mu$), de modo que o limite projetivo $X=\varprojlim_\lambda X_\lambda$ exista na categoria dos $S$-pré-esquemas \sref{IV.8.2.3}.
Seja $Y$ um $S$-pré-esquema e, para todo $\lambda\geqslant\alpha$, seja $e_\lambda:\operatorname{Hom}_S(X_\lambda,Y)\to\operatorname{Hom}_S(X,Y)$ a aplicação que faz corresponder a cada $S$-morfismo $f_\lambda:X_\lambda\to Y$ o morfismo $f=f_\lambda\circ v_\lambda$, onde $v_\lambda:X\to X_\lambda$ é o morfismo canônico.
A família $(e_\lambda)$ é um sistema indutivo de aplicações, que define, portanto, uma aplicação canônica
\begin{equation}
\label{IV.8.13.1.1}
\varinjlim_\lambda\operatorname{Hom}_S(X_\lambda,Y)\to\operatorname{Hom}_S(X,Y). \tag{8.13.1.1}
\end{equation}

Suponhamos que $X_\alpha$ seja quase-compacto (resp. quase-compacto e quase-separado) e que o morfismo estrutural $Y\to S$ seja localmente de tipo finito (resp. localmente de apresentação finita).
Então a aplicação \sref{IV.8.13.1.1} é injetiva (resp. bijetiva).
\end{proposition}

\begin{proof}
Para $\lambda\geqslant\alpha$, ponhamos $Z_\lambda=Y\times_S X_\lambda$, de modo que $Z_\lambda=Z_\alpha\times_{X_\alpha}X_\lambda$.
Do mesmo modo, ponhamos $Z=Y\times_S X=Z_\alpha\times_{X_\alpha}X$; decorre então de \sref{IV.8.2.5} que, pondo também $w_{\lambda\mu}=1\times v_{\lambda\mu}:Z_\mu\to Z_\lambda$ para $\alpha\leqslant\lambda\leqslant\mu$ e $w_\lambda=1\times v_\lambda:Z\to Z_\lambda$ para $\alpha\leqslant\lambda$, o pré-esquema $Z$ é o limite projetivo do sistema $(Z_\lambda,w_{\lambda\mu})$, e os $w_\lambda$ são os morfismos canônicos correspondentes.
Além disso, o morfismo $Z_\alpha\to X_\alpha$ é localmente de tipo finito (resp. localmente de apresentação finita) por \sref{IV.1.3.4} e \sref{IV.1.4.3}.
Por fim, por \sref[I]{I.3.3.14}, tem-se
\[
\operatorname{Hom}_S(X_\lambda,Y)=\operatorname{Hom}_{X_\lambda}(X_\lambda,Z_\lambda) \quad\text{e}\quad \operatorname{Hom}_S(X,Y)=\operatorname{Hom}_X(X,Z)
\]
Basta agora aplicar \sref{IV.8.8.2},~(i), tomando $X_\lambda=S_\lambda$ e substituindo $Y_\lambda$ por $Z_\lambda$.
\end{proof}

\begin{corollary}[8.13.2]
\label{IV.8.13.2}
Com as notações de \sref{IV.8.13.1}, suponhamos que $X_\alpha$ seja quase-compacto e quase-separado, e que os $v_{\alpha\lambda}$ sejam afins para $\alpha\leqslant\lambda$; suponhamos, além disso, que $Y=\varprojlim_\rho Y_\rho$, onde $(Y_\rho,t_{\rho\sigma})$ é um sistema projetivo filtrante de $S$-pré-esquemas tal que, para cada $\rho$, o morfismo estrutural $Y_\rho\to S$ seja localmente de apresentação finita.
Então há uma bijeção canônica
\begin{equation}
\label{IV.8.13.2.1}
\operatorname{Hom}_S(X,Y)\xrightarrow{\sim}\varprojlim_\rho\left(\varinjlim_\lambda\operatorname{Hom}_S(X_\lambda,Y_\rho)\right). \tag{8.13.2.1}
\end{equation}
\end{corollary}

\begin{proof}
Com efeito, como $Y$ é o limite projetivo dos $Y_\rho$, a aplicação canônica $\operatorname{Hom}_S(X,Y)\to\varprojlim_\rho\operatorname{Hom}_S(X,Y_\rho)$ é bijetiva; por outro lado, as hipóteses fornecem, para cada $\rho$, uma bijeção canônica $\operatorname{Hom}_S(X,Y_\rho)\xrightarrow{\sim}\varinjlim_\lambda\operatorname{Hom}_S(X_\lambda,Y_\rho)$ em virtude de \sref{IV.8.13.1}; daí a conclusão.
\end{proof}

\begin{env}[8.13.3]
\label{IV.8.13.3}
Os resultados precedentes permitem interpretar, na teoria dos pré-esquemas, as noções de ``provariedade'' e ``proesquema'' que surgem em certas aplicações (por exemplo, na teoria local dos corpos de classes, seguindo as ideias de Serre~[39], ou na teoria de N\'eron da redução de
\oldpage[IV-3]{50}
variedades abelianas~[32]).
Recordemos brevemente aqui a noção de \emph{pro-objeto} de uma categoria, remetendo ao cap.~V para desenvolvimentos mais amplos (aliás, não utilizaremos antes do cap.~V a interpretação que segue, e o leitor pode, portanto, omiti-la até o fim deste número).
Dada uma categoria $\mathcal{C}$, a categoria $\mathbf{Pro}(\mathcal{C})$ dos \emph{pro-objetos} de $\mathcal{C}$ tem como objetos os \emph{sistemas projetivos} (no universo considerado) $\mathbf{X}=(X_\mu)_{\mu\in M}$ de objetos de $\mathcal{C}$ cujos conjuntos de índices (que dependem do sistema projetivo) se supõem conjuntos pré-ordenados filtrantes.
Dados dois pro-objetos desse tipo $\mathbf{X}=(X_\mu)_{\mu\in M}$ e $\mathbf{X}'=(X'_{\mu'})_{\mu'\in M'}$, os \emph{morfismos} de $\mathbf{X}$ em $\mathbf{X}'$ são, por definição, os elementos do conjunto $\varprojlim_{\mu'}(\varinjlim_\mu\operatorname{Hom}(X_\mu,X'_{\mu'}))$; é imediato que se podem tomar esses conjuntos como conjuntos de morfismos. A composição de sistemas de morfismos $u_{\mu'}^\mu:X_\mu\to X'_{\mu'}$ e $u_{\mu''}^{\prime\mu'}:X'_{\mu'}\to X''_{\mu''}$, que são indutivos em relação ao índice superior e projetivos em relação ao índice inferior, efetua-se ``argumento por argumento'', isto é, considerando o sistema $u_{\mu''}^{\prime\mu'}\circ u_{\mu'}^\mu$.
\end{env}

\begin{env}[8.13.4]
\label{IV.8.13.4}
Seja $S$ um pré-esquema quase-compacto e quase-separado, e denotemos por $\mathcal{C}$ a subcategoria plena da categoria $(\mathbf{Sch})_{/S}$ dos $S$-pré-esquemas, formada pelos $S$-pré-esquemas $X$ que possuem a seguinte propriedade: o morfismo estrutural $X\to S$ se fatora como $X\xrightarrow{g}X_0\xrightarrow{f}S$, onde $g:X\to X_0$ é afim e $f:X_0\to S$ é \emph{de apresentação finita}; diremos, abreviadamente, que os pré-esquemas de $\mathcal{C}$ são \emph{essencialmente afins sobre} $S$.
Consideremos, por outro lado, a subcategoria plena $\mathcal{C}'_0$ de $(\mathbf{Sch})_{/S}$ formada pelos $S$-pré-esquemas \emph{de apresentação finita}, e a categoria $\mathbf{Pro}(\mathcal{C}'_0)$ de pro-objetos de $\mathcal{C}'_0$.
Diremos que um objeto $\mathbf{X}=(X_\mu)_{\mu\in M}$ de $\mathbf{Pro}(\mathcal{C}'_0)$ é \emph{essencialmente afim} se existe $\gamma\in M$ tal que, para $\mu\geqslant\gamma$, o morfismo de transição $v_{\gamma\mu}:X_\mu\to X_\gamma$ é afim (o que implica que, para $\gamma\leqslant\nu\leqslant\mu$, $v_{\nu\mu}:X_\mu\to X_\nu$ é afim).
Observe-se que um objeto de $\mathbf{Pro}(\mathcal{C}'_0)$ isomorfo a um objeto essencialmente afim \emph{não é} necessariamente ele próprio essencialmente afim.
Denotaremos por $\mathcal{C}'$ a subcategoria plena de $\mathbf{Pro}(\mathcal{C}'_0)$ formada pelos pro-objetos essencialmente afins.

Posto isso, decorre de \sref{IV.8.2.2} e \sref{IV.8.2.3} que, para todo objeto $\mathbf{X}=(X_\mu)_{\mu\in M}$ de $\mathcal{C}'$, existe o $S$-pré-esquema $X=\varprojlim_\mu X_\mu$; além disso, como o morfismo $X\to X_\mu$ é afim para todo $\mu$ suficientemente grande \sref{IV.8.2.2}, $X$ é, por definição, \emph{essencialmente afim sobre} $S$.
Ponhamos $X=L(\mathbf{X})$; mostremos que assim se definiu um \emph{funtor} canônico
\begin{equation}
\label{IV.8.13.4.1}
L:\mathcal{C}'\to\mathcal{C} \tag{8.13.4.1}
\end{equation}

Para dois objetos $\mathbf{X}=(X_\mu)$ e $\mathbf{X}'=(X'_{\mu'})$ de $\mathcal{C}'$, tem-se, para cada $\mu'$, uma aplicação canônica
\[
\varinjlim_\mu\operatorname{Hom}_S(X_\mu,X'_{\mu'})\to\operatorname{Hom}_S(\varprojlim_\mu X_\mu,X'_{\mu'})
\]
definida em \sref{IV.8.13.1.1}; por outro lado, por definição do limite projetivo, existe uma bijeção canônica
\[
\varprojlim_{\mu'}\operatorname{Hom}_S(\varprojlim_\mu X_\mu,X'_{\mu'})\xrightarrow{\sim}\operatorname{Hom}_S(\varprojlim_\mu X_\mu,\varprojlim_{\mu'}X'_{\mu'})
\]
\oldpage[IV-3]{51}

da qual se obtém uma aplicação canônica
\begin{equation}
\label{IV.8.13.4.2}
\varprojlim_{\mu'}\varinjlim_\mu\operatorname{Hom}_S(X_\mu,X'_{\mu'})\to\operatorname{Hom}_S(\varprojlim_\mu X_\mu,\varprojlim_{\mu'}X'_{\mu'}) \tag{8.13.4.2}
\end{equation}
evidentemente funtorial em $\mathbf{X}$ e $\mathbf{X}'$, e que completa a definição do funtor $L$.
\end{env}

\begin{proposition}[8.13.5]
\label{IV.8.13.5}
Sob as hipóteses e com as notações de \sref{IV.8.13.4}, o funtor $L$ é completamente fiel.
Se, além disso, $S$ é um pré-esquema noetheriano (o que já implica que $S$ é quase-compacto e quase-separado \sref{IV.1.2.8}), $L$ é \emph{uma equivalência de categorias}.
\end{proposition}

\begin{proof}
Dizer que $L$ é completamente fiel significa que a aplicação \sref{IV.8.13.4.2} é bijetiva, o que é um caso particular de \sref{IV.8.13.2}.
Com efeito, como os morfismos estruturais $X_\mu\to S$ são de apresentação finita, são em particular quase-compactos e quase-separados; portanto, os $X_\mu$ são quase-compactos e quase-separados.

Para demonstrar que, quando $S$ é noetheriano, $L$ é uma equivalência de categorias, basta, uma vez que já sabemos que $L$ é completamente fiel, provar que todo $S$-pré-esquema essencialmente afim $X$ é $S$-\emph{isomorfo} a um objeto da forma $L(\mathbf{X})$, onde $\mathbf{X}\in\mathcal{C}'$ $(\mathbf{0}_{\text{III}},~8.1.5)$.
Ora, por hipótese existe uma fatoração $X\xrightarrow{g}X_0\xrightarrow{f}S$ do morfismo estrutural, sendo $f$ de apresentação finita e $g$ afim.
Podemos, portanto, escrever $X=\operatorname{Spec}(\sh{A})$, onde $\sh{A}$ é uma $\sh{O}_{X_0}$-Álgebra quase-coerente \sref[II]{II.1.3.1}.
Ora, como $X_0$ é noetheriano (pois é de tipo finito sobre $S$ e $S$ é noetheriano), $\sh{A}$ é o limite indutivo filtrante da família $(\sh{A}_\lambda)$ de suas sub-$\sh{O}_{X_0}$-Álgebras quase-coerentes de tipo finito \sref[I]{I.9.6.6}.
Ponhamos $X_\lambda=\operatorname{Spec}(\sh{A}_\lambda)$; os morfismos $X_\lambda\to X_0$ são de tipo finito e, portanto, de apresentação finita, visto que $X_0$ é noetheriano; por conseguinte, o mesmo vale para os morfismos compostos $X_\lambda\to X_0\to S$ \sref{IV.1.6.2}; em outras palavras, os $X_\lambda$ pertencem a $\mathcal{C}'_0$ e, como os morfismos $X_\lambda\to X_0$ são afins, $\mathbf{X}=(X_\lambda)$ é um objeto de $\mathcal{C}'$ cujo limite projetivo existe e é $S$-isomorfo a $X$ por \sref{IV.8.2.2}.
Isto conclui a demonstração.
\end{proof}

\begin{remark}[8.13.6]
\label{IV.8.13.6}
Decorre de \sref{IV.1.6.2} e \sref[II]{II.1.6.2} que, se $X$ e $Y$ são essencialmente afins sobre $S$, o mesmo vale para $X\times_S Y$.
Conclui-se, por exemplo, de $(\mathbf{0}_{\text{III}},~8.2.5)$ que um $\mathcal{C}$-grupo nada mais é do que um $(\mathbf{Sch})_{/S}$-grupo que é um pré-esquema essencialmente afim sobre $S$.
Por outro lado, os produtos finitos existem na categoria $\mathcal{C}'$: com efeito, se $\mathbf{X}=(X_\mu)_{\mu\in M}$ e $\mathbf{Y}=(Y_\rho)_{\rho\in R}$ são dois objetos de $\mathcal{C}'$, os produtos $X_\mu\times_S Y_\rho$ são $S$-pré-esquemas de apresentação finita e, tomando como morfismos de transição $X_\nu\times_S Y_\sigma\to X_\mu\times_S Y_\rho$ os produtos dos morfismos de transição $X_\nu\to X_\mu$ e $Y_\sigma\to Y_\rho$, vê-se imediatamente que $(X_\mu\times_S Y_\rho)$ é um produto de $\mathbf{X}$ e $\mathbf{Y}$ em $\mathbf{Pro}(\mathcal{C}'_0)$; além disso, por \sref[II]{II.1.6.2}, os morfismos de transição assim definidos são afins para $\mu$ e $\rho$ suficientemente grandes e, portanto, o produto $\mathbf{X}\times\mathbf{Y}$ também pertence a $\mathcal{C}'$.
Conclui-se então, como acima, que um $\mathcal{C}'$-grupo é um $\mathbf{Pro}(\mathcal{C}'_0)$-grupo essencialmente afim.
Deduz-se de \sref{IV.8.13.5} que as categorias dos $\mathcal{C}'$-grupos e dos $\mathcal{C}$-grupos são \emph{equivalentes} quando $S$ é \emph{noetheriano}.
Parece plausível que, quando $S$ é o espectro de um corpo $k$, a categoria dos $\mathcal{C}$-grupos seja \emph{equivalente} à dos $\boldsymbol{K}$-\emph{grupos}, onde $\boldsymbol{K}$ é a categoria dos $S$-pré-esquemas \emph{quase-compactos}; em outras palavras, todo pré-esquema em grupos quase-compacto sobre $k$ seria essencialmente afim.
Por outro lado, se denotarmos
\oldpage[IV-3]{52}
por $\mathcal{C}'_0$-$\boldsymbol{gr}$ a categoria dos $\mathcal{C}'_0$-grupos, é muito provável que a categoria dos $\mathcal{C}'$-grupos seja equivalente à subcategoria plena de $\mathbf{Pro}(\mathcal{C}'_0\text{-}\boldsymbol{gr})$ formada pelos ``grupos proalgébricos essencialmente afins'', isto é, pelos sistemas projetivos $\mathbf{G}=(G_\mu)_{\mu\in M}$, onde os $G_\mu$ são grupos algébricos sobre $k$ e os morfismos de transição $G_\nu\to G_\mu$ são afins para $\mu$ suficientemente grande (o que também se pode exprimir dizendo que $\mathbf{G}$ é uma extensão de um grupo algébrico por um grupo proalgébrico afim).
A conjunção dessas duas conjecturas é, ademais, equivalente à afirmação seguinte: todo pré-esquema em grupos quase-compacto sobre $k$ é uma extensão de um ``grupo algébrico'' (isto é, um pré-esquema em grupos de tipo finito sobre $k$) por um pré-esquema em grupos \emph{afim} sobre $k$.
Como os únicos grupos proalgébricos encontrados até agora na prática são, de fato, essencialmente afins, seria, portanto, vantajoso substituir o estudo dos grupos proalgébricos gerais (introduzidos e estudados por Serre~[40]) pelo dos esquemas em grupos quase-compactos sobre $k$, cuja definição é conceitualmente muito mais simples.
\end{remark}

\subsection{Caracterização de um pré-esquema localmente de apresentação finita sobre outro em termos do funtor que ele representa}
\label{subsection:IV.8.14}

\begin{env}[8.14.1]
\label{IV.8.14.1}
Dado um pré-esquema $S$, diremos, como em \sref{IV.8.13.4}, que um sistema projetivo filtrante $(X_\lambda,v_{\lambda\mu})$ de $S$-pré-esquemas é \emph{essencialmente afim} se existe $\alpha$ tal que $v_{\alpha\lambda}$ é um morfismo afim para $\lambda\geqslant\alpha$.

O enunciado seguinte, que será especialmente útil no cap.~V, precisa \sref{IV.8.8.2},~(i) fornecendo uma recíproca:
\end{env}

\begin{proposition}[8.14.2]
\label{IV.8.14.2}
Seja $S$ um pré-esquema e seja $f:X\to S$ um morfismo.
Para todo $S$-pré-esquema $T$, ponhamos
\[
h_X(T)=\operatorname{Hom}_S(T,X)
\]
de modo que $h_X$ é um funtor contravariante da categoria $(\mathbf{Sch})_{/S}$ dos $S$-pré-esquemas para a categoria $\mathbf{Ens}$ dos conjuntos $(\mathbf{0}_{\text{III}},~8.1.1)$, e $X$ é um objeto que representa esse funtor $(\mathbf{0}_{\text{III}},~8.1.8)$.
As condições seguintes são equivalentes:
\begin{enumerate}[label=\emph{\alph*)}]
  \item $f$ é localmente de apresentação finita.
  \item Para todo sistema projetivo filtrante $(Z_\lambda)$ de $S$-pré-esquemas, essencialmente afim \sref{IV.8.13.4} e formado por pré-esquemas quase-compactos e quase-separados, a aplicação canônica \sref{IV.8.13.1.1}
\begin{equation}
\label{IV.8.14.2.1}
\varinjlim_\lambda h_X(Z_\lambda)\to h_X(\varprojlim_\lambda Z_\lambda) \tag{8.14.2.1}
\end{equation}
é bijetiva.
  \item Para todo sistema projetivo filtrante $(Z_\lambda)$ de $S$-pré-esquemas tal que os $Z_\lambda$ sejam esquemas afins, a aplicação \sref{IV.8.14.2.1} é bijetiva.
  \item[\rm c$'$)] Para todo aberto afim $U$ de $S$ e todo sistema projetivo filtrante $(Z_\lambda)$ de $U$-pré-esquemas tal que os $Z_\lambda$ sejam esquemas afins, a aplicação \sref{IV.8.14.2.1} é bijetiva.
\end{enumerate}
\end{proposition}

\begin{proof}
O fato de que \emph{a)} implica \emph{b)} é precisamente \sref{IV.8.13.1}; é trivial que \emph{b)} implica \emph{c)} e que \emph{c)} implica \emph{c$'$)}.
Resta ver que \emph{c$'$)} implica \emph{a)} e, como a propriedade \emph{a)} é local sobre $S$, podemos nos restringir ao caso em que $S$ é \emph{afim}.
\oldpage[IV-3]{53}

Suponhamos primeiro que $X$ também seja afim; a afirmação a demonstrar é então equivalente à seguinte:

\begin{corollary}[8.14.2.2]
\label{IV.8.14.2.2}
Sejam $A$ um anel e $B$ uma $A$-álgebra.
Para que, para todo sistema indutivo filtrante $(C_\lambda)$ de $A$-álgebras, a aplicação canônica
\begin{equation}
\label{IV.8.14.2.3}
\varinjlim_\lambda\operatorname{Hom}_{A\text{-}\mathrm{alg.}}(B,C_\lambda)\to\operatorname{Hom}_{A\text{-}\mathrm{alg.}}(B,\varinjlim_\lambda C_\lambda) \tag{8.14.2.3}
\end{equation}
seja bijetiva, é necessário e suficiente que $B$ seja uma $A$-álgebra de apresentação finita.
\end{corollary}

Resta apenas demonstrar que a condição é necessária.
Tomemos primeiro, para $(C_\lambda)$, o sistema indutivo filtrante das sub-$A$-álgebras de tipo finito de $B$, de modo que $\varinjlim_\lambda C_\lambda=B$.
O fato de que \sref{IV.8.14.2.3} seja bijetiva implica, em particular, que a aplicação identidade $1_B:B\to B$ se fatora como $B\to C_\lambda\to B$ para um $\lambda$ conveniente, o que implica que $C_\lambda=B$; portanto, $B$ é uma $A$-álgebra \emph{de tipo finito}.
Escrevamos $B=C/\mathfrak{J}$, onde $C=A[T_1,\ldots,T_n]$ e $\mathfrak{J}$ é um ideal de $C$. O ideal $\mathfrak{J}$ é o limite indutivo filtrante de seus ideais de tipo finito $\mathfrak{J}_\lambda\subset\mathfrak{J}$; pondo $C_\lambda=C/\mathfrak{J}_\lambda$ e utilizando a exatidão dos limites indutivos filtrantes, vê-se que $B$ é novamente isomorfo ao limite indutivo do sistema indutivo filtrante $(C_\lambda)$.
Existem, portanto, um $\lambda$ e um $A$-homomorfismo $u:B\to C_\lambda$ tais que o composto $B\xrightarrow{u}C_\lambda\xrightarrow{p_\lambda}B$ (onde $p_\lambda$ é o homomorfismo canônico) é a identidade.
Seja $q_\lambda:C\to C_\lambda$ o homomorfismo canônico, e ponhamos
$t_i=p_\lambda(q_\lambda(T_i))$.
Então $p_\lambda(u(t_i))=p_\lambda(q_\lambda(T_i))$, ou seja,
$u(t_i)-q_\lambda(T_i)\in\mathfrak{J}/\mathfrak{J}_\lambda$.
Existe, portanto, um $\mu\geqslant\lambda$ tal que esses $n$ elementos pertencem a
$\mathfrak{J}_\mu/\mathfrak{J}_\lambda$ ($1\leqslant i\leqslant n$).
Se $p_{\mu\lambda}:C_\lambda\to C_\mu$ é o homomorfismo canônico, então
$p_{\mu\lambda}(u(t_i))=p_{\mu\lambda}(q_\lambda(T_i))=q_\mu(T_i)$.
Substituindo $\lambda$ por $\mu$ e $u$ por $p_{\mu\lambda}\circ u$, podemos, portanto,
supor que $u(t_i)=q_\lambda(T_i)$ para todo $i$; e, se
$r=p_\lambda\circ q_\lambda$ é o homomorfismo canônico
$C\to C/\mathfrak{J}=B$, podemos escrever $u(r(T_i))=q_\lambda(T_i)$ para todo $i$,
de onde $q_\lambda=u\circ r$.
Isso implica necessariamente que $\mathfrak{J}=\mathfrak{J}_\lambda$: se $z\in\mathfrak{J}$, então $r(z)=0$ e, portanto, $q_\lambda(z)=u(r(z))=0$ e $z\in\mathfrak{J}_\lambda$. Assim, $B=C_\lambda$.

Passemos agora ao caso em que $S$ é afim e $X$ arbitrário; trata-se de demonstrar que todo aberto afim $V$ de $X$ é de apresentação finita sobre $S$ e, em virtude do que acaba de ser demonstrado, basta provar que, para todo sistema projetivo filtrante $(Z_\lambda)$ de $S$-pré-esquemas afins, a aplicação
\begin{equation}
\label{IV.8.14.2.4}
\varinjlim_\lambda\operatorname{Hom}_S(Z_\lambda,V)\to\operatorname{Hom}_S(\varprojlim_\lambda Z_\lambda,V) \tag{8.14.2.4}
\end{equation}
é bijetiva.
É imediato que essa aplicação é injetiva, pois, se $(v_\lambda)$, $(v'_\lambda)$ são dois sistemas indutivos de $S$-homomorfismos $v_\lambda:Z_\lambda\to V$, $v'_\lambda:Z_\lambda\to V$ tais que os morfismos correspondentes
\[
Z\xrightarrow{u_\lambda}Z_\lambda\xrightarrow{v_\lambda}V,\qquad
Z\xrightarrow{u_\lambda}Z_\lambda\xrightarrow{v'_\lambda}V
\]
são iguais (sendo $u_\lambda$ o morfismo canônico), então os morfismos
\[
Z\xrightarrow{u_\lambda}Z_\lambda\xrightarrow{v_\lambda}V\xrightarrow{j}X,\qquad Z\xrightarrow{u_\lambda}Z_\lambda\xrightarrow{v'_\lambda}V\xrightarrow{j}X
\]
(onde $j$ é a injeção canônica) são iguais, o que implica, por hipótese, $j\circ v_\lambda=j\circ v'_\lambda$ para um $\lambda$ conveniente e, portanto, $v_\lambda=v'_\lambda$.

\oldpage[IV-3]{54}

Resta demonstrar que \sref{IV.8.14.2.4} é sobrejetiva.
Seja $v:Z\to V$ um $S$-morfismo.
Por hipótese, existem um índice $\lambda$ e um $S$-morfismo
$w_\lambda:Z_\lambda\to X$ tais que $j\circ v$ se fatora como
\[
 Z\xrightarrow{u_\lambda}Z_\lambda\xrightarrow{w_\lambda}X.
\]
Basta demonstrar que existe um $\mu\geqslant\lambda$ tal que
\[
 Z_\mu\xrightarrow{u_{\lambda\mu}}Z_\lambda\xrightarrow{w_\lambda}X
\]
(onde $u_{\lambda\mu}$ é o morfismo de transição) se fatora como
\[
 Z_\mu\xrightarrow{v_\mu}V\xrightarrow{j}X.
\]
Para todo $\mu\geqslant\lambda$, ponhamos
$U_\mu=u_{\lambda\mu}^{-1}(w_\lambda^{-1}(V))$.
Então
\[
 u_\mu^{-1}(U_\mu)=u_\lambda^{-1}(U_\lambda)
 =u_\lambda^{-1}(w_\lambda^{-1}(V))
 =v^{-1}(V)=Z.
\]
Os $Z_\lambda$ são quase-compactos e os $U_\mu$, por serem abertos, são
ind-construtíveis \sref{IV.1.9.6}; portanto, \sref{IV.8.3.4} fornece um
$\mu\geqslant\lambda$ para o qual $U_\mu=Z_\mu$.
\end{proof}

\begin{remark}[8.14.3]
\label{IV.8.14.3}
O fato de que a aplicação \sref{IV.8.14.2.1} seja injetiva quando $f$ é localmente de tipo
finito \sref{IV.8.8.2},~(i) leva naturalmente a perguntar se esse resultado admite também uma recíproca.
Não é assim, nem mesmo quando $S$ e $X$ são afins, pois existem monomorfismos
$X\to S$ que não são de tipo finito \sref[I]{I.2.4.2}, e estes refutam tal
conjectura.
\end{remark}
